% Sensibilisation sur la nécessité des deux phénomènes biologiques complémentaires (méiose et fécondation) pour la pérennité de l'espèce — SVT Tle D — Cours signé OnBuch+
% Programme officiel MINESEC. Compiler avec : tectonic N03-meiose-fecondation-perennite-espece.tex
\newcommand{\DOCMATIERE}{SVT}
\newcommand{\DOCNIVEAU}{Tle D}
\newcommand{\DOCMODULE}{Module I : Le Monde Vivant — Reproduction sexuée chez les Mammifères et les Spermaphytes}
\newcommand{\DOCLECON}{N03}
\newcommand{\DOCTITRE}{Sensibilisation sur la nécessité des deux phénomènes biologiques complémentaires (méiose et fécondation) pour la pérennité de l'espèce}
% OnBuch+ — préambule « cours signés » ENRICHI (compilé avec tectonic / XeLaTeX)
% Système visuel « Pop » d'OnBuch+ : fond crème (jamais blanc), bordures encre
% épaisses, ombres « dures » décalées, titres Archivo Black en capitales.
% Version MATHÉMATIQUES : graphes (pgfplots/tikz), boîtes pédagogiques
% (définition, propriété, méthode, attention, le savais-tu, exercice résolu,
% exercice, TP, bilan), page de garde et signature OnBuch+ ancrée en bas.
%
% Macros attendues dans main.tex AVANT \input{preamble} :
%   \DOCMATIERE  \DOCNIVEAU  \DOCMODULE  \DOCLECON (numéro)  \DOCTITRE
\documentclass[11pt]{article}

\usepackage{fontspec}
\usepackage[french]{babel}
\usepackage[a4paper,margin=2cm,top=3.4cm,bottom=2.8cm,headheight=1.4cm,footskip=1.3cm]{geometry}
\usepackage{amsmath,amssymb,amsfonts}
\usepackage{mathtools} % \xmapsto, \coloneqq... souvent utilisés par les modèles
\usepackage{cancel} % simplifications barrées (bilans de forces)
% \abs{z} (module, valeur absolue) et \norm{u} (norme), utilisés par les modèles
\providecommand{\abs}{}\let\abs\relax\DeclarePairedDelimiter{\abs}{\lvert}{\rvert}
\providecommand{\norm}{}\let\norm\relax\DeclarePairedDelimiter{\norm}{\lVert}{\rVert}
\usepackage[table]{xcolor} % [table] : fournit aussi \rowcolors (alternance de lignes)
\usepackage{pagecolor}
\usepackage{tikz}
\usetikzlibrary{arrows.meta,positioning,calc,shapes.geometric,shapes.misc,shapes.arrows,decorations.pathmorphing,decorations.pathreplacing,decorations.markings,patterns,shadows,mindmap,trees,fit,backgrounds,matrix,intersections,angles,quotes}
\usepackage{pgfplots}
\usepackage[version=4]{mhchem} % équations nucléaires \ce{^{235}_{92}U}
\usepackage[european,siunitx]{circuitikz} % schémas électriques
\pgfplotsset{compat=1.18}
\usepgfplotslibrary{fillbetween,groupplots} % aires entre courbes (calcul intégral)
\usepackage{enumitem}
\usepackage{fancyhdr}
% [most] charge toutes les bibliothèques tcolorbox (listings, minted, raster,
% documentation...) et gonfle inutilement la mémoire TeX sur un document long
% (~300 boîtes) : on ne charge que ce qui est réellement utilisé.
\usepackage[skins,breakable]{tcolorbox}
\usepackage{graphicx}
\usepackage{colortbl}
\usepackage{booktabs}
\usepackage{tabularx}
\usepackage{diagbox} % case d'en-tête diagonale des tableaux à double entrée
% Colonnes extensibles centrée / gauche / droite (C, L, R), souvent utilisées par les modèles
\newcolumntype{C}{>{\centering\arraybackslash}X}
\newcolumntype{L}{>{\raggedright\arraybackslash}X}
\newcolumntype{R}{>{\raggedleft\arraybackslash}X}
\usepackage{array}
\usepackage{multicol}
\usepackage{siunitx}
% parse-numbers=false : accepte aussi \SI{2,5 \times 10^{-3}}{...}
\sisetup{locale=FR,output-decimal-marker={,},group-minimum-digits=4,parse-numbers=false}
\DeclareSIUnit\ohm{\ensuremath{\symup{\Omega}}} % U+2126 absent de la police
\DeclareSIUnit\division{div} % réglages d'oscilloscope (ms/div, V/div)
\DeclareSIUnit\tr{tr} % tours (vitesse de rotation, ex. \SI{1500}{\tr\per\minute})
\DeclareSIUnit\molar{M} % concentration molaire (ex. \SI{3}{\molar}, \SI{1}{\micro\molar})
\DeclareSIUnit\an{an} % année (le modèle confond parfois avec \year, un compteur TeX) — ex. \SI{5}{\kilo\meter\per\an}
\DeclareSIUnit\FCFA{FCFA} % franc CFA (ex. \SI{500}{\FCFA\per\kilo\gram})
\DeclareSIUnit\degreeBrix{\textdegree Brix} % teneur en sucre d'un sirop/jus (réfractométrie)
\DeclareSIUnit\month{mois} % le modèle utilise parfois \month, un compteur TeX, comme unité de durée
\usepackage{qrcode}
% \degree : commande fréquemment utilisée par les modèles pour un angle ou une
% température en dehors de \SI{}{} (non fournie par les paquets chargés).
\providecommand{\degree}{^{\circ}}
% \qed (fin de démonstration), utilisé par les modèles sans amsthm
\providecommand{\qed}{\hfill$\square$}
% \barre : double barre des valeurs interdites dans un tableau de variations (tkz-tab)
\providecommand{\barre}{\ensuremath{\|}}
% environnement proof (amsthm non chargé) utilisé par les modèles
\ifdefined\proof\else\newenvironment{proof}[1][Démonstration]{\par\noindent\textit{#1.}\ }{\qed\par}\fi
\usepackage{lastpage}
\usepackage{hyperref}
\hypersetup{hidelinks}

% ============================================================
% Palette « Pop » OnBuch+ — mêmes hex que la classe P (plus_ui.dart)
% ============================================================
\definecolor{poporange}{HTML}{F59321}
\definecolor{popdark}{HTML}{E07A0C}
\definecolor{popcream}{HTML}{FFF6E8}
\definecolor{popink}{HTML}{1C1714}
\definecolor{poppurple}{HTML}{7A5AE0}
\definecolor{popgreen}{HTML}{2E9E6A}
\definecolor{poppink}{HTML}{E0457B}
\definecolor{popblue}{HTML}{2F7DE1}
\definecolor{popgold}{HTML}{C98A12}
\definecolor{popmuted}{HTML}{8A7F76}
\definecolor{popline}{HTML}{EADFD2}
\definecolor{popgray}{HTML}{8A7F76}
\colorlet{popgrey}{popgray}
% Alias tolérants (le modèle nomme parfois les couleurs différemment)
\colorlet{popwhite}{white}
\colorlet{popblack}{popink}
\colorlet{popred}{poppink}
\colorlet{popyellow}{popgold}
% Teintes claires (fonds de tableaux/encadrés) — le modèle nomme parfois
% les couleurs avec un suffixe L (ex. popblueL) sans les avoir définies.
\colorlet{poporangeL}{poporange!20}
\colorlet{popdarkL}{popdark!20}
\colorlet{poppurpleL}{poppurple!20}
\colorlet{popgreenL}{popgreen!20}
\colorlet{poppinkL}{poppink!20}
\colorlet{popblueL}{popblue!20}
\colorlet{popgoldL}{popgold!20}
\colorlet{popredL}{popred!20}
\colorlet{popgrayL}{popgray!20}
% Teintes claires pour fonds de tableaux / zones
\colorlet{popblueL}{popblue!10!white}
\colorlet{poporangeL}{poporange!14!white}
\colorlet{popgreenL}{popgreen!12!white}
\colorlet{poppurpleL}{poppurple!12!white}
\colorlet{poppinkL}{poppink!12!white}
\colorlet{popgoldL}{popgold!14!white}

\pagecolor{popcream}

% ============================================================
% Typographie
% ============================================================
\defaultfontfeatures{Ligatures=TeX}
\setmainfont{PlusJakartaSans-Medium.ttf}[Path=../../../fonts/,BoldFont=PlusJakartaSans-Bold.ttf]
% Maths en sans-serif (Fira Math) avec chiffres et lettres droites de Plus
% Jakarta, pour que les formules se fondent dans le texte.
\usepackage[math-style=ISO,bold-style=ISO]{unicode-math}
\setmathfont{FiraMath-Regular.otf}[Path=../../../fonts/]
\setmathfont{PlusJakartaSans-Medium.ttf}[Path=../../../fonts/,range={up/num,up/{Latin,latin},\mathup}]
\usepackage{icomma} % virgule décimale sans espace en mode math
% Caractères mathématiques Unicode absents des polices de texte (Plus Jakarta,
% Archivo Black) : rendus par leur équivalent mathématique, en texte comme en
% formule — sinon ils s'affichent en carré vide (ex. « Arithmétique dans ℤ »).
\usepackage{newunicodechar}
\newunicodechar{ℝ}{\ensuremath{\mathbb{R}}}
\newunicodechar{ℤ}{\ensuremath{\mathbb{Z}}}
\newunicodechar{ℂ}{\ensuremath{\mathbb{C}}}
\newunicodechar{ℕ}{\ensuremath{\mathbb{N}}}
\newunicodechar{ℚ}{\ensuremath{\mathbb{Q}}}
\newunicodechar{𝔻}{\ensuremath{\mathbb{D}}}
\newunicodechar{α}{\ensuremath{\alpha}}
\newunicodechar{β}{\ensuremath{\beta}}
\newunicodechar{γ}{\ensuremath{\gamma}}
\newunicodechar{δ}{\ensuremath{\delta}}
\newunicodechar{ε}{\ensuremath{\varepsilon}}
\newunicodechar{θ}{\ensuremath{\theta}}
\newunicodechar{λ}{\ensuremath{\lambda}}
\newunicodechar{μ}{\ensuremath{\mu}}
\newunicodechar{σ}{\ensuremath{\sigma}}
\newunicodechar{τ}{\ensuremath{\tau}}
\newunicodechar{φ}{\ensuremath{\varphi}}
\newunicodechar{ψ}{\ensuremath{\psi}}
\newunicodechar{ω}{\ensuremath{\omega}}
\newunicodechar{Σ}{\ensuremath{\Sigma}}
% majuscules produites par \MakeUppercase dans les titres (α-aminés -> Α)
\newunicodechar{Α}{\ensuremath{\alpha}}
\newunicodechar{Β}{\ensuremath{\beta}}
\newunicodechar{Γ}{\ensuremath{\Gamma}}
\newunicodechar{Δ}{\ensuremath{\Delta}}
\newunicodechar{Ε}{\ensuremath{\varepsilon}}
\newunicodechar{Θ}{\ensuremath{\Theta}}
\newunicodechar{Λ}{\ensuremath{\Lambda}}
\newunicodechar{Μ}{\ensuremath{\mu}}
\newunicodechar{Τ}{\ensuremath{\tau}}
\newunicodechar{Φ}{\ensuremath{\Phi}}
\newunicodechar{Ψ}{\ensuremath{\Psi}}
\newunicodechar{Ω}{\ensuremath{\Omega}}
\newunicodechar{↦}{\ensuremath{\mapsto}}
\newunicodechar{∈}{\ensuremath{\in}}
\newunicodechar{∉}{\ensuremath{\notin}}
\newunicodechar{∧}{\ensuremath{\wedge}}
\newunicodechar{⇔}{\ensuremath{\Leftrightarrow}}
\newunicodechar{⟺}{\ensuremath{\Longleftrightarrow}}
\newunicodechar{⇒}{\ensuremath{\Rightarrow}}
\newunicodechar{⟹}{\ensuremath{\Longrightarrow}}
\newunicodechar{∘}{\ensuremath{\circ}}
\newunicodechar{∪}{\ensuremath{\cup}}
\newunicodechar{∩}{\ensuremath{\cap}}
\newunicodechar{≡}{\ensuremath{\equiv}}
\newunicodechar{⊂}{\ensuremath{\subset}}
\newunicodechar{⊥}{\ensuremath{\perp}}
\newunicodechar{∃}{\ensuremath{\exists}}
\newunicodechar{∀}{\ensuremath{\forall}}
\newunicodechar{∛}{\ensuremath{\sqrt[3]{\,}}}
\newunicodechar{∜}{\ensuremath{\sqrt[4]{\,}}}
\newunicodechar{⃗}{\ensuremath{{}^{\to}}}
\newunicodechar{ⁿ}{\ifmmode^{n}\else\textsuperscript{\ensuremath{n}}\fi}
\newunicodechar{ˣ}{\ifmmode^{x}\else\textsuperscript{\ensuremath{x}}\fi}
\newunicodechar{ᵇ}{\ifmmode^{b}\else\textsuperscript{\ensuremath{b}}\fi}
\newunicodechar{ʳ}{\ifmmode^{r}\else\textsuperscript{\ensuremath{r}}\fi}
\newunicodechar{ᵘ}{\ifmmode^{u}\else\textsuperscript{\ensuremath{u}}\fi}
\newunicodechar{ʸ}{\ifmmode^{y}\else\textsuperscript{\ensuremath{y}}\fi}
\newunicodechar{ᵃ}{\ifmmode^{a}\else\textsuperscript{\ensuremath{a}}\fi}
\newunicodechar{ᵉ}{\ifmmode^{e}\else\textsuperscript{\ensuremath{e}}\fi}
\newunicodechar{ᵏ}{\ifmmode^{k}\else\textsuperscript{\ensuremath{k}}\fi}
\newunicodechar{ᵖ}{\ifmmode^{p}\else\textsuperscript{\ensuremath{p}}\fi}
\newunicodechar{ᴺ}{\ifmmode^{N}\else\textsuperscript{\ensuremath{N}}\fi}
\newunicodechar{ᶜ}{\ifmmode^{c}\else\textsuperscript{\ensuremath{c}}\fi}
\newunicodechar{ʰ}{\ifmmode^{h}\else\textsuperscript{\ensuremath{h}}\fi}
\newunicodechar{ᵠ}{\ifmmode^{q}\else\textsuperscript{\ensuremath{q}}\fi}
\newunicodechar{ₙ}{\ifmmode_{n}\else\textsubscript{\ensuremath{n}}\fi}
\newunicodechar{ₐ}{\ifmmode_{a}\else\textsubscript{\ensuremath{a}}\fi}
\newunicodechar{ᵢ}{\ifmmode_{i}\else\textsubscript{\ensuremath{i}}\fi}
\newunicodechar{ᵣ}{\ifmmode_{r}\else\textsubscript{\ensuremath{r}}\fi}
\newunicodechar{ₖ}{\ifmmode_{k}\else\textsubscript{\ensuremath{k}}\fi}
\newunicodechar{ᵦ}{\ifmmode_{\beta}\else\textsubscript{\ensuremath{\beta}}\fi}
\newunicodechar{ₓ}{\ifmmode_{x}\else\textsubscript{\ensuremath{x}}\fi}
\newfontfamily\popdisplay{ArchivoBlack-Regular.ttf}[Path=../../../fonts/]
\newfontfamily\poplogo{SpaceGrotesk-Bold.ttf}[Path=../../../fonts/]
\newfontfamily\popsemi{PlusJakartaSans-SemiBold.ttf}[Path=../../../fonts/]
\newfontfamily\popxbold{PlusJakartaSans-ExtraBold.ttf}[Path=../../../fonts/]
\newfontfamily\popxboldls{PlusJakartaSans-ExtraBold.ttf}[Path=../../../fonts/,LetterSpace=3.0]

\setlist[enumerate]{leftmargin=1.7em,itemsep=5pt,topsep=4pt}
\setlist[itemize]{leftmargin=1.7em,itemsep=4pt,topsep=4pt,label={\color{poporange}\textbullet}}
\setlength{\parindent}{0pt}
\setlength{\parskip}{5pt}
\renewcommand{\arraystretch}{1.25}

% pgfplots : style Pop par défaut
\pgfplotsset{
  every axis/.append style={axis line style={popink,line width=1pt},tick style={popink},
    grid style={popline,line width=0.5pt},label style={font=\small},tick label style={font=\footnotesize}},
  popcurve/.style={poporange,line width=1.8pt,no markers,smooth},
}
% Même style disponible dans un \draw TikZ simple (hors pgfplots)
\tikzset{popcurve/.style={poporange,line width=1.8pt}}
\tikzset{popgrid/.style={popline,very thin}}
\colorlet{popcurve}{poporange} % le nom du style est parfois utilisé comme couleur

% ============================================================
% Pill / squiggle
% ============================================================
\newcommand{\poppill}[3]{%
  \tikz[baseline=-0.55ex]\node[draw=popink,line width=1.2pt,rounded corners=2.6pt,%
    fill=#2,inner xsep=6.5pt,inner ysep=3pt]{%
    \popxboldls\fontsize{7}{7}\selectfont\color{#3}\MakeUppercase{#1}};%
}
\newcommand{\popsquiggle}[1]{%
  \tikz[baseline=-0.4ex]{
    \draw[#1,line width=2.6pt,line cap=round]
      plot [smooth] coordinates {(0,0.09) (0.34,0.01) (0.68,0.21) (1.02,0.01) (1.36,0.10)};
  }%
}
% Mot-clé surligné (marqueur orange sous le texte)
\newcommand{\cle}[1]{{\popxbold\color{popdark}#1}}

% ============================================================
% En-tête / pied
% ============================================================
\pagestyle{fancy}
\fancyhf{}
\renewcommand{\headrulewidth}{0pt}
\renewcommand{\footrulewidth}{0pt}
\lhead{\raisebox{-0.15ex}{\poplogo\fontsize{13}{13}\selectfont\color{popink}OnBuch\color{poporange}+}}
\rhead{\poppill{\DOCMATIERE\ \textbullet\ \DOCNIVEAU\ \textbullet\ Leçon \DOCLECON}{white}{popink}}
\lfoot{\popxbold\fontsize{7.6}{7.6}\selectfont\color{popmuted}\MakeUppercase{Cours signé OnBuch+}\ \textbullet\ {\popsemi\DOCTITRE}}
\rfoot{\tikz[baseline=-0.6ex]\node[rounded corners=8.5pt,fill=popink,minimum height=17pt,minimum width=17pt,inner xsep=5pt,inner sep=0]{\popxbold\fontsize{8}{8}\selectfont\color{popcream}\thepage\,/\,\pageref*{LastPage}};}

% ============================================================
% Titres
% ============================================================
\newcommand{\chaptitre}[2]{%
  \begin{tcolorbox}[enhanced,colback=poporange,colframe=popink,boxrule=2.6pt,arc=11pt,
      drop shadow=popink,left=15pt,right=15pt,top=12pt,bottom=11pt,before skip=3pt,after skip=18pt]
    \poppill{#2}{popink}{popcream}\par\vspace{9pt}
    {\raggedright\hyphenpenalty=10000\exhyphenpenalty=10000\popdisplay\fontsize{22}{24}\selectfont\color{popcream}\MakeUppercase{#1}\par}
  \end{tcolorbox}
}
% Section numérotée : pastille orange avec le numéro + titre Archivo
\newcounter{popsec}
\newcommand{\coursec}[1]{%
  \stepcounter{popsec}\setcounter{popsub}{0}%
  \par\vspace{16pt}\noindent
  \begin{minipage}{\linewidth}
  \tikz[baseline=-0.7ex]\node[draw=popink,line width=1.6pt,fill=poporange,rounded corners=4pt,minimum size=22pt,inner sep=0]{\popdisplay\fontsize{12}{12}\selectfont\color{popcream}\thepopsec};%
  \hspace{8pt}{\popdisplay\fontsize{15}{17}\selectfont\color{popink}\MakeUppercase{#1}}\par
  \vspace{4pt}\noindent{\color{poporange}\rule{36pt}{3.6pt}}
  \end{minipage}\par\nopagebreak\vspace{8pt}
}
\newcounter{popsub}
\newcommand{\courssub}[1]{%
  \stepcounter{popsub}%
  \par\vspace{9pt}\noindent{\popxbold\fontsize{11.5}{13}\selectfont\color{popink}{\color{poporange}\thepopsec.\thepopsub}\hspace{6pt}#1}\par\nopagebreak\vspace{4pt}
}

% ============================================================
% Boîtes pédagogiques — même recette : fond blanc, bordure épaisse colorée,
% ombre dure encre, badge plein. Titre optionnel [#1].
% ============================================================
\tcbset{popbase/.style={breakable,enhanced,colback=white,boxrule=2.3pt,arc=9pt,
  shadow={3pt}{-3pt}{0pt}{popink},left=14pt,right=14pt,top=11pt,bottom=11pt,before skip=11pt,after skip=15pt,
  fontupper=\popsemi\fontsize{10.6}{14.5}\selectfont\color{popink},
  fontlower=\popsemi\fontsize{10.6}{14.5}\selectfont\color{popink}}}
% Coche dessinée (le glyphe ✓ n'existe pas dans les polices du document)
\newcommand{\popcheck}[1]{\tikz[baseline=-0.5ex]\draw[#1,line width=1.6pt,line cap=round,line join=round](-0.13,0.0)--(-0.04,-0.09)--(0.13,0.1);}
\providecommand{\coche}{\popcheck{popgreen}}
\providecommand{\resultat}[1]{\par\smallskip\noindent\fcolorbox{popgreen}{popgreenL}{\parbox{\dimexpr\linewidth-2\fboxsep-2\fboxrule\relax}{\textbf{Résultat.} #1}}\par\smallskip}
\newcommand{\popboxhead}[3]{\poppill{#1}{#2}{white}\ifx\relax#3\relax\else\hspace{6pt}{\popxbold\fontsize{10.6}{12}\selectfont\color{popink}#3}\fi\par\vspace{7pt}}

\newtcolorbox{aretenir}[1][]{popbase,colframe=popgreen,before upper={\popboxhead{À retenir}{popgreen}{#1}}}
\newtcolorbox{exemplebox}[1][]{popbase,colframe=poporange,before upper={\popboxhead{Exemple}{poporange}{#1}}}
\newtcolorbox{definition}[1][]{popbase,colframe=popblue,before upper={\popboxhead{Définition}{popblue}{#1}}}
\newtcolorbox{propriete}[1][]{popbase,colframe=poppurple,before upper={\popboxhead{Propriété}{poppurple}{#1}}}
\newtcolorbox{methode}[1][]{popbase,colframe=popgold,before upper={\popboxhead{Méthode}{popgold}{#1}}}
\newtcolorbox{attention}[1][]{popbase,colframe=poppink,before upper={\popboxhead{Attention}{poppink}{#1}}}
\newtcolorbox{experience}[1][]{popbase,colframe=popblue,colback=popblueL,before upper={\popboxhead{Expérience}{popblue}{#1}}}
% « Le savais-tu ? » : carte violette pleine (contexte, culture, Cameroun)
\newtcolorbox{savaistu}[1][]{popbase,colback=poppurple,colframe=popink,
  fontupper=\popsemi\fontsize{10.4}{14.2}\selectfont\color{white},
  before upper={\poppill{Le savais-tu\,?}{white}{poppurple}\ifx\relax#1\relax\else\hspace{6pt}{\popxbold\color{white}#1}\fi\par\vspace{7pt}}}
% Exercice résolu : énoncé en haut, \tcblower puis la solution sur fond crème
\newcounter{popexo}\newcounter{popexr}
\newtcolorbox{exoresolu}[1][]{popbase,colframe=poporange,colbacklower=popcream,
  segmentation style={popink,line width=1.2pt,dashed},
  before upper={\stepcounter{popexr}\popboxhead{Exercice résolu \thepopexr}{poporange}{#1}},
  before lower={\poppill{Solution}{popink}{popcream}\par\vspace{6pt}}}
% Exercice (non résolu en place) — la correction est dans la partie Corrigés
\newtcolorbox{exercice}[1][]{popbase,colframe=popink,
  before upper={\stepcounter{popexo}\popboxhead{Exercice \thepopexo}{popink}{#1}}}
% Corrigé
\newtcolorbox{corrige}[1][]{popbase,colframe=popgreen,colback=popgreenL,
  before upper={\popboxhead{Corrigé}{popgreen}{#1}}}
% Objectifs / prérequis / bilan
\newtcolorbox{objectifs}{popbase,colframe=popink,colback=poporangeL,
  before upper={\popboxhead{Objectifs de la leçon}{popink}{}}}
\newtcolorbox{prerequis}{popbase,colframe=popink,colback=popblueL,
  before upper={\popboxhead{Prérequis}{popblue}{}}}
\newtcolorbox{bilan}{popbase,colframe=popink,colback=white,boxrule=2.8pt,
  before upper={\popboxhead{Fiche bilan}{poporange}{L'essentiel à connaître pour l'examen}}}
% Figure encadrée avec légende
\newtcolorbox{popfigure}[1][]{enhanced,breakable=false,colback=white,colframe=popline,boxrule=1.4pt,arc=8pt,
  left=10pt,right=10pt,top=10pt,bottom=8pt,before skip=10pt,after skip=12pt,halign=center,
  lower separated=false,halign lower=center,
  fontlower=\popsemi\fontsize{9}{11.5}\selectfont\color{popmuted},
  #1}
\newcommand{\legende}[1]{\par\vspace{6pt}{\popsemi\fontsize{9}{11.5}\selectfont\color{popmuted}#1}}

% ============================================================
% Signature OnBuch+
% ============================================================
\newcommand{\poplogoinline}[1]{{\poplogo\fontsize{#1}{#1}\selectfont\color{popink}OnBuch\color{poporange}+}}
\newcommand{\signatureonbuch}{%
  \begin{tcolorbox}[enhanced,colback=popink,colframe=popink,boxrule=2.2pt,arc=10pt,
      drop shadow=poporange,left=14pt,right=14pt,top=10pt,bottom=10pt,before skip=0pt,after skip=0pt]
    \tikz[baseline=-0.7ex]\node[circle,fill=popgreen,draw=popcream,line width=1.3pt,minimum size=22pt,inner sep=0]{\popcheck{white}};%
    \hspace{9pt}\begin{minipage}[c]{0.62\linewidth}
      {\popxbold\fontsize{12}{13}\selectfont\color{popcream}Signé\ \textbullet\ {\poplogo OnBuch\color{poporange}+}}\par\vspace{2pt}
      {\raggedright\popsemi\fontsize{8}{10.5}\selectfont\color{popline}\MakeUppercase{Équipe pédagogique OnBuch+}\\\MakeUppercase{Conforme au programme officiel MINESEC}\par}
    \end{minipage}\hfill
    {\poplogo\fontsize{22}{22}\selectfont\color{popcream}OnBuch\color{poporange}+}
  \end{tcolorbox}%
}

% ============================================================
% Page de garde
% #1 titre · #2 matière · #3 niveau · #4 module · #5 n° leçon · #6 durée ·
% #7 description / objectifs (texte ou itemize)
% ============================================================
\newcommand{\pagegarde}[7]{%
  \thispagestyle{empty}
  \newgeometry{a4paper,margin=1.7cm,top=1.6cm,bottom=1.4cm}%
  \begin{tikzpicture}[remember picture,overlay]
    \fill[poporange!18] ([xshift=-3.2cm,yshift=-3.6cm]current page.north east) circle (4.6cm);
    \fill[poppurple!14] ([xshift=2.4cm,yshift=7.6cm]current page.south west) circle (3.8cm);
  \end{tikzpicture}%
  {\poplogo\fontsize{30}{30}\selectfont\color{popink}OnBuch\color{poporange}+}\hfill
  \raisebox{0.6ex}{\poppill{Cours signé \textbullet\ Premium}{popink}{popcream}}\par
  \vspace{1pt}\popsquiggle{poppurple}\par
  \vspace{8pt}
  \begin{tcolorbox}[enhanced,colback=poporange,colframe=popink,boxrule=2.8pt,arc=13pt,
      drop shadow=popink,left=18pt,right=18pt,top=12pt,bottom=13pt,after skip=10pt]
    \poppill{#2 \textbullet\ #3}{popink}{popcream}\hspace{5pt}\poppill{#4}{white}{popink}\par\vspace{8pt}
    {\popdisplay\fontsize{14}{14}\selectfont\color{popink}LEÇON #5}\par\vspace{4pt}
    {\raggedright\hyphenpenalty=10000\exhyphenpenalty=10000\popdisplay\fontsize{25}{27}\selectfont\color{popcream}\MakeUppercase{#1}\par}
    \vspace{8pt}
    {\popxbold\fontsize{9.5}{9.5}\selectfont\color{popink}\MakeUppercase{Durée conseillée : #6}}
  \end{tcolorbox}
  \begin{tcolorbox}[enhanced,colback=white,colframe=popink,boxrule=2.2pt,arc=10pt,
      drop shadow=popink,left=16pt,right=16pt,top=10pt,bottom=10pt,after skip=8pt]
    \poppill{Dans cette leçon}{poppurple}{white}\par\vspace{5pt}
    \popsemi\fontsize{9.3}{12.6}\selectfont\color{popink}#7
  \end{tcolorbox}
  \noindent\begin{minipage}[t]{0.487\linewidth}
    \begin{tcolorbox}[enhanced,colback=popgreen,colframe=popink,boxrule=2.2pt,arc=10pt,
        drop shadow=popink,left=11pt,right=11pt,top=10pt,bottom=10pt,height=3.1cm,valign=center]
      \tcbox[colback=white,colframe=popink,boxrule=1.3pt,arc=5pt,boxsep=0pt,nobeforeafter,
        left=3pt,right=3pt,top=3pt,bottom=3pt,valign=center]{\qrcode[height=1.9cm]{https://play.google.com/store/apps/details?id=com.luvvix.onbuch&hl=fr}}%
      \hspace{8pt}\begin{minipage}[c]{0.5\linewidth}
        {\popdisplay\fontsize{11}{12.5}\selectfont\color{white}\MakeUppercase{Télécharge OnBuch}}\par\vspace{4pt}
        {\raggedright\popsemi\fontsize{8.4}{11}\selectfont\color{white}Gratuit \textbullet\ Google Play\\Tuteur IA, annales, concours, résultats\par}
      \end{minipage}
    \end{tcolorbox}
  \end{minipage}\hfill
  \begin{minipage}[t]{0.487\linewidth}
    \begin{tcolorbox}[enhanced,colback=poppurple,colframe=popink,boxrule=2.2pt,arc=10pt,
        drop shadow=popink,left=11pt,right=11pt,top=10pt,bottom=10pt,height=3.1cm,valign=center]
      \tikz[baseline=-0.7ex]\node[circle,fill=white,draw=popink,line width=1.3pt,minimum size=1.6cm,inner sep=0]{\popdisplay\fontsize{24}{24}\selectfont\color{poppurple}+};%
      \hspace{8pt}\begin{minipage}[c]{0.6\linewidth}
        {\popdisplay\fontsize{11}{12.5}\selectfont\color{white}\MakeUppercase{Passe à OnBuch+}}\par\vspace{4pt}
        {\raggedright\popsemi\fontsize{8.4}{11}\selectfont\color{white}Tous les cours signés, Tuteur IA illimité, fiches PDF — onglet OnBuch+ de l'app\par}
      \end{minipage}
    \end{tcolorbox}
  \end{minipage}
  \vspace{8pt}
  \popcontenu
  \vfill
  \signatureonbuch
  \restoregeometry
  \newpage
}

% Rangée « Ce cours contient » (page de garde)
\newcommand{\poptile}[3]{% #1 couleur #2 gros texte #3 légende
  \begin{tcolorbox}[enhanced,colback=#1,colframe=popink,boxrule=2pt,arc=9pt,shadow={2.5pt}{-2.5pt}{0pt}{popink},
      left=6pt,right=6pt,top=9pt,bottom=9pt,halign=center,valign=center,height=1.75cm,width=\linewidth,nobeforeafter]
    {\popdisplay\fontsize{12.5}{13}\selectfont\color{white}\MakeUppercase{#2}}\par\vspace{3pt}
    {\centering\popsemi\fontsize{7.8}{9.5}\selectfont\color{white}#3\par}
  \end{tcolorbox}}
\newcommand{\popcontenu}{%
  \par\noindent{\popxboldls\fontsize{7.5}{7.5}\selectfont\color{popmuted}CE COURS SIGNÉ CONTIENT}\par\vspace{7pt}
  \noindent\begin{minipage}[t]{0.235\linewidth}\poptile{popblue}{Cours}{complet, illustré, pas à pas}\end{minipage}\hfill
  \begin{minipage}[t]{0.235\linewidth}\poptile{poporange}{Méthodes}{et exercices résolus}\end{minipage}\hfill
  \begin{minipage}[t]{0.235\linewidth}\poptile{poppink}{Exercices}{type examen + corrigés}\end{minipage}\hfill
  \begin{minipage}[t]{0.235\linewidth}\poptile{popgreen}{Fiche bilan}{l'essentiel à retenir}\end{minipage}\par}

% Sommaire
\newtcolorbox{sommaire}{enhanced,colback=white,colframe=popink,boxrule=2.2pt,arc=10pt,
  drop shadow=popink,left=18pt,right=18pt,top=13pt,bottom=13pt,before skip=6pt,after skip=20pt,
  before upper={\poppill{Sommaire}{poppurple}{white}\par\vspace{9pt}\popsemi\fontsize{10.6}{14.5}\selectfont\color{popink}}}

% Signature de fin de cours — poussée tout en bas de la dernière page
\newcommand{\findecours}{%
  \par\vfill\nopagebreak
  \begin{center}\popsquiggle{poporange}\end{center}\vspace{4pt}
  \signatureonbuch
}
\AtBeginDocument{\def\llbracket{\mathopen{[\mkern-3.5mu[}}\def\rrbracket{\mathclose{]\mkern-3.5mu]}}} % intervalles d'entiers (glyphe ⟦ absent de la police)
% Glyphes absents de Fira Math : redéfinis à partir de glyphes disponibles
\AtBeginDocument{%
  \def\setminus{\mathbin{\backslash}}\let\smallsetminus\setminus
  \def\vdots{\mathord{\vbox{\baselineskip3.2pt\lineskiplimit0pt\kern4pt\hbox{.}\hbox{.}\hbox{.}}}}%
  \def\ddots{\mathinner{\mkern1mu\raise6.5pt\hbox{.}\mkern2mu\raise3.6pt\hbox{.}\mkern2mu\raise0.7pt\hbox{.}\mkern1mu}}%
  \def\triangle{\mathord{\scalebox{0.9}{$\symup{\Delta}$}}}%
  \def\nabla{\mathord{\raisebox{\depth}{\scalebox{1}[-1]{$\symup{\Delta}$}}}}%
  \def\checkmark{\popcheck{popdark}}%
  \def\star{\mathbin{\ast}}\def\bigstar{\mathord{\ast}}%
  \def\diamond{\mathbin{\scalebox{0.6}{$\Diamond$}}}%
  \def\bigcup{\mathop{\mathchoice{\raisebox{-0.3em}{\scalebox{1.7}{$\cup$}}}{\scalebox{1.25}{$\cup$}}{\cup}{\cup}}\displaylimits}%
  \def\bigcap{\mathop{\mathchoice{\raisebox{-0.3em}{\scalebox{1.7}{$\cap$}}}{\scalebox{1.25}{$\cap$}}{\cap}{\cap}}\displaylimits}%
  \def\lesseqqgtr{\mathrel{\lessgtr}}\def\bigoplus{\mathop{\oplus}\displaylimits}%
}
\providecommand{\ssi}{si et seulement si} % abréviation fréquente
\AtBeginDocument{\providecommand{\wideparen}{\overparen}} % arc de cercle

%% FIGFIX-BEGIN v5 — correctifs globaux des figures (docs/onbuch-generation/tools/figfix)
\usepackage{adjustbox}\usepackage{etoolbox}\tcbuselibrary{raster}
% toute figure TikZ/pgfplots plus large que la ligne est réduite à la largeur disponible
\BeforeBeginEnvironment{tikzpicture}{\begin{adjustbox}{max width=\linewidth}}
\AfterEndEnvironment{tikzpicture}{\end{adjustbox}}
% légendes pgfplots lisibles par-dessus les courbes
\pgfplotsset{every axis/.append style={legend style={fill=white,fill opacity=0.92,text opacity=1,draw=black!15,inner sep=3pt}}}
% étiquettes d'arêtes (syntaxe edge["..."]) avec fond blanc
\tikzset{every edge quotes/.append style={fill=white,inner sep=1.2pt}}
%% FIGFIX-END
\begin{document}

\pagegarde{Sensibilisation sur la nécessité des deux phénomènes biologiques complémentaires (méiose et fécondation) pour la pérennité de l'espèce}{SVT}{Tle D}{Module I : Le Monde Vivant — Reproduction sexuée chez les Mammifères et les Spermaphytes}{N03}{10 h}{Cette leçon établit, à partir de faits expérimentaux et d'observations microscopiques, les mécanismes de la méiose, de la gamétogenèse et de la fécondation chez les Mammifères et les Spermaphytes. Elle montre comment la complémentarité de la méiose (réduction chromosomique, brassage) et de la fécondation (restauration de la diploïdie) assure à la fois la stabilité et la diversité génétiques, conditions de la pérennité de l'espèce.
\vspace{4pt}
\begin{multicols}{2}\small
\begin{itemize}
\item À la fin de cette leçon, je sais décrire la structure et le rôle des gonades (testicules, ovaires ; étamines, pistil) chez les Mammifères et les Spermaphytes.
\item À la fin de cette leçon, je sais réaliser la dissection d'un Mammifère et d'une fleur, mettre en évidence les organes reproducteurs et en réaliser des schémas légendés.
\item À la fin de cette leçon, je sais décrire le comportement des chromosomes au cours des deux divisions de la méiose et reconnaître ses étapes sur des électronographies.
\item À la fin de cette leçon, je sais interpréter la courbe d'évolution de la quantité d'ADN au cours de la méiose et de la gamétogenèse.
\item À la fin de cette leçon, je sais identifier les cellules de la lignée germinale (spermatogonies, spermatocytes, spermatides, ovogonies, ovocytes) à leurs différents stades.
\item À la fin de cette leçon, je sais décrire et comparer les étapes de la fécondation chez les Mammifères et chez les Spermaphytes, y compris la double fécondation.
\end{itemize}\end{multicols}}

\begin{sommaire}
\begin{multicols}{2}
\begin{enumerate}
\item Les organes de la reproduction : gonades des Mammifères et des Spermaphytes
\item La méiose : comportement des chromosomes
\item Évolution de la quantité d'ADN au cours de la méiose et de la gamétogenèse
\item La gamétogenèse chez les Mammifères et les Spermaphytes
\item La fécondation chez les Mammifères et la double fécondation chez les Spermaphytes
\item Méiose et fécondation : deux phénomènes complémentaires au service de la pérennité de l'espèce
\item Activité d'intégration
\item Méthodes et exercices résolus
\item Exercices
\item Corrigés des exercices
\item Fiche bilan
\end{enumerate}
\end{multicols}
\end{sommaire}

\begin{objectifs}
\begin{itemize}
\item À la fin de cette leçon, je sais décrire la structure et le rôle des gonades (testicules, ovaires ; étamines, pistil) chez les Mammifères et les Spermaphytes.
\item À la fin de cette leçon, je sais réaliser la dissection d'un Mammifère et d'une fleur, mettre en évidence les organes reproducteurs et en réaliser des schémas légendés.
\item À la fin de cette leçon, je sais décrire le comportement des chromosomes au cours des deux divisions de la méiose et reconnaître ses étapes sur des électronographies.
\item À la fin de cette leçon, je sais interpréter la courbe d'évolution de la quantité d'ADN au cours de la méiose et de la gamétogenèse.
\item À la fin de cette leçon, je sais identifier les cellules de la lignée germinale (spermatogonies, spermatocytes, spermatides, ovogonies, ovocytes) à leurs différents stades.
\item À la fin de cette leçon, je sais décrire et comparer les étapes de la fécondation chez les Mammifères et chez les Spermaphytes, y compris la double fécondation.
\item À la fin de cette leçon, je sais interpréter la courbe d'évolution de la quantité d'ADN au cours de la fécondation.
\item À la fin de cette leçon, je sais expliquer comment la complémentarité méiose-fécondation assure la stabilité du stock chromosomique et la diversité génétique, donc la pérennité de l'espèce.
\end{itemize}
\end{objectifs}

\begin{prerequis}
\begin{itemize}
\item Structure du chromosome et notion de caryotype (2n, n) vues en classe de Seconde/Première
\item La mitose et son rôle dans la croissance (Première)
\item Notion de reproduction sexuée et d'organes reproducteurs (cycle 3 et Seconde)
\item L'ADN comme support de l'information génétique (Première)
\item Lecture et exploitation de courbes et de documents scientifiques (transversal)
\end{itemize}
\end{prerequis}

\newpage

\coursec{Les organes de la reproduction : gonades des Mammifères et des Spermaphytes}

\textbf{Situation d'accroche.} Au marché de Mokolo à Yaoundé, une vendeuse de plantains explique que ses plants donnent toujours des fruits identiques, tandis que le maïs semé à côté produit des épis très différents d'une année à l'autre. Dans un village de l'Adamaoua, un éleveur constate que ses zébus croisés avec une race améliorée donnent des veaux tous différents, mais toujours reconnaissables comme des bovins. Comment la reproduction sexuée parvient-elle à la fois à \emph{conserver} les caractères de l'espèce et à \emph{produire} des individus tous différents ? Pourquoi la nature a-t-elle « inventé » deux phénomènes aussi complexes que la méiose et la fécondation ? C'est à cette double énigme, essentielle pour la pérennité des espèces, que cette leçon va répondre.

\textbf{Problématique.} Pour comprendre la reproduction sexuée, il faut d'abord identifier les organes qui produisent les cellules reproductrices. Où et comment les gamètes sont-ils fabriqués chez les Mammifères et chez les Spermaphytes ? Quelles structures assurent cette production ?

\courssub{La reproduction sexuée : définition et cycle diploïde dominant}

\begin{definition}[Reproduction sexuée]
La \cle{reproduction sexuée} est un mode de reproduction qui met en jeu l'union de deux cellules spécialisées, les \cle{gamètes} (un gamète mâle et un gamète femelle), au cours d'un phénomène appelé \cle{fécondation}. Il en résulte une cellule unique, la \cle{cellule-œuf} (ou \cle{zygote}), qui se développe en un nouvel individu.
\end{definition}

Chez les Mammifères et les Spermaphytes, l'individu adulte est \cle{diploïde} : chacune de ses cellules possède $2n$ chromosomes, c'est-à-dire $n$ paires de chromosomes \emph{homologues} (un chromosome de chaque paire vient du père, l'autre de la mère). Chez l'Homme, $2n = 46$ chromosomes ; chez le maïs, $2n = 20$ ; chez le riz, $2n = 24$. On dit que ces espèces présentent un \cle{cycle diploïde dominant} : seuls les gamètes sont \cle{haploïdes} ($n$ chromosomes).

\begin{attention}[Le problème du maintien du nombre de chromosomes]
Si les gamètes étaient produits par de simples mitoses, ils seraient diploïdes ($2n$) : à chaque génération, la fécondation doublerait le nombre de chromosomes ($2n + 2n = 4n$, puis $8n$\ldots), ce qui est incompatible avec la stabilité de l'espèce. Il faut donc une division spéciale qui \emph{réduise de moitié} le nombre de chromosomes : c'est la \cle{méiose}. La fécondation, elle, \emph{restaure} la diploïdie. Méiose et fécondation sont donc deux phénomènes \emph{complémentaires} : c'est le fil conducteur de toute la leçon.
\end{attention}

\begin{exemplebox}[Deux modes de reproduction au marché]
Le plantain se multiplie par \emph{reproduction asexuée} (rejets, boutures) : les descendants sont génétiquement identiques au parent. Le maïs et les zébus se reproduisent par \emph{reproduction sexuée} : le brassage génétique lié à la méiose et à la fécondation produit des individus tous différents, tout en maintenant le nombre de chromosomes et les caractères généraux de l'espèce.
\end{exemplebox}

\courssub{Les organes reproducteurs du Mammifère mâle : dissection du rat}

\begin{experience}[Dissection de l'appareil génital du rat mâle]
\textbf{Matériel :} rat mâle adulte fraîchement sacrifié (ou conservé au formol), plateau de dissection, ciseaux fins, pinces, aiguilles à étaler, loupe binoculaire.

\textbf{Protocole :}
\begin{enumerate}
\item Fixer l'animal sur le dos, membres écartés.
\item Pratiquer une incision médiane de la peau de l'abdomen, puis de la paroi musculaire, et écarter les replis.
\item Refouler délicatement les anses intestinales vers le haut pour dégager la région postérieure de la cavité abdominale.
\item Observer à la loupe les organes pairs situés près des reins et dans la région pelvienne.
\end{enumerate}

\textbf{Observations :} on distingue deux \cle{testicules}, petits organes ovoïdes blanchâtres, logés dans des poches cutanées appelées \cle{bourses} (scrotum), à l'extérieur de l'abdomen. Chaque testicule est coiffé par un organe allongé et sinueux, l'\cle{épididyme}, qui se prolonge par un tube fin et blanchâtre, le \cle{canal déférent}, remontant vers l'urètre. On observe aussi des glandes annexes (vésicules séminales, prostate) à la base de la vessie.

\textbf{Interprétation :} la position externe des testicules s'explique : la production des spermatozoïdes exige une température d'environ $2\ ^\circ\mathrm{C}$ inférieure à la température interne du corps (environ $37\ ^\circ\mathrm{C}$ chez l'Homme). L'épididyme est le lieu de \emph{maturation} et de stockage des spermatozoïdes ; le canal déférent assure leur transport lors de l'éjaculation.
\end{experience}

\begin{popfigure}
\begin{tikzpicture}[scale=0.92, every node/.style={font=\small}]
% ---- MALE ----
\begin{scope}
\node[font=\bfseries] at (2.6,6.6) {APPAREIL GÉNITAL MÂLE (vue latérale)};
% vessie
\draw[thick, popblue] (1.2,4.6) ellipse (0.9 and 0.7);
% testicule
\draw[thick, poporange, fill=poporangeL] (4.6,1.1) ellipse (0.75 and 0.55);
% épididyme
\draw[thick, popgreen, decorate, decoration={coil, amplitude=1.5pt, segment length=3pt}] (3.9,1.15) .. controls (3.7,1.9) and (4.6,1.9) .. (5.3,1.3);
% canal déférent
\draw[thick, popgreen] (5.3,1.3) .. controls (5.9,2.4) and (4.4,3.6) .. (2.4,4.1);
% urètre
\draw[thick, poppurple] (1.9,4.0) .. controls (2.6,3.0) and (3.4,2.6) .. (5.6,2.2);
\draw[thick, poppurple] (5.6,2.2) -- (6.3,2.1);
% légendes
\draw[->, popmuted] (1.2,5.4) -- (1.2,5.9) node[above]{1. Vessie};
\draw[->, popmuted] (4.6,0.55) -- (4.6,0.1) node[below]{2. Testicule (dans les bourses)};
\draw[->, popmuted] (5.2,1.7) -- (6.4,1.3) node[right]{3. Épididyme};
\draw[->, popmuted] (4.6,3.2) -- (6.4,3.4) node[right]{4. Canal déférent};
\draw[->, popmuted] (6.0,2.15) -- (6.4,2.5) node[right]{5. Urètre};
\end{scope}
% ---- FEMELLE ----
\begin{scope}[xshift=8.6cm]
\node[font=\bfseries] at (2.6,6.6) {APPAREIL GÉNITAL FEMELLE (vue dorsale)};
% ovaires
\draw[thick, poporange, fill=poporangeL] (0.9,4.9) ellipse (0.45 and 0.35);
\draw[thick, poporange, fill=poporangeL] (4.3,4.9) ellipse (0.45 and 0.35);
% trompes
\draw[thick, popgreen, decorate, decoration={coil, amplitude=1.2pt, segment length=3pt}] (1.3,4.8) .. controls (1.9,4.3) .. (2.2,3.9);
\draw[thick, popgreen, decorate, decoration={coil, amplitude=1.2pt, segment length=3pt}] (3.9,4.8) .. controls (3.3,4.3) .. (3.0,3.9);
% cornes utérines
\draw[thick, popblue] (2.2,3.9) .. controls (2.3,3.2) .. (2.5,2.6);
\draw[thick, popblue] (3.0,3.9) .. controls (2.9,3.2) .. (2.7,2.6);
% corps utérus
\draw[thick, popblue] (2.5,2.6) -- (2.55,1.7);
\draw[thick, popblue] (2.7,2.6) -- (2.65,1.7);
\draw[thick, popblue] (2.55,1.7) -- (2.65,1.7);
% vagin
\draw[thick, poppurple] (2.55,1.7) -- (2.55,0.7);
\draw[thick, poppurple] (2.65,1.7) -- (2.65,0.7);
\draw[thick, poppurple] (2.55,0.7) .. controls (2.6,0.55) .. (2.65,0.7);
% légendes
\draw[->, popmuted] (0.9,5.25) -- (0.9,5.8) node[above]{1. Ovaire};
\draw[->, popmuted] (1.6,4.5) -- (0.6,3.6) node[left]{2. Trompe (oviducte)};
\draw[->, popmuted] (2.35,3.2) -- (0.6,2.6) node[left]{3. Corne de l'utérus};
\draw[->, popmuted] (2.6,1.4) -- (4.4,1.4) node[right]{4. Corps de l'utérus};
\draw[->, popmuted] (2.6,0.7) -- (4.4,0.6) node[right]{5. Vagin};
\end{scope}
\end{tikzpicture}
\legende{Schémas comparatifs des appareils génitaux du rat. À gauche : appareil génital mâle (1. vessie ; 2. testicule ; 3. épididyme ; 4. canal déférent ; 5. urètre). À droite : appareil génital femelle (1. ovaire ; 2. trompe ; 3. corne utérine ; 4. corps de l'utérus ; 5. vagin). Les gonades (testicules et ovaires) sont en orange.}
\end{popfigure}

\courssub{Les organes reproducteurs du Mammifère femelle}

La dissection d'une rate (ou l'observation d'un modèle anatomique) met en évidence deux \cle{ovaires}, petits organes ovoïdes situés dans la cavité abdominale, près des reins. Chaque ovaire est voisin de l'extrémité élargie en pavillon d'un conduit fin, la \cle{trompe} (ou oviducte), qui s'ouvre dans l'\cle{utérus}. Chez la rate, l'utérus comporte deux longues \emph{cornes utérines} (utérus \emph{bicorne}, adapté aux portées nombreuses) ; chez la femme, l'utérus est unique (utérus \emph{simple}). L'utérus débouche dans le \cle{vagin}.

\begin{aretenir}[Rôles des organes reproducteurs des Mammifères]
\begin{itemize}
\item \textbf{Testicules} (gonades mâles) : production des \cle{spermatozoïdes} et sécrétion d'hormones mâles (testostérone).
\item \textbf{Ovaires} (gonades femelles) : production des \cle{ovocytes} (gamètes femelles) et sécrétion d'hormones femelles (œstrogènes, progestérone).
\item \textbf{Voies génitales} (épididymes, canaux déférents ; trompes, utérus, vagin) : transport, maturation, rencontre des gamètes et, chez la femelle, développement de l'embryon.
\end{itemize}
\end{aretenir}

\courssub{Observation histologique : où naissent les gamètes ?}

La dissection montre \emph{où} sont les gonades, mais pas \emph{où} se forment les gamètes. Pour cela, il faut observer des coupes fines au microscope optique.

\begin{experience}[Observation d'une coupe de testicule]
\textbf{Protocole :} observer au microscope (objectif $\times 40$) une coupe histologique de testicule de rat colorée à l'hématéine-éosine.

\textbf{Observations :} le testicule est formé d'un enchevêtrement de tubes enroulés, les \cle{tubes séminifères}. En coupe transversale, chaque tube montre une \emph{gradation remarquable} des cellules, de la paroi vers la lumière centrale :
\begin{itemize}
\item contre la membrane basale : de grandes cellules, les \cle{spermatogonies} ($2n$) ;
\item plus près de la lumière : les \cle{spermatocytes} (I puis II) ;
\item encore plus près : de petites cellules rondes, les \cle{spermatides} ($n$) ;
\item dans la lumière : des \cle{spermatozoïdes} ($n$), cellules flagellées libres.
\end{itemize}

\textbf{Interprétation :} cette gradation traduit une \emph{différenciation progressive} : les gamètes mâles se forment à partir des cellules de la paroi et migrent vers la lumière au fur et à mesure de leur maturation. Le tube séminifère est donc le \emph{lieu de production des spermatozoïdes}.
\end{experience}

\begin{popfigure}
\begin{tikzpicture}[scale=1.0]
% paroi du tube
\draw[thick, popblue, fill=popblueL] (0,0) circle (3.4);
\draw[thick, popdark, fill=popcream] (0,0) circle (2.9);
% lumière
\draw[thick, popmuted, fill=white] (0,0) circle (0.9);
% spermatogonies (contre la paroi)
\foreach \a in {0,30,...,330}{
  \draw[fill=popgreenL, draw=popgreen] (\a:2.72) circle (0.16);
}
% spermatocytes
\foreach \a in {15,60,105,150,195,240,285,330}{
  \draw[fill=poporangeL, draw=poporange] (\a:2.25) circle (0.19);
}
% spermatides
\foreach \a in {0,45,...,315}{
  \draw[fill=poppinkL, draw=poppink] (\a:1.65) circle (0.11);
}
% spermatozoïdes dans la lumière
\foreach \a/\r in {20/0.45, 110/0.5, 200/0.4, 290/0.55, 70/0.3, 250/0.25}{
  \draw[fill=popdark, draw=popdark] (\a:\r) circle (0.06);
  \draw[popdark] (\a:\r) -- ({\a+40}:{\r+0.35});
}
% légendes
\draw[->, popmuted] (150:2.72) -- (-4.6,2.6) node[left]{1. Spermatogonies ($2n$)};
\draw[->, popmuted] (105:2.25) -- (-1.6,4.1) node[above]{2. Spermatocytes};
\draw[->, popmuted] (45:1.65) -- (3.6,3.0) node[right]{3. Spermatides ($n$)};
\draw[->, popmuted] (20:0.5) -- (3.6,0.9) node[right]{4. Spermatozoïdes ($n$)};
\draw[->, popmuted] (-30:3.15) -- (3.6,-2.2) node[right]{5. Paroi du tube séminifère};
\node at (0,-4.0) {};
\end{tikzpicture}
\legende{Coupe transversale schématique d'un tube séminifère de Mammifère. La gradation des cellules de la paroi (5) vers la lumière centrale traduit les étapes de la spermatogenèse : spermatogonies (1), spermatocytes (2), spermatides (3), spermatozoïdes flagellés (4) libérés dans la lumière.}
\end{popfigure}

De même, l'observation d'une coupe d'\cle{ovaire} de Mammifère montre des structures sphériques de tailles variées, les \cle{follicules ovariens}. Chaque follicule contient une grosse cellule, l'\cle{ovocyte}, entourée de cellules nourricières. Les follicules sont à différents stades de maturation (follicules primordiaux, follicules en croissance, follicule mûr ou \emph{follicule de De Graaf}) ; au moment de l'\cle{ovulation}, le follicule mûr libère l'ovocyte dans le pavillon de la trompe. L'ovaire est donc le \emph{lieu de production des gamètes femelles}.

\begin{savaistu}[Une production très inégale]
Chez l'Homme, les testicules produisent environ $100$ millions de spermatozoïdes \emph{par jour}, soit plus de $1\,000$ par seconde, et ce de la puberté jusqu'à un âge avancé. À l'inverse, la femme naît avec un stock déjà constitué d'environ $1$ à $2$ millions de follicules, dont seulement $400$ à $500$ arriveront à maturation : elle ne libère en général qu'\emph{un seul ovocyte par cycle} (environ tous les $28$ jours), de la puberté à la ménopause. Le gamète mâle est minuscule (tête d'environ $5\ \mu\mathrm{m}$ de long) et mobile ; l'ovocyte est l'une des plus grosses cellules du corps humain (environ $120\ \mu\mathrm{m}$ de diamètre), riche en réserves.
\end{savaistu}

\courssub{Les organes reproducteurs de la fleur de Spermaphyte}

\begin{experience}[Dissection d'une fleur d'hibiscus]
\textbf{Matériel :} fleur d'hibiscus (ou de pois de senteur), pince fine, aiguille, lame de rasoir, loupe binoculaire, microscope optique.

\textbf{Protocole :}
\begin{enumerate}
\item Observer la fleur entière et identifier de l'extérieur vers l'intérieur : le pédoncule, les sépales (verts), les pétales (colorés).
\item Écarter la corolle et dégager la colonne centrale.
\item Prélever une étamine et observer l'anthère à la loupe, puis écraser un grain de pollen entre lame et lamelle pour l'observer au microscope.
\item Pratiquer une coupe longitudinale de l'ovaire et observer les ovules à la loupe.
\end{enumerate}

\textbf{Observations :}
\begin{itemize}
\item Les \cle{étamines}, organes reproducteurs \emph{mâles}, sont formées chacune d'un pédicule fin, le \cle{filet}, surmonté d'un renflement, l'\cle{anthère}. L'anthère contient des loges polliniques remplies de \cle{grains de pollen}, petites sphères ornementées visibles au microscope.
\item Le \cle{pistil}, organe reproducteur \emph{femelle} situé au centre, comprend trois parties : le \cle{stigmate} (extrémité supérieure, souvent collante, qui reçoit le pollen), le \cle{style} (colonne reliant le stigmate à l'ovaire) et l'\cle{ovaire} (partie renflée à la base).
\item La coupe de l'ovaire montre qu'il contient de petites structures ovoïdes fixées au placenta : les \cle{ovules}. Chaque ovule renferme un \emph{sac embryonnaire} contenant notamment l'\cle{oosphère} (le véritable gamète femelle).
\end{itemize}

\textbf{Interprétation :} chez les Spermaphytes, les gonades mâles sont les \emph{anthères} (elles produisent les grains de pollen, qui élaboreront les gamètes mâles) et la gonade femelle est l'\emph{ovaire de la fleur} (ses ovules contiennent les gamètes femelles).
\end{experience}

\begin{popfigure}
\begin{tikzpicture}[scale=1.05, every node/.style={font=\small}]
% pédoncule
\draw[thick, popgreen] (0,-2.6) -- (0,-1.2);
% réceptacle
\draw[thick, popgreen, fill=popgreenL] (-0.5,-1.2) .. controls (-0.5,-0.9) and (0.5,-0.9) .. (0.5,-1.2) -- cycle;
% sépales
\draw[thick, popgreen, fill=popgreenL] (-0.5,-1.1) .. controls (-1.6,-0.6) and (-1.9,0.2) .. (-1.7,0.7) .. controls (-1.2,0.1) and (-0.7,-0.5) .. (-0.3,-1.0) -- cycle;
\draw[thick, popgreen, fill=popgreenL] (0.5,-1.1) .. controls (1.6,-0.6) and (1.9,0.2) .. (1.7,0.7) .. controls (1.2,0.1) and (0.7,-0.5) .. (0.3,-1.0) -- cycle;
% pétales
\draw[thick, poppink, fill=poppinkL] (-0.4,-0.9) .. controls (-2.3,0.3) and (-2.6,1.8) .. (-2.1,2.6) .. controls (-1.3,1.6) and (-0.6,0.4) .. (-0.15,-0.7) -- cycle;
\draw[thick, poppink, fill=poppinkL] (0.4,-0.9) .. controls (2.3,0.3) and (2.6,1.8) .. (2.1,2.6) .. controls (1.3,1.6) and (0.6,0.4) .. (0.15,-0.7) -- cycle;
% pistil : ovaire
\draw[thick, poporange, fill=poporangeL] (0,-0.75) ellipse (0.42 and 0.5);
% ovules dans l'ovaire
\draw[fill=popgoldL, draw=popgold] (-0.15,-0.7) circle (0.1);
\draw[fill=popgoldL, draw=popgold] (0.15,-0.85) circle (0.1);
\draw[fill=popgoldL, draw=popgold] (0.02,-0.5) circle (0.1);
% style
\draw[thick, poporange] (-0.07,-0.25) -- (-0.07,2.2);
\draw[thick, poporange] (0.07,-0.25) -- (0.07,2.2);
% stigmate
\draw[thick, poporange, fill=poporangeL] (0,2.35) circle (0.18);
% étamines
\draw[thick, popblue] (-0.35,-0.9) .. controls (-1.0,0.6) .. (-1.15,1.7);
\draw[thick, popblue, fill=popblueL] (-1.15,1.95) ellipse (0.13 and 0.28);
\draw[thick, popblue] (0.35,-0.9) .. controls (1.0,0.6) .. (1.15,1.7);
\draw[thick, popblue, fill=popblueL] (1.15,1.95) ellipse (0.13 and 0.28);
% légendes
\draw[->, popmuted] (0,2.53) -- (0,3.1) node[above]{1. Stigmate};
\draw[->, popmuted] (0.07,1.4) -- (2.9,1.5) node[right]{2. Style};
\draw[->, popmuted] (0.42,-0.6) -- (2.9,-0.4) node[right]{3. Ovaire};
\draw[->, popmuted] (0.15,-0.85) -- (2.9,-1.3) node[right]{4. Ovule};
\draw[->, popmuted] (1.15,2.23) -- (2.9,2.6) node[right]{5. Anthère};
\draw[->, popmuted] (0.95,0.8) -- (2.9,0.7) node[right]{6. Filet};
\draw[->, popmuted] (-2.1,2.0) -- (-3.2,2.6) node[left]{7. Pétale};
\draw[->, popmuted] (-1.75,0.3) -- (-3.2,0.3) node[left]{8. Sépale};
\draw[->, popmuted] (0,-2.3) -- (-3.2,-2.3) node[left]{9. Pédoncule};
% accolades
\node[popblue, font=\footnotesize\bfseries] at (4.6,1.65) {étamine};
\node[poporange, font=\footnotesize\bfseries] at (4.6,-0.15) {pistil};
\end{tikzpicture}
\legende{Coupe longitudinale schématique d'une fleur de Spermaphyte. Le pistil (1. stigmate ; 2. style ; 3. ovaire ; 4. ovules) est l'organe reproducteur femelle ; les étamines (5. anthère ; 6. filet) sont les organes reproducteurs mâles. 7. pétale ; 8. sépale ; 9. pédoncule floral.}
\end{popfigure}

\begin{attention}[Ne pas confondre les deux « ovules »]
Erreur classique au Baccalauréat : le mot \emph{ovule} désigne deux réalités différentes.
\begin{itemize}
\item Chez les \textbf{Mammifères}, l'ovule (ou ovocyte) est le \emph{gamète femelle} lui-même, une cellule haploïde.
\item Chez les \textbf{Spermaphytes}, l'ovule est une \emph{structure pluricellulaire} contenue dans l'ovaire de la fleur ; il renferme le sac embryonnaire qui contient l'\emph{oosphère}, véritable gamète femelle. Après fécondation, l'ovule se transforme en \cle{graine} et l'ovaire en \cle{fruit}.
\end{itemize}
Dans une copie, précise toujours de quel « ovule » tu parles.
\end{attention}

\begin{methode}[Règles du schéma anatomique en SVT]
Un schéma d'observation noté au Baccalauréat doit respecter des règles strictes :
\begin{enumerate}
\item Dessiner au crayon à papier, à \emph{traits nets et continus}, sans ombres ni hachures fantaisistes ; les couleurs ne sont tolérées que si elles correspondent à l'observation réelle.
\item Donner au schéma un \emph{titre en majuscules}, souligné, précis (ex. : COUPE LONGITUDINALE D'UNE FLEUR D'HIBISCUS).
\item Tracer les traits de légende à la règle, \emph{horizontalement}, sans qu'ils se croisent, et écrire les légendes horizontalement, alignées de part et d'autre du schéma.
\item Respecter les proportions réelles des structures observées.
\item Ne jamais mélanger sur un même schéma des éléments vus à des grossissements différents ; indiquer le grossissement si l'énoncé le demande.
\end{enumerate}
\end{methode}

\courssub{Conclusion : les gonades, organes producteurs de gamètes}

\begin{definition}[Gonade]
Une \cle{gonade} est l'organe qui produit les gamètes. Chez les Mammifères, les gonades sont les \cle{testicules} (mâle) et les \cle{ovaires} (femelle). Chez les Spermaphytes, les gonades sont les \cle{anthères} (mâle) et l'\cle{ovaire de la fleur} (femelle). La formation des gamètes dans les gonades s'appelle la \cle{gamétogenèse} ; elle fait intervenir la méiose, que nous étudierons dans la prochaine section.
\end{definition}

\begin{center}
\begin{tabularx}{\linewidth}{|l|X|X|}
\hline
\rowcolor{popblueL} & \textbf{Mammifères} & \textbf{Spermaphytes} \\
\hline
Gonade mâle & Testicules (tubes séminifères) & Anthères (loges polliniques) \\
\hline
Gamète mâle & Spermatozoïde (mobile, flagellé) & Noyaux gamétiques du grain de pollen \\
\hline
Gonade femelle & Ovaires (follicules ovariens) & Ovaire de la fleur (ovules) \\
\hline
Gamète femelle & Ovocyte (gros, riche en réserves) & Oosphère (dans le sac embryonnaire) \\
\hline
\end{tabularx}
\end{center}

\begin{exoresolu}[Identifier les gonades et justifier]
\textbf{Énoncé.} Lors d'une séance de travaux pratiques au lycée de Ngaoundéré, un élève observe au microscope une coupe d'organe montrant de nombreux tubes dont la paroi présente des cellules de plus en plus petites vers la lumière, celle-ci contenant des cellules flagellées. (a) De quel organe s'agit-il et à quel groupe de vertébrés peut-il appartenir ? (b) Quelle est la fonction de cet organe ? (c) Quel équivalent fonctionnel cet organe possède-t-il chez une plante à fleurs ?
\tcblower
\textbf{Solution.} (a) Il s'agit d'une coupe de \emph{testicule} : l'organisation en tubes séminifères, la gradation des cellules de la paroi vers la lumière (spermatogonies $\rightarrow$ spermatocytes $\rightarrow$ spermatides) et la présence de spermatozoïdes flagellés dans la lumière sont caractéristiques. Cet organe appartient à un \emph{Mammifère mâle} (par exemple le rat). (b) Le testicule est la \emph{gonade mâle} : il produit les spermatozoïdes (gamètes mâles haploïdes) par spermatogenèse, et sécrète des hormones sexuelles mâles. (c) Chez les Spermaphytes, l'équivalent fonctionnel du testicule est l'\emph{anthère} de l'étamine, dont les loges polliniques produisent les grains de pollen porteurs des gamètes mâles.
\end{exoresolu}

\begin{exercice}[Pour t'entraîner]
(a) Réalise le schéma légendé de la coupe longitudinale d'une fleur de pois de senteur en respectant les règles du schéma anatomique. (b) Explique pourquoi, chez le rat mâle, les testicules sont situés dans les bourses et non dans l'abdomen. (c) Un élève écrit : « l'ovaire de la fleur produit les ovules, qui sont les gamètes femelles ». Corrige cette phrase.
\end{exercice}

\begin{corrige}[Exercice]
(a) Le schéma doit comporter un titre en majuscules, des traits nets, des légendes horizontales non croisées : stigmate, style, ovaire, ovule, anthère, filet, pétales, sépales, pédoncule. (b) La spermatogenèse exige une température d'environ $2\ ^\circ\mathrm{C}$ inférieure à la température abdominale ; la position externe des testicules dans les bourses assure ce refroidissement. (c) Phrase corrigée : « l'ovaire de la fleur contient les ovules ; chaque ovule renferme un sac embryonnaire qui contient l'oosphere, véritable gamète femelle. Après fécondation, l'ovule devient la graine et l'ovaire devient le fruit. »
\end{corrige}

\coursec{La méiose : comportement des chromosomes}

Dans la section précédente, nous avons décrit les gonades : testicules et ovaires chez les Mammifères, étamines et ovaire chez les Spermaphytes. Ces organes abritent les \cle{cellules de la lignée germinale}, qui produisent les gamètes. Mais une question fondamentale se pose immédiatement : comment une cellule diploïde peut-elle fabriquer des gamètes haploïdes ? C'est toute l'histoire de la \cle{méiose}, que nous allons maintenant raconter chromosome par chromosome.

\courssub{Le problème posé : pourquoi une division réductionnelle ?}

\begin{experience}[Un raisonnement par l'absurde]
\textbf{Hypothèse à tester.} Imaginons que les gamètes soient produits par \emph{mitose}, comme toutes les autres cellules du corps. Une mitose conserve le nombre de chromosomes : une cellule $2n$ donne deux cellules $2n$.

\textbf{Conséquence logique.} Les gamètes seraient alors diploïdes ($2n$). À la fécondation, la fusion de deux gamètes $2n$ donnerait un œuf $4n$ ; à la génération suivante, $8n$, puis $16n$... Le nombre de chromosomes \emph{doublerait à chaque génération}.

\textbf{Vérification expérimentale.} Or l'observation des caryotypes montre que le nombre de chromosomes d'une espèce est \emph{constant} d'une génération à l'autre : $2n = 46$ chez l'Homme, $2n = 8$ chez la Drosophile, $2n = 20$ chez le Maïs. L'hypothèse est donc \emph{fausse}.

\textbf{Conclusion.} Il doit exister, quelque part dans le cycle reproducteur, une division \emph{réductionnelle} qui divise par deux le stock chromosomique : c'est la \cle{méiose}.
\end{experience}

Ce raisonnement, simple mais puissant, est celui qu'August \textsc{Weismann} formula dès 1887, avant même que la méiose soit observée au microscope. L'observation de coupes de testicules de Mammifères (riches en cellules en division) et d'électronographies confirma ensuite l'existence de deux divisions successives au comportement chromosomique original.

\begin{definition}[La méiose]
La \cle{méiose} est une succession de \textbf{deux divisions cellulaires} précédées d'\textbf{une seule réplication de l'ADN}, qui fait passer une cellule mère diploïde ($2n$ chromosomes, chacun à deux chromatides après réplication) à \textbf{quatre cellules filles haploïdes} ($n$ chromosomes), génétiquement différentes entre elles et différentes de la cellule mère.
\end{definition}

\begin{attention}[Le piège du vocabulaire]
Ne confonds jamais \emph{réplication de l'ADN} (qui a lieu une seule fois, avant la méiose, en phase S) et \emph{divisions} (qui sont deux). C'est précisément le décalage « une réplication pour deux divisions » qui explique la réduction du stock chromosomique.
\end{attention}

\courssub{La première division : la division réductionnelle}

\textbf{Prophase I : l'appariement des homologues et les enjambements.}

C'est l'étape la plus longue et la plus riche en événements. Les chromosomes, déjà répliqués (chacun formé de \textbf{deux chromatides sœurs}), se condensent. Puis survient un phénomène sans équivalent en mitose : les \cle{chromosomes homologues} (l'un d'origine paternelle, l'autre d'origine maternelle) s'apparient intimement, point par point, pour former des \cle{bivalents} (ou tétrades, car chaque bivalent compte quatre chromatides). Au sein des bivalents, des échanges de segments de chromatides non sœurs peuvent se produire : ce sont les \cle{enjambements} ou \cle{crossing-over}. La membrane nucléaire disparaît en fin de prophase I.

\begin{popfigure}
\begin{center}
\begin{tikzpicture}[scale=1.0, every node/.style={font=\small}]
% Homologue paternel (bleu) : deux chromatides
\draw[line width=3pt, popblue] (0,1.6) -- (0,0.4);
\draw[line width=3pt, popblue] (0.35,1.6) -- (0.35,0.4);
% Homologue maternel (orange)
\draw[line width=3pt, poporange] (1.2,1.6) -- (1.2,0.4);
\draw[line width=3pt, poporange] (1.55,1.6) -- (1.55,0.4);
% centromères
\fill[popdark] (0,1.0) circle (2pt); \fill[popdark] (0.35,1.0) circle (2pt);
\fill[popdark] (1.2,1.0) circle (2pt); \fill[popdark] (1.55,1.0) circle (2pt);
\node[below] at (0.17,0.3) {\scriptsize homologue paternel};
\node[below] at (1.37,0.3) {\scriptsize homologue maternel};
\node[above] at (0.77,1.7) {\textbf{Avant l'enjambement}};
% flèche
\draw[-{Stealth}, very thick, popdark] (2.4,1.0) -- (3.4,1.0);
\node[above, font=\scriptsize] at (2.9,1.0) {crossing-over};
% Après : échange des extrémités
\draw[line width=3pt, popblue] (4.2,1.6) -- (4.2,1.15);
\draw[line width=3pt, poporange] (4.2,1.0) -- (4.2,0.4);
\draw[line width=3pt, popblue] (4.55,1.6) -- (4.55,0.4);
\draw[line width=3pt, poporange] (5.4,1.6) -- (5.4,0.4);
\draw[line width=3pt, poporange] (5.75,1.6) -- (5.75,1.15);
\draw[line width=3pt, popblue] (5.75,1.0) -- (5.75,0.4);
\fill[popdark] (4.2,1.0) circle (2pt); \fill[popdark] (4.55,1.0) circle (2pt);
\fill[popdark] (5.4,1.0) circle (2pt); \fill[popdark] (5.75,1.0) circle (2pt);
\node[above] at (5.0,1.7) {\textbf{Après l'enjambement}};
\node[below, align=center] at (5.0,0.3) {\scriptsize chromatides recombinées\\ \scriptsize (segments échangés)};
% chiasma
\draw[popink, thick, dashed] (4.2,1.08) -- (5.75,1.08);
\node[popink, right, font=\scriptsize] at (6.0,1.08) {chiasma};
\end{tikzpicture}
\end{center}
\legende{Schéma d'un crossing-over en prophase I. Deux chromosomes homologues (bleu : paternel ; orange : maternel), chacun à deux chromatides, échangent des segments de chromatides non sœurs au niveau d'un chiasma. Il en résulte des chromatides \emph{recombinées}, porteuses d'allèles des deux parents.}
\end{popfigure}

\textbf{Métaphase I.} Les bivalents se placent \emph{par paires} sur la plaque équatoriale du fuseau. L'orientation de chaque paire est \textbf{aléatoire} et indépendante des autres paires : c'est la clé du brassage interchromosomique (voir plus loin).

\textbf{Anaphase I.} Les fibres du fuseau se raccourcissent et séparent les \emph{chromosomes homologues} (et non les chromatides !). Chaque chromosome migrant vers un pôle est encore formé de \textbf{deux chromatides}. C'est ici, et nulle part ailleurs, que s'opère la \cle{réduction chromosomique} : chaque futur lot ne contient plus que $n$ chromosomes.

\textbf{Télophase I.} Les chromosomes arrivent aux pôles ; une division cytoplasmique (cytodiérèse) donne \textbf{deux cellules filles haploïdes} ($n$ chromosomes à deux chromatides chacun). L'enveloppe nucléaire peut se reformer ou non selon les espèces.

\courssub{L'interméiose : une interphase sans réplication}

\begin{aretenir}[Point capital]
Entre les deux divisions, l'\cle{interméiose} est une interphase \textbf{dépourvue de phase S} : il n'y a \textbf{aucune réplication de l'ADN}. Les chromosomes entrent donc directement dans la seconde division avec leurs deux chromatides. C'est l'absence de cette réplication qui garantit que les quatre cellules finales seront haploïdes.
\end{aretenir}

\courssub{La deuxième division : la division équationnelle}

Cette seconde division ressemble à s'y méprendre à une \emph{mitose}, mais elle s'opère sur des cellules déjà haploïdes.

\begin{itemize}
\item \textbf{Prophase II :} les chromosomes (toujours à deux chromatides) se recondensent ; le fuseau se forme. Il n'y a \emph{pas} d'appariement, car il n'y a plus de paires d'homologues.
\item \textbf{Métaphase II :} les $n$ chromosomes s'alignent \emph{individuellement} sur la plaque équatoriale.
\item \textbf{Anaphase II :} les centromères se divisent et les \emph{chromatides sœurs} se séparent, migrant vers les pôles opposés. Chaque chromatide devient un chromosome à une chromatide.
\item \textbf{Télophase II :} les lots de $n$ chromosomes simples se constituent ; la cytodiérèse aboutit à \textbf{quatre cellules haploïdes} ($n$ chromosomes à une chromatide).
\end{itemize}

\begin{popfigure}
\begin{center}
\begin{tikzpicture}[scale=0.92, every node/.style={font=\scriptsize}]
% Macro pour une cellule : ellipse + chromosomes
% Etape 1 : cellule mère 2n=4 (avant réplication simplifiée : 2 paires)
\draw[popline, thick] (0,0) ellipse (1.05 and 0.75);
\draw[line width=2.2pt, popblue] (-0.45,0.35) -- (-0.45,-0.35);
\draw[line width=2.2pt, poporange] (-0.15,0.35) -- (-0.15,-0.35);
\draw[line width=2.2pt, poppurple] (0.25,0.25) -- (0.25,-0.25);
\draw[line width=2.2pt, popgreen] (0.55,0.25) -- (0.55,-0.25);
\node[below, align=center] at (0,-0.85) {Cellule mère\\ $2n=4$};
% Etape 2 : prophase I : bivalents
\draw[popline, thick] (2.6,0) ellipse (1.05 and 0.75);
\draw[line width=2.2pt, popblue] (2.0,0.4) -- (2.0,-0.4);
\draw[line width=2.2pt, poporange] (2.22,0.4) -- (2.22,-0.4);
\draw[line width=2.2pt, poppurple] (3.0,0.3) -- (3.0,-0.3);
\draw[line width=2.2pt, popgreen] (3.22,0.3) -- (3.22,-0.3);
\draw[popink, thick] (2.11,0.1) circle (0.06);
\node[below, align=center] at (2.6,-0.85) {Prophase I\\ bivalents, crossing-over};
% Etape 3 : métaphase I
\draw[popline, thick] (5.2,0) ellipse (1.05 and 0.75);
\draw[popmuted, dashed] (4.25,0) -- (6.15,0);
\draw[line width=2.2pt, popblue] (4.85,0.35) -- (4.85,-0.05);
\draw[line width=2.2pt, poporange] (4.85,0.05) -- (4.85,-0.35);
\draw[line width=2.2pt, poppurple] (5.55,0.3) -- (5.55,-0.04);
\draw[line width=2.2pt, popgreen] (5.55,0.04) -- (5.55,-0.3);
\node[below, align=center] at (5.2,-0.85) {Métaphase I\\ bivalents à l'équateur};
% Etape 4 : anaphase I
\draw[popline, thick] (7.8,0) ellipse (1.05 and 0.75);
\draw[line width=2.2pt, popblue] (7.35,0.55) -- (7.35,0.2);
\draw[line width=2.2pt, poppurple] (7.75,0.5) -- (7.75,0.2);
\draw[line width=2.2pt, poporange] (7.35,-0.55) -- (7.35,-0.2);
\draw[line width=2.2pt, popgreen] (7.75,-0.5) -- (7.75,-0.2);
\draw[-{Stealth}, popmuted] (7.0,0.1) -- (7.0,0.45);
\draw[-{Stealth}, popmuted] (7.0,-0.1) -- (7.0,-0.45);
\node[below, align=center] at (7.8,-0.85) {Anaphase I\\ séparation des homologues};
% Etape 5 : télophase I (deux cellules)
\draw[popline, thick] (10.0,0.28) ellipse (0.62 and 0.42);
\draw[popline, thick] (10.0,-0.55) ellipse (0.62 and 0.42);
\draw[line width=2pt, popblue] (9.85,0.45) -- (9.85,0.12);
\draw[line width=2pt, poppurple] (10.15,0.42) -- (10.15,0.14);
\draw[line width=2pt, poporange] (9.85,-0.4) -- (9.85,-0.72);
\draw[line width=2pt, popgreen] (10.15,-0.43) -- (10.15,-0.7);
\node[below, align=center] at (10.0,-1.1) {Télophase I\\ 2 cellules $n$};
% Ligne 2
% Etape 6 : prophase II / métaphase II
\draw[popline, thick] (1.3,-2.6) ellipse (0.62 and 0.42);
\draw[popline, thick] (2.6,-2.6) ellipse (0.62 and 0.42);
\draw[popmuted, dashed] (0.75,-2.6) -- (1.85,-2.6);
\draw[popmuted, dashed] (2.05,-2.6) -- (3.15,-2.6);
\draw[line width=2pt, popblue] (1.2,-2.45) -- (1.2,-2.75);
\draw[line width=2pt, poppurple] (1.45,-2.47) -- (1.45,-2.73);
\draw[line width=2pt, poporange] (2.5,-2.45) -- (2.5,-2.75);
\draw[line width=2pt, popgreen] (2.75,-2.47) -- (2.75,-2.73);
\node[below, align=center] at (1.95,-3.15) {Prophase II -- Métaphase II\\ chromosomes à l'équateur};
% Etape 7 : anaphase II
\draw[popline, thick] (4.6,-2.6) ellipse (0.62 and 0.42);
\draw[popline, thick] (5.9,-2.6) ellipse (0.62 and 0.42);
\draw[line width=2pt, popblue] (4.45,-2.35) -- (4.45,-2.5);
\draw[line width=2pt, popblue] (4.45,-2.7) -- (4.45,-2.85);
\draw[line width=2pt, poppurple] (4.72,-2.37) -- (4.72,-2.5);
\draw[line width=2pt, poppurple] (4.72,-2.7) -- (4.72,-2.83);
\draw[line width=2pt, poporange] (5.75,-2.35) -- (5.75,-2.5);
\draw[line width=2pt, poporange] (5.75,-2.7) -- (5.75,-2.85);
\draw[line width=2pt, popgreen] (6.02,-2.37) -- (6.02,-2.5);
\draw[line width=2pt, popgreen] (6.02,-2.7) -- (6.02,-2.83);
\node[below, align=center] at (5.25,-3.15) {Anaphase II\\ séparation des chromatides};
% Etape 8 : télophase II : 4 cellules
\draw[popline, thick] (8.1,-2.35) ellipse (0.42 and 0.3);
\draw[popline, thick] (9.0,-2.35) ellipse (0.42 and 0.3);
\draw[popline, thick] (8.1,-3.0) ellipse (0.42 and 0.3);
\draw[popline, thick] (9.0,-3.0) ellipse (0.42 and 0.3);
\draw[line width=1.6pt, popblue] (8.02,-2.28) -- (8.02,-2.42);
\draw[line width=1.6pt, poppurple] (8.2,-2.29) -- (8.2,-2.41);
\draw[line width=1.6pt, popblue] (8.92,-2.28) -- (8.92,-2.42);
\draw[line width=1.6pt, poppurple] (9.1,-2.29) -- (9.1,-2.41);
\draw[line width=1.6pt, poporange] (8.02,-2.93) -- (8.02,-3.07);
\draw[line width=1.6pt, popgreen] (8.2,-2.94) -- (8.2,-3.06);
\draw[line width=1.6pt, poporange] (8.92,-2.93) -- (8.92,-3.07);
\draw[line width=1.6pt, popgreen] (9.1,-2.94) -- (9.1,-3.06);
\node[below, align=center] at (8.55,-3.4) {Télophase II\\ 4 cellules $n$ différentes};
% flèches de progression
\draw[-{Stealth}, very thick, popgold] (1.15,0) -- (1.45,0);
\draw[-{Stealth}, very thick, popgold] (3.75,0) -- (4.05,0);
\draw[-{Stealth}, very thick, popgold] (6.35,0) -- (6.65,0);
\draw[-{Stealth}, very thick, popgold] (8.95,0) -- (9.25,0);
\draw[-{Stealth}, very thick, popgold] (10.0,-1.15) -- (10.0,-1.6) -- (3.4,-1.6) -- (3.4,-2.15);
\draw[-{Stealth}, very thick, popgold] (3.35,-2.6) -- (3.85,-2.6);
\draw[-{Stealth}, very thick, popgold] (6.65,-2.6) -- (7.55,-2.6);
\end{tikzpicture}
\end{center}
\legende{Les huit étapes de la méiose pour une cellule $2n = 4$ (deux paires d'homologues : grande paire bleu/orange, petite paire violet/vert). Première ligne : division réductionnelle (séparation des homologues en anaphase I). Deuxième ligne : division équationnelle (séparation des chromatides en anaphase II). Bilan : quatre cellules haploïdes $n = 2$, génétiquement distinctes.}
\end{popfigure}

\begin{methode}[Reconnaître une étape de méiose sur un document (électronographie, coupe de testicule)]
Face à une figure de division, pose-toi trois questions, dans l'ordre :
\begin{enumerate}
\item \textbf{Combien de lots chromosomiques} voit-on, et où ? Des chromosomes qui migrent vers deux pôles indiquent une anaphase ; des chromosomes alignés au centre indiquent une métaphase.
\item \textbf{Les chromosomes sont-ils à une ou deux chromatides ?} Des chromosomes à deux chromatides qui \emph{se séparent} = anaphase I ; des chromatides sœurs qui se séparent = anaphase II (ou mitose).
\item \textbf{Y a-t-il des bivalents ?} La présence de chromosomes appariés en paires est la signature \emph{exclusive} de la prophase I ou de la métaphase I. En mitose et en méiose II, les chromosomes sont toujours isolés.
\end{enumerate}
Critère décisif supplémentaire : si la cellule ne contient que $n$ chromosomes (pas de paires d'homologues), tu es forcément dans la \emph{seconde} division.
\end{methode}

\begin{attention}[L'erreur classique du Baccalauréat]
Beaucoup d'élèves écrivent que « la réduction chromosomique a lieu en anaphase II ». C'est \textbf{faux} : la réduction (passage de $2n$ à $n$) a lieu en \cle{anaphase I}, lors de la séparation des \emph{chromosomes homologues}. L'anaphase II ne fait que séparer des chromatides sœurs au sein de cellules \emph{déjà haploïdes} : elle ne change pas le niveau de ploïdie. Retiens : anaphase I = homologues ; anaphase II = chromatides.
\end{attention}

\courssub{Le brassage génétique : source de la diversité des gamètes}

La méiose ne se contente pas de diviser le stock chromosomique par deux : elle \emph{brasse} les allèles, de deux façons complémentaires.

\textbf{Le brassage interchromosomique.} En métaphase I, l'orientation de chaque bivalent sur la plaque équatoriale est \emph{indépendante} de celle des autres bivalents. Pour une paire d'homologues, il y a deux orientations possibles ; pour $n$ paires, il y a donc $2^{n}$ répartitions possibles des homologues paternels et maternels dans les gamètes.

\begin{popfigure}
\begin{center}
\begin{tikzpicture}[scale=0.95, every node/.style={font=\scriptsize}]
% Cas 1
\draw[popline, thick] (0,0) ellipse (0.8 and 0.6);
\draw[popmuted, dashed] (-0.7,0) -- (0.7,0);
\draw[line width=2.2pt, popblue] (-0.25,0.4) -- (-0.25,0.05);
\draw[line width=2.2pt, poporange] (-0.25,-0.4) -- (-0.25,-0.05);
\draw[line width=2.2pt, poppurple] (0.25,0.35) -- (0.25,0.05);
\draw[line width=2.2pt, popgreen] (0.25,-0.35) -- (0.25,-0.05);
\node[above] at (0,0.7) {\textbf{Orientation 1}};
% Cas 2
\draw[popline, thick] (2.6,0) ellipse (0.8 and 0.6);
\draw[popmuted, dashed] (1.9,0) -- (3.3,0);
\draw[line width=2.2pt, popblue] (2.35,0.4) -- (2.35,0.05);
\draw[line width=2.2pt, poporange] (2.35,-0.4) -- (2.35,-0.05);
\draw[line width=2.2pt, popgreen] (2.85,0.35) -- (2.85,0.05);
\draw[line width=2.2pt, poppurple] (2.85,-0.35) -- (2.85,-0.05);
\node[above] at (2.6,0.7) {\textbf{Orientation 2}};
% flèches vers gamètes
\draw[-{Stealth}, very thick, popgold] (0,-0.7) -- (0,-1.3);
\draw[-{Stealth}, very thick, popgold] (2.6,-0.7) -- (2.6,-1.3);
% Gamètes orientation 1
\draw[popline, thick] (-0.55,-1.9) ellipse (0.4 and 0.3);
\draw[line width=1.8pt, popblue] (-0.62,-1.82) -- (-0.62,-1.98);
\draw[line width=1.8pt, poppurple] (-0.45,-1.83) -- (-0.45,-1.97);
\draw[popline, thick] (0.55,-1.9) ellipse (0.4 and 0.3);
\draw[line width=1.8pt, poporange] (0.48,-1.82) -- (0.48,-1.98);
\draw[line width=1.8pt, popgreen] (0.65,-1.83) -- (0.65,-1.97);
% Gamètes orientation 2
\draw[popline, thick] (2.05,-1.9) ellipse (0.4 and 0.3);
\draw[line width=1.8pt, popblue] (1.98,-1.82) -- (1.98,-1.98);
\draw[line width=1.8pt, popgreen] (2.15,-1.83) -- (2.15,-1.97);
\draw[popline, thick] (3.15,-1.9) ellipse (0.4 and 0.3);
\draw[line width=1.8pt, poporange] (3.08,-1.82) -- (3.08,-1.98);
\draw[line width=1.8pt, poppurple] (3.25,-1.83) -- (3.25,-1.97);
\node[below, align=center] at (1.3,-2.35) {4 types de gamètes : bleu+violet \;|\; orange+vert \;|\; bleu+vert \;|\; orange+violet};
\end{tikzpicture}
\end{center}
\legende{Brassage interchromosomique pour $2n = 4$. Selon l'orientation aléatoire des deux bivalents en métaphase I (deux possibilités par paire, donc $2^{2} = 4$ combinaisons), une même cellule mère peut produire quatre types de gamètes différents.}
\end{popfigure}

\textbf{Le brassage intrachromosomique.} Les crossing-over de la prophase I mélangent les allèles \emph{au sein d'un même chromosome} : les chromatides recombinées portent des associations d'allèles nouvelles, qui n'existaient ni chez le père ni chez la mère. Le nombre de combinaisons possibles devient alors, en pratique, \emph{illimité}.

\begin{exemplebox}[L'immense diversité des gamètes humains]
Chez l'Homme, $2n = 46$, donc $n = 23$ paires d'homologues. Le \emph{seul} brassage interchromosomique permet à un individu de produire :
\[ 2^{23} = \num{8388608} \approx 8,4 \text{ millions de types de gamètes différents.} \]
Un couple peut donc, sans même invoquer les crossing-over, engendrer $2^{23} \times 2^{23} = 2^{46} \approx 7 \times 10^{13}$ combinaisons génétiques différentes, soit environ \num{70000} milliards ! Avec le brassage intrachromosomique, chaque enfant (hors vrais jumeaux) est un être génétiquement unique.
\end{exemplebox}

\begin{propriete}[Double rôle de la méiose (admis)]
La méiose assure la \textbf{pérennité de l'espèce} en réduisant de moitié le stock chromosomique ($2n \to n$), ce qui permet à la fécondation de restaurer la diploïdie sans augmenter le nombre de chromosomes d'une génération à l'autre. Simultanément, elle crée la \textbf{diversité génétique} des gamètes par les brassages \cle{interchromosomique} (répartition aléatoire des homologues) et \cle{intrachromosomique} (crossing-over).
\end{propriete}

\courssub{Bilan et comparaison avec la mitose}

\begin{aretenir}[Bilan chiffré de la méiose]
Une cellule mère $2n$ (chromosomes à deux chromatides après réplication) donne, au terme de deux divisions, \textbf{quatre cellules haploïdes} $n$ (chromosomes à une chromatide), toutes \textbf{génétiquement différentes} entre elles et de la cellule mère.
\end{aretenir}

\begin{center}
\begin{tabularx}{\linewidth}{|l|X|X|}
\hline
\rowcolor{popblueL} \textbf{Critère} & \textbf{Mitose} & \textbf{Méiose} \\
\hline
Nombre de divisions & 1 & 2 (réductionnelle puis équationnelle) \\
\hline
Réplications de l'ADN & 1 & 1 (aucune entre les deux divisions) \\
\hline
Appariement des homologues (bivalents) & Non & Oui, en prophase I \\
\hline
Crossing-over & Non & Possible, en prophase I \\
\hline
Nombre de cellules filles & 2 & 4 \\
\hline
Ploïdie des cellules filles & $2n$ (identique à la mère) & $n$ (moitié) \\
\hline
État génétique des cellules filles & Identiques entre elles et à la mère & Toutes différentes (brassages) \\
\hline
Rôle biologique & Croissance, renouvellement cellulaire, reproduction asexuée & Production des gamètes (ou des spores) \\
\hline
\end{tabularx}
\end{center}

\begin{exoresolu}[Analyse de figures de division]
Sur une coupe de testicule de Rat ($2n = 42$), on observe deux cellules en division. Cellule A : 21 structures doubles alignées par paires à l'équateur. Cellule B : des chromosomes à deux chromatides se séparent en deux lots de 21 chromosomes chacun migrant vers les pôles. Identifier le stade de chaque cellule et justifier.
\tcblower
\textbf{Cellule A.} On compte 21 \emph{paires} de structures doubles : ce sont des \textbf{bivalents} (21 paires d'homologues, chaque chromosome à deux chromatides), alignés sur la plaque équatoriale. La présence de bivalents est la signature de la première division : il s'agit d'une \cle{métaphase I}.

\textbf{Cellule B.} Des chromosomes \emph{à deux chromatides} se séparent, et chaque lot migrant compte 21 chromosomes ($n$). La séparation concerne les \emph{homologues} (les chromatides restent soudées) et le stock passe de $2n = 42$ à $n = 21$ par pôle : il s'agit d'une \cle{anaphase I}, étape de la réduction chromosomique. Si c'était une anaphase II, on verrait des chromatides sœurs (chromosomes simples) se séparer dans des cellules ne contenant que 21 chromosomes au départ.
\end{exoresolu}

\begin{exercice}[Pour t'entraîner]
Chez le Maïs ($2n = 20$) : 1) Combien de bivalents observe-t-on en prophase I ? 2) Combien de chromosomes et de chromatides compte chaque cellule à l'issue de la télophase I ? 3) Combien de types de gamètes le seul brassage interchromosomique permet-il ? 4) Expliquer pourquoi il serait faux de dire que la méiose comporte « deux réplications de l'ADN pour deux divisions ».
\end{exercice}

\begin{corrige}[Exercice : le Maïs]
1) $n = 10$, donc \textbf{10 bivalents}. 2) À l'issue de la télophase I : \textbf{10 chromosomes à 2 chromatides} par cellule (soit 20 chromatides). 3) $2^{10} = \num{1024}$ types de gamètes. 4) Il n'y a qu'\textbf{une seule réplication} (avant la méiose) ; l'interméiose n'en comporte aucune. S'il y avait deux réplications, les cellules finales resteraient diploïdes et la méiose perdrait son rôle réductionnel.
\end{corrige}

\begin{savaistu}[Et chez les Spermaphytes ?]
Chez les plantes à fleurs, la méiose ne produit pas directement les gamètes mais des \textbf{spores} haploïdes : les microspores (dans les anthères) et la mégaspore (dans l'ovaire). Ces spores subissent ensuite quelques mitoses pour former les gamétophytes, au sein desquels naîtront les gamètes. La logique reste identique : la méiose y réduit toujours le stock chromosomique de moitié, préparant la double fécondation que nous étudierons dans la section 5.
\end{savaistu}

Ainsi, la méiose apparaît comme une mécanique chromosomique d'une précision remarquable : une réplication, deux divisions, quatre cellules haploïdes toutes différentes. Reste à comprendre \emph{comment} ces cellules se transforment en gamètes fonctionnels — spermatozoïdes, ovules, grains de pollen — et \emph{comment} la quantité d'ADN évolue à chaque étape : ce sera l'objet des sections 3 et 4.

\coursec{Évolution de la quantité d'ADN au cours de la méiose et de la gamétogenèse}

Dans la section précédente, nous avons suivi pas à pas le comportement des \cle{chromosomes} au cours des deux divisions de la méiose : réplication, appariement des chromosomes homologues, \cle{crossing-over}, puis deux divisions successives. Mais une question reste en suspens : comment les biologistes ont-ils \emph{mesuré} ces événements ? Les chromosomes ne se voient qu'au microscope, à certains stades seulement. En revanche, on peut \textbf{doser la quantité d'ADN} contenue dans le noyau de chaque cellule, à n'importe quel moment. Cette mesure, réalisée sur des milliers de cellules d'un tube séminifère, fournit une courbe remarquable qui constitue la \emph{signature quantitative} de la méiose. C'est cette courbe que nous allons maintenant construire, exploiter et relier au comportement chromosomique.

\courssub{Mise en évidence expérimentale : le dosage de l'ADN dans les cellules d'un tube séminifère}

\textbf{Situation-problème.} Un tube séminifère de testicule de Mammifère contient, alignées de la périphérie vers la lumière, toutes les cellules de la lignée germinale mâle : spermatogonies, spermatocytes, spermatides et spermatozoïdes. Peut-on, en mesurant la quantité d'ADN de chaque cellule, retrouver les grandes étapes de la méiose ?

\begin{experience}[Dosage de l'ADN par cytophotométrie de Feulgen]
\textbf{Matériel :} coupe histologique de testicule de Rat (ou de Mammifère de laboratoire), réactif de Schiff (coloration de Feulgen), cytophotomètre (spectrophotomètre adapté au microscope).

\textbf{Protocole :}
\begin{enumerate}
\item On réalise une coupe fine de testicule et on applique la \cle{coloration de Feulgen}, qui fixe un colorant \emph{spécifiquement} sur l'ADN, en quantité proportionnelle à la quantité d'ADN présente.
\item Pour chaque noyau, on mesure l'\textbf{absorption de la lumière} à travers le noyau coloré : plus le noyau absorbe, plus il contient d'ADN.
\item On mesure ainsi un grand nombre de cellules et on classe les résultats.
\end{enumerate}

\textbf{Résultats :} les cellules se répartissent en \textbf{quatre populations} de quantités d'ADN bien distinctes :
\begin{itemize}
\item des cellules contenant une quantité de référence, notée \cle{Q} (spermatogonies en G1, spermatocytes II après la première division) ;
\item des cellules contenant \cle{2Q} (spermatogonies en G2, spermatocytes I) ;
\item des cellules contenant \cle{Q/2} (spermatides et spermatozoïdes) ;
\item des valeurs \emph{intermédiaires} entre Q et 2Q, correspondant aux cellules en cours de réplication (phase S).
\end{itemize}

\textbf{Interprétation :} l'existence de paliers discrets (Q, 2Q, puis Q/2) montre que la quantité d'ADN varie par \emph{doublements} (réplication) et par \emph{divisions par deux} (divisions cellulaires). On peut donc reconstituer la chronologie complète de la méiose à partir des seules mesures d'ADN.
\end{experience}

\begin{attention}[Attention à la convention de notation]
Selon les manuels, la quantité de référence Q désigne la quantité d'ADN d'une cellule \textbf{diploïde en phase G1} (2n chromosomes à 1 chromatide). Avec cette convention : spermatide et spermatozoïde = Q/2. Certains ouvrages prennent comme référence la quantité haploïde ; les paliers deviennent alors 2Q, 4Q, 2Q, Q. \textbf{Lis toujours la légende de l'axe des ordonnées} avant d'exploiter une courbe : c'est la première chose que l'examinateur vérifie.
\end{attention}

\begin{savaistu}[La coloration de Feulgen, un siècle de bons et loyaux services]
Mise au point en 1924 par les chimistes allemands Robert Feulgen et Heinrich Rossenbeck, la \cle{réaction de Feulgen} repose sur une hydrolyse acide qui libère des groupes aldéhydes de la désoxyribose de l'ADN ; ceux-ci réagissent alors avec le réactif de Schiff pour donner une coloration \emph{fuchsia} (rose-violet intense). Comme l'ARN est détruit par l'hydrolyse, seul l'ADN est coloré : la réaction est \textbf{spécifique} et \textbf{proportionnelle} (loi de Beer-Lambert), ce qui autorise un véritable \emph{dosage} de l'ADN \emph{in situ}, noyau par noyau. C'est grâce à cette technique que l'on a pu démontrer la constance de la quantité d'ADN dans les cellules somatiques d'une espèce — un argument décisif, dès les années 1940, en faveur du rôle de l'ADN comme support de l'hérédité.
\end{savaistu}

\courssub{La courbe d'évolution de la quantité d'ADN : construction et lecture}

En reportant la quantité d'ADN par cellule en fonction du temps, depuis une spermatogonie jusqu'à la spermatide, on obtient la courbe suivante.

\begin{popfigure}
\begin{center}
\begin{tikzpicture}
\begin{axis}[
  width=0.95\linewidth, height=7.5cm,
  xlabel={Temps (phases successives)},
  ylabel={Quantité d'ADN par cellule},
  xmin=0, xmax=10.4, ymin=0, ymax=2.6,
  xtick=\empty,
  ytick={0.5,1,2},
  yticklabels={Q/2, Q, 2Q},
  axis lines=left,
  every axis x label/.style={at={(ticklabel* cs:1.02)}, anchor=west},
  every axis y label/.style={at={(ticklabel* cs:1.02)}, anchor=south},
]
% G1 : palier Q
\addplot[popcurve, domain=0:1.5, samples=20] {1};
% S : montée Q -> 2Q
\addplot[popcurve, domain=1.5:3.5, samples=60] {1 + (x-1.5)/2};
% G2 : palier 2Q
\addplot[popcurve, domain=3.5:5, samples=20] {2};
% Division I : chute 2Q -> Q
\addplot[popcurve, domain=5:6, samples=30] {2 - (x-5)};
% Intercinèse : palier Q
\addplot[popcurve, domain=6:7.5, samples=20] {1};
% Division II : chute Q -> Q/2
\addplot[popcurve, domain=7.5:8.5, samples=30] {1 - (x-7.5)/2};
% Spermiogenèse : palier Q/2
\addplot[popcurve, domain=8.5:10.2, samples=20] {0.5};
% Annotations
\node[popblue, font=\small] at (axis cs:0.75,1.18) {G1};
\node[popblue, font=\small] at (axis cs:2.5,1.75) {S};
\node[popblue, font=\small] at (axis cs:4.25,2.18) {G2};
\node[poporange, font=\small, align=center] at (axis cs:5.5,1.75) {Division I};
\node[popgreen, font=\small, align=center] at (axis cs:6.75,1.18) {intercinèse\\ (pas de S)};
\node[poporange, font=\small, align=center] at (axis cs:8,1.05) {Division II};
\node[poppurple, font=\small, align=center] at (axis cs:9.35,0.72) {spermiogenèse};
% Pointillés de repère
\draw[dashed, popmuted] (axis cs:1.5,0) -- (axis cs:1.5,1);
\draw[dashed, popmuted] (axis cs:3.5,0) -- (axis cs:3.5,2);
\draw[dashed, popmuted] (axis cs:5,0) -- (axis cs:5,2);
\draw[dashed, popmuted] (axis cs:7.5,0) -- (axis cs:7.5,1);
\draw[dashed, popmuted] (axis cs:8.5,0) -- (axis cs:8.5,0.5);
\end{axis}
\end{tikzpicture}
\end{center}
\legende{Courbe d'évolution de la quantité d'ADN par cellule au cours de la spermatogenèse. \textbf{G1} : spermatogonie, quantité Q (2n chromosomes à 1 chromatide). \textbf{S} : réplication, la quantité double de Q à 2Q. \textbf{G2} : 2Q (2n chromosomes à 2 chromatides). \textbf{Division I} : chute de 2Q à Q (n chromosomes à 2 chromatides). \textbf{Intercinèse} : palier à Q, \emph{sans phase S}. \textbf{Division II} : chute de Q à Q/2 (n chromosomes à 1 chromatide). \textbf{Spermiogenèse} : différenciation de la spermatide en spermatozoïde, sans variation de la quantité d'ADN.}
\end{popfigure}

\begin{definition}[Phase S, phase G1, phase G2]
L'\cle{interphase} qui précède la méiose comporte trois phases :
\begin{itemize}
\item la phase \cle{G1} (\emph{gap} 1) : la cellule croît, chaque chromosome est formé d'\textbf{une seule chromatide}, la quantité d'ADN vaut Q ;
\item la phase \cle{S} (\emph{synthèse}) : chaque molécule d'ADN est \textbf{répliquée} ; chaque chromosome devient un chromosome à \textbf{deux chromatides sœurs} ; la quantité d'ADN passe de Q à 2Q ;
\item la phase \cle{G2} : la cellule, à 2Q, achève sa préparation à la division.
\end{itemize}
\end{definition}

\begin{propriete}[Lecture quantitative de la méiose]
Au cours de la méiose :
\begin{itemize}
\item il n'existe qu'\textbf{une seule phase S}, située \emph{avant} la première division : la quantité d'ADN double une unique fois (Q $\rightarrow$ 2Q) ;
\item chaque \textbf{chute brutale} de la courbe correspond à \textbf{une division cellulaire} : la première division ramène la quantité de 2Q à Q, la seconde de Q à Q/2 ;
\item entre les deux divisions, l'\cle{intercinèse} (ou interphase II) est une phase \textbf{courte, sans réplication d'ADN} : la courbe reste au palier Q.
\end{itemize}
Bilan : \textbf{une réplication suivie de deux divisions} divise par quatre la quantité d'ADN par cellule : $2Q \rightarrow Q \rightarrow Q/2$, soit un passage global de Q à Q/2 par rapport à la cellule de départ.
\end{propriete}

\begin{methode}[Exploiter une courbe d'évolution de la quantité d'ADN]
Face à une courbe d'ADN au baccalauréat, procède toujours dans le même ordre :
\begin{enumerate}
\item \textbf{Repère le doublement} (montée progressive) : c'est la \textbf{phase S}, moment de la réplication de l'ADN. Il n'y en a qu'une.
\item \textbf{Repère chaque chute brutale} : chaque chute = \textbf{une division cellulaire}. Deux chutes = méiose ; une seule chute = mitose.
\item \textbf{Compte les paliers} : Q (départ), 2Q (après S), Q (après division I), Q/2 (après division II).
\item \textbf{Vérifie l'absence de remontée entre les deux chutes} : cela prouve qu'il n'y a pas de phase S entre les deux divisions.
\item \textbf{Conclus} en reliant chaque palier à un stade cellulaire (spermatogonie, spermatocyte I, spermatocyte II, spermatide).
\end{enumerate}
\end{methode}

\courssub{Correspondance entre la quantité d'ADN et le comportement des chromosomes}

Le point délicat — et le plus riche au baccalauréat — est de relier chaque palier de la courbe à l'\textbf{état des chromosomes} : nombre de chromosomes ($n$ ou $2n$) et nombre de chromatides par chromosome (1 ou 2). Rappelons la règle fondamentale.

\begin{aretenir}[Quantité d'ADN $\neq$ nombre de chromosomes]
\begin{itemize}
\item Le \textbf{nombre de chromosomes} se compte par les \textbf{centromères} : un chromosome à 2 chromatides ne compte que pour \emph{un} chromosome (un seul centromère).
\item La \textbf{quantité d'ADN} se compte par les \textbf{molécules d'ADN} : chaque chromatide contient \emph{une} molécule d'ADN double brin.
\end{itemize}
Donc : un chromosome à 1 chromatide = 1 molécule d'ADN ; un chromosome à 2 chromatides = 2 molécules d'ADN. C'est pourquoi une cellule peut être \textbf{haploïde} ($n$ chromosomes) tout en contenant une quantité d'ADN \textbf{Q} (si ses chromosomes ont 2 chromatides).
\end{aretenir}

\begin{center}
\rowcolors{2}{popblueL}{white}
\begin{tabularx}{\linewidth}{|l|c|c|c|X|}
\hline
\rowcolor{popblue!40} \textbf{Stade cellulaire} & \textbf{Chromosomes} & \textbf{Chromatides/chromosome} & \textbf{ADN} & \textbf{Palier de la courbe} \\
\hline
Spermatogonie en G1 & $2n$ & 1 & Q & palier initial \\
\hline
Spermatogonie en S & $2n$ & 1 $\rightarrow$ 2 & Q $\rightarrow$ 2Q & montée \\
\hline
Spermatocyte I (G2, méiose I) & $2n$ & 2 & 2Q & palier haut \\
\hline
Spermatocyte II (après division I) & $n$ & 2 & Q & palier intermédiaire \\
\hline
Spermatide (après division II) & $n$ & 1 & Q/2 & palier bas \\
\hline
Spermatozoïde (après spermiogenèse) & $n$ & 1 & Q/2 & palier bas (inchangé) \\
\hline
\end{tabularx}
\end{center}

Analysons les deux chutes de la courbe à la lumière de ce tableau :

\begin{itemize}
\item \textbf{Première division (2Q $\rightarrow$ Q) :} les chromosomes \emph{homologues} se séparent (anaphase I). Chaque cellule fille reçoit $n$ chromosomes, mais chacun est encore formé de \textbf{2 chromatides}. La quantité d'ADN est divisée par deux \emph{parce que le nombre de chromosomes a été divisé par deux}, alors que chaque chromosome reste double.
\item \textbf{Seconde division (Q $\rightarrow$ Q/2) :} les \emph{chromatides sœurs} se séparent (anaphase II). Le nombre de chromosomes reste $n$, mais chaque chromosome passe de 2 chromatides à 1 chromatide : c'est cette fois la \emph{structure} des chromosomes qui change.
\end{itemize}

\begin{attention}[L'erreur classique : confondre la chute 2Q $\rightarrow$ Q avec le passage 2n $\rightarrow$ n]
Beaucoup d'élèves écrivent : « après la première division, la cellule contient Q/2 car elle est devenue haploïde ». \textbf{C'est faux.} Après la première division, la cellule est bien haploïde ($n$ chromosomes), mais chaque chromosome possède encore \textbf{2 chromatides}, donc \textbf{2 molécules d'ADN} : la quantité d'ADN vaut \textbf{Q}, et non Q/2. La chute de 2Q à Q traduit la séparation des homologues ; la chute de Q à Q/2 traduit la séparation des chromatides. Retiens l'image : \emph{la première division divise le nombre de chromosomes, la seconde « dédouble » les chromosomes}.
\end{attention}

\courssub{Prolongement de la courbe pendant la gamétogenèse : la spermiogenèse}

La méiose s'achève avec la formation des spermatides ($n$, 1 chromatide, Q/2). Mais la gamétogenèse n'est pas terminée : la spermatide, cellule ronde et immobile, doit se \textbf{transformer} en spermatozoïde flagellé. C'est la \cle{spermiogenèse} : condensation du noyau, formation de l'acrosome à partir de l'appareil de Golgi, développement du flagelle à partir du centriole, élimination de l'essentiel du cytoplasme.

\begin{propriete}[La spermiogenèse ne modifie pas la quantité d'ADN]
La spermiogenèse est une \textbf{différenciation cellulaire sans division ni réplication} : la courbe d'ADN reste sur le palier Q/2 pendant toute cette phase. Le spermatozoïde mature contient donc la même quantité d'ADN que la spermatide dont il dérive : Q/2, soit $n$ chromosomes à 1 chromatide.
\end{propriete}

\begin{exemplebox}[Un exemple chiffré : la spermatogenèse humaine]
Chez l'Homme ($2n = 46$ chromosomes), les mesures cytophotométriques donnent :
\begin{itemize}
\item spermatogonie en G1 : \SI{6,4}{pg} d'ADN par noyau (= Q) ;
\item spermatocyte I en G2 : \SI{12,8}{pg} (= 2Q) après réplication ;
\item spermatocyte II : \SI{6,4}{pg} (= Q), avec seulement $n = 23$ chromosomes à 2 chromatides ;
\item spermatide et spermatozoïde : \SI{3,2}{pg} (= Q/2), avec $n = 23$ chromosomes à 1 chromatide.
\end{itemize}
On vérifie que $\dfrac{\num{6,4}}{2} = \num{3,2}$ pg : la méiose a bien divisé par deux la quantité d'ADN, préparant la restauration de la diploïdie (\SI{6,4}{pg}) lors de la fécondation.
\end{exemplebox}

\begin{exoresolu}[Exploitation d'une courbe d'évolution de la quantité d'ADN]
\textbf{Énoncé.} On a mesuré la quantité d'ADN par noyau dans les cellules d'un tube séminifère de Rat. On obtient quatre groupes de cellules contenant respectivement : 5 unités arbitraires (u.a.), 10 u.a., des valeurs comprises entre 5 et 10 u.a., et 2,5 u.a.
\begin{enumerate}
\item Attribuer à chaque groupe le stade de la lignée germinale correspondant.
\item Un élève affirme : « les cellules à 5 u.a. sont toutes haploïdes ». Cette affirmation est-elle correcte ? Justifier.
\item Expliquer pourquoi on ne trouve jamais de cellules dont la quantité d'ADN est comprise entre 2,5 et 5 u.a.
\end{enumerate}
\tcblower
\textbf{Solution détaillée.}
\begin{enumerate}
\item \textbf{Cellules à 5 u.a. (= Q) :} spermatogonies en G1 \emph{et} spermatocytes II (après la première division). \textbf{Cellules à 10 u.a. (= 2Q) :} spermatogonies en G2 et spermatocytes I en méiose I. \textbf{Valeurs entre 5 et 10 u.a. :} cellules en \textbf{phase S}, en cours de réplication de leur ADN. \textbf{Cellules à 2,5 u.a. (= Q/2) :} spermatides et spermatozoïdes.
\item \textbf{L'affirmation est fausse.} Le palier Q regroupe deux types de cellules : les spermatogonies en G1, qui sont \textbf{diploïdes} ($2n$ chromosomes à 1 chromatide), et les spermatocytes II, qui sont \textbf{haploïdes} ($n$ chromosomes à 2 chromatides). La quantité d'ADN seule ne permet pas de connaître le nombre de chromosomes : il faut aussi connaître la structure des chromosomes (1 ou 2 chromatides). C'est toute la différence entre Q (quantité d'ADN) et $n$ (nombre de chromosomes).
\item Une valeur intermédiaire entre Q/2 et Q signifierait une \textbf{réplication d'ADN entre la fin de la méiose et la formation des gamètes}, c'est-à-dire une phase S pendant la spermiogenèse ou l'intercinèse. Or l'intercinèse est dépourvue de phase S et la spermiogenèse est une simple différenciation : l'ADN n'y est jamais répliqué. L'absence de telles valeurs \textbf{prouve expérimentalement} qu'il n'existe qu'une seule phase S, située avant la première division.
\end{enumerate}
\end{exoresolu}

\begin{exercice}[À toi de jouer]
Chez le Chimpanzé, $2n = 48$ chromosomes. Une spermatogonie en G1 contient \SI{7,0}{pg} d'ADN.
\begin{enumerate}
\item Compléter le tableau suivant : spermatocyte I en métaphase I, spermatocyte II en métaphase II, spermatide — pour chacun, donner le nombre de chromosomes, le nombre de chromatides par chromosome et la quantité d'ADN en pg.
\item Tracer l'allure de la courbe d'évolution de la quantité d'ADN de la spermatogonie au spermatozoïde, en annotant les phases.
\item Que deviendrait la quantité d'ADN du zygote si la méiose n'existait pas, au bout de trois générations de fécondations ? Conclure sur le rôle de la méiose.
\end{enumerate}
\end{exercice}

\begin{corrige}[Exercice : À toi de jouer]
\begin{enumerate}
\item Spermatocyte I en métaphase I : $2n = 48$ chromosomes à 2 chromatides, ADN = $2Q = \SI{14,0}{pg}$. Spermatocyte II en métaphase II : $n = 24$ chromosomes à 2 chromatides, ADN = $Q = \SI{7,0}{pg}$. Spermatide : $n = 24$ chromosomes à 1 chromatide, ADN = $Q/2 = \SI{3,5}{pg}$.
\item Courbe : palier à \SI{7,0}{pg} (G1), montée progressive jusqu'à \SI{14,0}{pg} (S), palier (G2), chute à \SI{7,0}{pg} (division I), palier (intercinèse sans S), chute à \SI{3,5}{pg} (division II), palier constant (spermiogenèse).
\item Sans méiose, chaque fécondation doublerait la quantité d'ADN : génération 1 : \SI{14,0}{pg} ; génération 2 : \SI{28,0}{pg} ; génération 3 : \SI{56,0}{pg}. La quantité d'ADN doublerait indéfiniment à chaque génération, ce qui est incompatible avec la \textbf{constance du caryotype} de l'espèce. La méiose, en divisant par deux la quantité d'ADN des gamètes, compense exactement le doublement réalisé à la fécondation : c'est le cœur de la complémentarité méiose--fécondation, que nous détaillerons dans les sections suivantes.
\end{enumerate}
\end{corrige}

\begin{aretenir}[L'essentiel de la section]
\begin{itemize}
\item La quantité d'ADN par cellule évolue par paliers : \textbf{Q} (G1) $\rightarrow$ \textbf{2Q} (après la phase S) $\rightarrow$ \textbf{Q} (après la division I) $\rightarrow$ \textbf{Q/2} (après la division II), puis reste à Q/2 pendant la spermiogenèse.
\item Une seule réplication (phase S) précède deux divisions ; il n'y a \textbf{pas de phase S entre les deux divisions}.
\item \textbf{Q et $n$ sont deux notions distinctes} : $n$ compte les centromères, Q compte les molécules d'ADN. Un spermatocyte II est haploïde ($n$) mais contient Q car ses chromosomes ont 2 chromatides.
\item La division I sépare les homologues (2n $\rightarrow$ n) ; la division II sépare les chromatides (chromosomes doubles $\rightarrow$ simples).
\item Chez l'Homme : \SI{6,4}{pg} (spermatogonie G1) $\rightarrow$ \SI{12,8}{pg} $\rightarrow$ \SI{6,4}{pg} $\rightarrow$ \SI{3,2}{pg} (spermatozoïde).
\end{itemize}
\end{aretenir}

\coursec{La gamétogenèse chez les Mammifères et les Spermaphytes}

Dans la section précédente, nous avons suivi l'évolution de la quantité d'ADN au cours de la méiose : une cellule diploïde à $2n$ chromosomes et \textbf{4Q d'ADN} (après réplication) donne quatre cellules haploïdes à $n$ chromosomes et \textbf{Q d'ADN}. Mais une cellule haploïde n'est pas encore un gamète fonctionnel : un spermatozoïde doit acquérir un flagelle pour nager, un ovule doit accumuler des réserves pour nourrir le futur embryon. Entre la cellule souche diploïde de la gonade et le gamète mature, il se déroule toute une histoire cellulaire : c'est la \cle{gamétogenèse}, que nous étudions maintenant chez les Mammifères puis chez les Spermaphytes.

\courssub{Définition et étapes générales de la gamétogenèse}

\begin{definition}[Gamétogenèse]
La \cle{gamétogenèse} est l'ensemble des phénomènes conduisant à la formation des \cle{gamètes} (cellules reproductrices haploïdes) à partir des cellules souches diploïdes des gonades. Elle comprend quatre phases :
\begin{enumerate}
\item la \cle{multiplication} : les cellules souches (gonies) se divisent par \textbf{mitoses}, ce qui maintient le stock de cellules diploïdes ($2n$, $2Q$) ;
\item la \cle{croissance} : les cellules augmentent de volume et répliquent leur ADN (passage de $2Q$ à \textbf{4Q}) ;
\item la \cle{maturation} : c'est la \textbf{méiose}, qui réduit la ploïdie de $2n$ à $n$ et la quantité d'ADN de $4Q$ à $Q$ ;
\item la \cle{différenciation} : les cellules haploïdes ($n$, $Q$) acquièrent leur structure définitive de gamète.
\end{enumerate}
\end{definition}

\begin{attention}[Ne pas confondre gamétogenèse et méiose]
La méiose n'est qu'\textbf{une étape} (la maturation) de la gamétogenèse. Dire « la gamétogenèse est une méiose » est une erreur classique au Baccalauréat : la gamétogenèse inclut aussi des mitoses (multiplication), une croissance (réplication de l'ADN) et une différenciation.
\end{attention}

\begin{methode}[Interpréter la courbe d'évolution de la quantité d'ADN au cours de la gamétogenèse]
Pour analyser une courbe représentant la quantité d'ADN en fonction du temps au cours de la gamétogenèse :
\begin{enumerate}
\item \textbf{Phase de multiplication (mitoses) :} plateau à \textbf{2Q} (cellules $2n$ en G1) ; pas de variation de la ploïdie.
\item \textbf{Phase de croissance (phase S) :} montée brutale de \textbf{2Q à 4Q} (réplication de l'ADN) ; la cellule est $2n$ mais avec des chromosomes à deux chromatides.
\item \textbf{Phase de maturation (méiose I) :} chute de \textbf{4Q à 2Q} (séparation des chromosomes homologues) ; passage de $2n$ à $n$ chromosomes (chacun à deux chromatides).
\item \textbf{Phase de maturation (méiose II) :} chute de \textbf{2Q à Q} (séparation des chromatides) ; cellules $n$ à une chromatide.
\item \textbf{Phase de différenciation :} plateau à \textbf{Q} ; pas de division, mise en place des structures spécifiques du gamète.
\end{enumerate}
\end{methode}

\courssub{La spermatogenèse chez les Mammifères}

\begin{experience}[Observation histologique : coupe de tube séminifère (rat ou verrat)]
\textbf{Matériel :} Lame mince de testicule colorée (hématoxyline-éosine), microscope optique.\par
\textbf{Protocole :} Observer à grossissement fort (×400) la paroi d'un tube séminifère.\par
\textbf{Observations :} Les cellules sont \textbf{ordonnées de la paroi vers la lumière} du tube, comme sur une chaîne de fabrication.
\begin{itemize}
\item Au contact de la membrane basale : grosses cellules à noyau rond, chromatine fine : \textbf{spermatogonies} ($2n$).
\item Vers la lumière : cellules très grosses, noyau volumineux à chromatine filamenteuse (prophase I) : \textbf{spermatocytes I} ($2n$, $4Q$).
\item Plus loin : cellules plus petites, noyau dense : \textbf{spermatocytes II} ($n$, $2Q$) puis \textbf{spermatides} ($n$, $Q$), rondes puis allongées.
\item Dans la lumière : \textbf{spermatozoïdes} ($n$, $Q$) aux flagelles mobiles.
\end{itemize}
\textbf{Interprétation :} Cette disposition spatiale reflète la chronologie de la spermatogenèse : les divisions mitotiques et le début de la méiose ont lieu près de la paroi (niche stem cell), la différenciation finale (spermiogenèse) se fait vers la lumière où les spermatozoïdes sont libérés (spermiation).
\end{experience}

\begin{methode}[Identifier une cellule de la lignée germinale sur une coupe]
Pour reconnaître une cellule de la lignée spermatique, utiliser trois critères :
\begin{enumerate}
\item sa \textbf{position} : les cellules jeunes (spermatogonies, spermatocytes I) sont contre la paroi ; les cellules évoluées (spermatides, spermatozoïdes) sont proches de la lumière ;
\item sa \textbf{taille} : les spermatocytes I sont les plus grosses cellules de la lignée ; les spermatides sont les plus petites ;
\item son \textbf{noyau} : noyau volumineux à chromatine fine chez les spermatocytes ; noyau petit, dense et parfois allongé chez les spermatides en cours de spermiogenèse.
\end{enumerate}
\end{methode}

\textbf{Déroulement de la spermatogenèse.}\par
\begin{itemize}
\item Les \cle{spermatogonies} ($2n$, $2Q$), situées contre la paroi du tube, se multiplient par mitoses tout au long de la vie de l'homme, à partir de la puberté.
\item Certaines entrent en croissance (phase S) et deviennent des \cle{spermatocytes de premier ordre} ou spermatocytes I ($2n$ chromosomes, \textbf{4Q d'ADN}).
\item Chaque spermatocyte I accomplit la \textbf{première division de méiose} et donne deux \cle{spermatocytes de second ordre} ou spermatocytes II ($n$ chromosomes, \textbf{2Q d'ADN}, chromosomes à deux chromatides).
\item Chaque spermatocyte II accomplit la \textbf{deuxième division de méiose} et donne deux \cle{spermatides} ($n$ chromosomes à une chromatide, \textbf{Q d'ADN}).
\item Les spermatides se transforment ensuite en \cle{spermatozoïdes} ($n$, $Q$) sans aucune division : c'est la \cle{spermiogenèse}.
\end{itemize}

\begin{definition}[Spermiogenèse]
La \cle{spermiogenèse} est la phase de \textbf{différenciation} qui transforme les spermatides haploïdes ($n$, $Q$) en spermatozoïdes mobiles, \textbf{sans division cellulaire ni variation de la quantité d'ADN}. Au cours de la spermiogenèse : le noyau se condense et s'allonge, l'\cle{acrosome} (vésicule d'enzymes issue de l'appareil de Golgi, indispensable à la pénétration de l'ovule) se forme à l'extrémité antérieure, les mitochondries s'assemblent en gaine spiralée dans la pièce intermédiaire, et un \cle{flagelle} se développe à partir du centriole. Le cytoplasme excédentaire est éliminé (phagocyté par les cellules de Sertoli).
\end{definition}

\begin{aretenir}[Bilan de la spermatogenèse]
Un spermatocyte I ($2n$, $4Q$) donne, par méiose, \textbf{quatre spermatides} ($n$, $Q$), qui deviennent \textbf{quatre spermatozoïdes fonctionnels}. La spermatogenèse est \textbf{continue} de la puberté à la mort et produit des millions de spermatozoïdes chaque jour (environ $1 \times 10^8$ par jour chez l'homme adulte).
\end{aretenir}

\courssub{L'oogenèse chez les Mammifères}

\textbf{Observation.}\par
Sur des coupes d'ovaire de lapine ou de truie, on n'observe jamais les alignements ordonnés du testicule, mais des \textbf{follicules} de tailles variées, contenant chacun une seule grosse cellule : l'ovocyte. De plus, les divisions y sont \textbf{inégales} : on observe à côté de l'ovocyte de minuscules cellules, les globules polaires, qui dégénèrent. Pourquoi cette différence avec la spermatogenèse ?

\textbf{Déroulement de l'oogenèse.}\par
\begin{itemize}
\item Pendant la \textbf{vie fœtale}, les \cle{ovogonies} ($2n$, $2Q$) se multiplient par mitoses, puis entrent en croissance et deviennent des \cle{ovocytes de premier ordre} (ovocytes I, $2n$, \textbf{4Q}). Ceux-ci \textbf{commencent la méiose mais se bloquent en prophase I} dès la naissance : la femelle naît avec son stock définitif d'ovocytes I (environ 400 000 à la puberté chez la femme, dont seulement 400 à 500 environ seront ovulés au cours de la vie).
\item À chaque cycle ovarien, à partir de la puberté, un ovocyte I (parfois plusieurs) achève la \textbf{première division de méiose} : il donne un gros \cle{ovocyte de second ordre} (ovocyte II, $n$, \textbf{2Q}) et un petit \cle{premier globule polaire} ($n$, \textbf{2Q}), pratiquement dépourvu de cytoplasme.
\item La \textbf{deuxième division de méiose} commence mais reste bloquée en métaphase II ; elle ne s'achève \textbf{qu'en cas de fécondation}, donnant l'\cle{ovule} ($n$, \textbf{Q}) et un \cle{deuxième globule polaire} ($n$, \textbf{Q}).
\end{itemize}

Les divisions inégales ont une signification adaptative claire : le cytoplasme, riche en réserves (vitellus) et en organites, est conservé dans une seule cellule, le futur gamète, pour assurer les premiers stades du développement embryonnaire. Les globules polaires, simples « déchets » de chromosomes, dégénèrent.

\begin{attention}[Erreur fréquente : le rendement de l'oogenèse]
Contrairement à la spermatogenèse, l'oogenèse ne produit \textbf{pas quatre ovules} ! Un ovocyte I donne \textbf{un seul gamète fonctionnel} (l'ovule) et deux ou trois globules polaires qui dégénèrent. De même, ne dis pas que la femme « fabrique des ovocytes toute sa vie » : le stock est constitué avant la naissance.
\end{attention}

\begin{savaistu}[Un blocage de plus de 40 ans]
Chez la femme, les ovocytes I restent bloqués en prophase I de la vie fœtale jusqu'à l'ovulation, parfois plus de 40 ans ! Or c'est en prophase I que les chromosomes homologues sont appariés et que se prépare leur séparation. Ce long blocage fragilise les mécanismes de disjonction : c'est pourquoi la fréquence des anomalies chromosomiques (comme la trisomie 21) augmente nettement avec l'âge maternel. C'est un enjeu de santé publique bien connu dans les maternités camerounaises, où les sages-femmes informent les femmes sur ce risque.
\end{savaistu}

\begin{popfigure}
\begin{tikzpicture}[
  cell/.style={circle, draw=popblue, fill=popblueL, minimum size=0.85cm, font=\scriptsize, align=center},
  cellp/.style={circle, draw=poppink, fill=poppinkL, minimum size=0.85cm, font=\scriptsize, align=center},
  small/.style={circle, draw=poppink, fill=poppinkL, minimum size=0.45cm, font=\tiny, align=center},
  etape/.style={font=\scriptsize\itshape, text=popmuted},
  fl/.style={-{Stealth[length=2.5mm]}, thick, popdark}
]
% ==== Spermatogenèse (gauche) ====
\node[font=\small\bfseries, text=popblue] at (0,4.6) {Spermatogenèse};
\node[cell, align=center] (sg) at (0,3.8){Spermatogonie\\ $2n$, $2Q$};
\node[cell, minimum size=1.05cm, align=center] (s1) at (0,2.6){Spermatocyte I\\ $2n$, $4Q$};
\node[cell, minimum size=0.8cm, align=center] (s2a) at (-1.1,1.3){Sperm. II\\ $n$, $2Q$};
\node[cell, minimum size=0.8cm, align=center] (s2b) at (1.1,1.3){Sperm. II\\ $n$, $2Q$};
\node[cell, minimum size=0.6cm, align=center] (st1) at (-1.9,0){Spermatide\\ $n$, $Q$};
\node[cell, minimum size=0.6cm, align=center] (st2) at (-0.5,0){Spermatide\\ $n$, $Q$};
\node[cell, minimum size=0.6cm, align=center] (st3) at (0.5,0){Spermatide\\ $n$, $Q$};
\node[cell, minimum size=0.6cm, align=center] (st4) at (1.9,0){Spermatide\\ $n$, $Q$};
\node[cell, minimum size=0.55cm, fill=popgreenL, draw=popgreen, align=center] (sz1) at (-1.9,-1.2){Spermatoz.\\ $n$, $Q$};
\node[cell, minimum size=0.55cm, fill=popgreenL, draw=popgreen, align=center] (sz2) at (-0.5,-1.2){Spermatoz.\\ $n$, $Q$};
\node[cell, minimum size=0.55cm, fill=popgreenL, draw=popgreen, align=center] (sz3) at (0.5,-1.2){Spermatoz.\\ $n$, $Q$};
\node[cell, minimum size=0.55cm, fill=popgreenL, draw=popgreen, align=center] (sz4) at (1.9,-1.2){Spermatoz.\\ $n$, $Q$};
\draw[fl] (sg) -- (s1) node[midway, right, etape]{croissance (S)};
\draw[fl] (s1) -- (s2a) node[midway, above left, etape]{méiose I};
\draw[fl] (s1) -- (s2b);
\draw[fl] (s2a) -- (st1) node[midway, left, etape]{méiose II};
\draw[fl] (s2a) -- (st2);
\draw[fl] (s2b) -- (st3);
\draw[fl] (s2b) -- (st4);
\draw[fl, popgreen] (st1) -- (sz1);
\draw[fl, popgreen] (st2) -- (sz2);
\draw[fl, popgreen] (st3) -- (sz3);
\draw[fl, popgreen] (st4) -- (sz4);
\node[etape, text=popgreen] at (2.9,-0.6) {spermiogenèse};
\node[font=\scriptsize\bfseries, text=popgreen] at (0,-2) {4 gamètes fonctionnels};
% ==== Oogenèse (droite) ====
\begin{scope}[xshift=8.2cm]
\node[font=\small\bfseries, text=poppink] at (0,4.6) {Oogenèse};
\node[cellp, align=center] (og) at (0,3.8){Ovogonie\\ $2n$, $2Q$};
\node[cellp, minimum size=1.05cm, align=center] (o1) at (0,2.6){Ovocyte I\\ $2n$, $4Q$};
\node[cellp, minimum size=0.95cm, align=center] (o2) at (-0.7,1.2){Ovocyte II\\ $n$, $2Q$};
\node[small, align=center] (gp1) at (1.4,1.2){GP1\\ $n$, $2Q$};
\node[cellp, minimum size=0.9cm, fill=popgreenL, draw=popgreen, align=center] (ov) at (-0.7,-0.3){Ovule\\ $n$, $Q$};
\node[small, align=center] (gp2) at (0.9,-0.3){GP2\\ $n$, $Q$};
\draw[fl] (og) -- (o1) node[midway, right, etape]{croissance (vie fœtale)};
\draw[fl] (o1) -- (o2) node[midway, above left, etape]{méiose I};
\draw[fl] (o1) -- (gp1);
\draw[fl] (o2) -- (ov) node[midway, left, etape]{méiose II};
\draw[fl] (o2) -- (gp2);
\node[etape] at (1.6,0.6) {globules polaires};
\node[etape] at (1.6,-1) {dégénèrent};
\node[font=\scriptsize\bfseries, text=popgreen] at (0,-2) {1 seul gamète fonctionnel};
\end{scope}
\end{tikzpicture}
\legende{Comparaison des arbres de la spermatogenèse (à gauche) et de l'oogenèse (à droite) chez les Mammifères. La spermatogenèse produit quatre spermatozoïdes fonctionnels par spermatocyte I ; l'oogenèse ne produit qu'un seul ovule fonctionnel, les globules polaires (GP1, GP2) dégénèrent. La spermiogenèse (flèches vertes) est une différenciation sans division. Quantités d'ADN indiquées : Q = contenu haploïde.}
\end{popfigure}

\begin{exoresolu}[Rendement comparé des deux gamétogenèses]
Énoncé : chez un Mammifère, on suit le devenir de 100 spermatocytes I et de 100 ovocytes I ayant tous achevé la méiose. Combien de gamètes fonctionnels obtient-on dans chaque cas ? Justifier.
\tcblower
\textbf{Solution.} Chaque spermatocyte I donne, par les deux divisions de méiose, quatre spermatides, toutes transformées en spermatozoïdes fonctionnels : $100 \times 4 = 400$ \textbf{spermatozoïdes fonctionnels}. Chaque ovocyte I donne un seul ovocyte II puis un seul ovule fonctionnel, les globules polaires dégénérant : $100 \times 1 = 100$ \textbf{ovules fonctionnels} (et 200 à 300 globules polaires non fonctionnels). Conclusion : à ploïdie égale, le rendement de la spermatogenèse est quatre fois supérieur ; l'oogenèse privilégie la qualité (cytoplasme riche en réserves) au détriment de la quantité.
\end{exoresolu}

\courssub{Comparaison des gamétogenèses mâle et femelle chez les Mammifères}

\begin{center}
\begin{tabularx}{\linewidth}{|l|X|X|}
\hline
\rowcolor{popblueL} \textbf{Critère} & \textbf{Spermatogenèse} & \textbf{Oogenèse} \\
\hline
Lieu & Tubes séminifères des testicules & Follicules ovariens de l'ovaire \\
\hline
Période & Continue, de la puberté à la mort & Cyclique, de la puberté à la ménopause ; débutée dans la vie fœtale \\
\hline
Divisions & Égales & Inégales (conservation du cytoplasme) \\
\hline
Rendement & 4 spermatozoïdes par spermatocyte I & 1 ovule + globules polaires par ovocyte I \\
\hline
Gamète & Petit, mobile (flagelle), sans réserves & Gros, immobile, riche en réserves \\
\hline
\end{tabularx}
\end{center}

\begin{popfigure}
\begin{tikzpicture}[scale=1]
% Frise mâle
\draw[-{Stealth}, very thick, popblue] (0,2.2) -- (12,2.2);
\node[font=\small\bfseries, text=popblue, anchor=west] at (0,2.9) {Spermatogenèse (homme)};
\fill[popblue] (3.2,2.2) circle (2pt) node[above=2pt, font=\scriptsize, text=popdark]{puberté ($\approx$ 13 ans)};
\fill[popblue] (11,2.2) circle (2pt) node[above=2pt, font=\scriptsize, text=popdark]{mort};
\draw[popblue, thick, decorate, decoration={brace, amplitude=4pt}] (3.2,1.95) -- (11,1.95) node[midway, below=6pt, font=\scriptsize, text=popdark]{production continue, jour et nuit};
% Frise femelle
\draw[-{Stealth}, very thick, poppink] (0,-0.4) -- (12,-0.4);
\node[font=\small\bfseries, text=poppink, anchor=west] at (0,0.3) {Oogenèse (femme)};
\fill[poppink] (0.6,-0.4) circle (2pt) node[below=2pt, font=\scriptsize, text=popdark, align=center]{vie fœtale :\\ ovocytes I bloqués\\ en prophase I};
\fill[poppink] (3.2,-0.4) circle (2pt) node[below=2pt, font=\scriptsize, text=popdark]{puberté ($\approx$ 12 ans)};
\fill[poppink] (8.6,-0.4) circle (2pt) node[below=2pt, font=\scriptsize, text=popdark]{ménopause ($\approx$ 50 ans)};
\draw[poppink, thick, decorate, decoration={brace, amplitude=4pt, mirror}] (3.2,-1.7) -- (8.6,-1.7) node[midway, below=6pt, font=\scriptsize, text=popdark, align=center]{cycles d'environ 28 jours :\\ 1 ovocyte II libéré par cycle};
\end{tikzpicture}
\legende{Frise comparative du calendrier des gamétogenèses chez l'Homme. Chez l'homme, la production de spermatozoïdes est continue de la puberté à la mort. Chez la femme, les ovocytes I sont constitués et bloqués en prophase I dès la vie fœtale ; un seul ovocyte achève la méiose I à chaque cycle, de la puberté à la ménopause.}
\end{popfigure}

\courssub{La gamétogenèse chez les Spermaphytes}

Chez les plantes à fleurs, les gonades sont les organes de la fleur : les \textbf{étamines} (organe mâle) et le \textbf{pistil} (organe femelle). L'observation d'une fleur de hibiscus ou de flamboyant, courante dans nos jardins, permet de les identifier facilement.

\textbf{Gamétogenèse mâle : formation du grain de pollen.}\par
Dans les loges polliniques de l'\cle{anthère}, les \cle{cellules-mères du pollen} (ou microsporocytes, $2n$, $4Q$) subissent la méiose : chacune donne \textbf{quatre microspores} haploïdes ($n$, $Q$), groupées en tétrade. Chaque microspore effectue ensuite \textbf{une mitose} et devient un \cle{grain de pollen} contenant \textbf{deux noyaux haploïdes} ($n$, $Q$) : le \cle{noyau végétatif} (qui dirigera la formation du tube pollinique) et le \cle{noyau générateur} (qui donnera plus tard, par mitose dans le tube pollinique, les deux noyaux spermatiques).

\textbf{Gamétogenèse femelle : formation du sac embryonnaire.}\par
Dans l'\cle{ovule} logé dans l'ovaire du pistil, la \cle{cellule-mère du sac embryonnaire} (ou mégasporocyte, $2n$, $4Q$) subit la méiose et donne \textbf{quatre mégaspores} haploïdes ($n$, $Q$), dont \textbf{trois dégénèrent}. La mégaspore survivante effectue \textbf{trois mitoses successives} sans cloisonnement, donnant le \cle{sac embryonnaire}, qui contient \textbf{huit noyaux haploïdes} ($n$, $Q$) répartis en sept cellules :
\begin{itemize}
\item l'\cle{oosphere} ($n$, $Q$) : le gamète femelle proprement dit ;
\item deux \cle{synergides} ($n$, $Q$), qui accompagnent l'oosphère ;
\item trois \cle{antipodes} ($n$, $Q$), à l'extrémité opposée ;
\item deux \cle{noyaux polaires} ($n$, $Q$), au centre (souvent fusionnés en un noyau diploïde secondaire $2n$, $2Q$ avant la fécondation).
\end{itemize}

\begin{aretenir}[L'essentiel à retenir]
Le détail des mitoses formatrices du sac embryonnaire est donné à titre informatif. Ce qu'il faut impérativement retenir : \textbf{toutes les cellules du grain de pollen et du sac embryonnaire sont haploïdes} ($n$), car elles dérivent de la méiose. Le grain de pollen possède deux noyaux ($n$) ; le sac embryonnaire possède huit noyaux ($n$), dont l'oosphère et les deux noyaux polaires, qui joueront un rôle clé dans la double fécondation (section suivante).
\end{aretenir}

\begin{popfigure}
\begin{tikzpicture}[scale=1]
% ==== Grain de pollen ====
\begin{scope}
\draw[very thick, poporange, fill=popgoldL] (0,0) circle (1.5);
\draw[poporange, thick] (0,0) circle (1.62);
\node[circle, draw=popblue, fill=popblueL, minimum size=0.95cm, font=\scriptsize, align=center] (nv) at (-0.45,0.35) {noyau\\ végétatif\\ $n$};
\node[ellipse, draw=poppurple, fill=poppurpleL, minimum width=0.75cm, minimum height=1.05cm, font=\scriptsize, align=center] (ng) at (0.7,-0.55) {noyau\\ générateur\\ $n$};
\node[font=\small\bfseries, text=poporange] at (0,-2.2) {Grain de pollen ($n$)};
\end{scope}
% ==== Sac embryonnaire ====
\begin{scope}[xshift=8.6cm]
\draw[very thick, popgreen, fill=popgreenL] (0,0) ellipse (1.7 and 2.6);
% oosphère + synergides (haut)
\node[circle, draw=popblue, fill=popblueL, minimum size=0.95cm, font=\scriptsize, align=center] (oo) at (0,1.75) {oosphere\\ $n$};
\node[ellipse, draw=popmuted, fill=popcream, minimum width=0.6cm, minimum height=0.85cm, font=\tiny, align=center] (s1) at (-0.85,1.6) {syner-\\ gide};
\node[ellipse, draw=popmuted, fill=popcream, minimum width=0.6cm, minimum height=0.85cm, font=\tiny, align=center] (s2) at (0.85,1.6) {syner-\\ gide};
% noyaux polaires (centre)
\node[circle, draw=poppurple, fill=poppurpleL, minimum size=0.7cm, font=\scriptsize] (np1) at (-0.35,0) {$n$};
\node[circle, draw=poppurple, fill=poppurpleL, minimum size=0.7cm, font=\scriptsize] (np2) at (0.35,0) {$n$};
\node[font=\tiny, text=popdark] at (0,-0.75) {2 noyaux polaires};
% antipodes (bas)
\node[circle, draw=popmuted, fill=popcream, minimum size=0.55cm, font=\tiny] (a1) at (-0.6,-1.9) {ant.};
\node[circle, draw=popmuted, fill=popcream, minimum size=0.55cm, font=\tiny] (a2) at (0,-2.05) {ant.};
\node[circle, draw=popmuted, fill=popcream, minimum size=0.55cm, font=\tiny] (a3) at (0.6,-1.9) {ant.};
\node[font=\small\bfseries, text=popgreen] at (0,-3.1) {Sac embryonnaire (8 noyaux, tous $n$)};
\end{scope}
\end{tikzpicture}
\legende{Gamètes et structures haploïdes chez les Spermaphytes. À gauche : grain de pollen à deux noyaux haploïdes (noyau végétatif et noyau générateur), issu d'une microspore par mitose. À droite : sac embryonnaire à huit noyaux haploïdes : l'oosphère (gamète femelle), deux synergides, trois antipodes (ant.) et deux noyaux polaires centraux.}
\end{popfigure}

\begin{exercice}[Pour t'entraîner]
Chez le maïs, la cellule-mère du sac embryonnaire possède $2n = 20$ chromosomes. 1) Indiquer le nombre de chromosomes d'une mégaspore, de l'oosphère et d'un noyau polaire. 2) Combien de sacs embryonnaires fonctionnels obtient-on à partir de 50 cellules-mères du sac embryonnaire ? Comparer avec le nombre de grains de pollen obtenus à partir de 50 cellules-mères du pollen.
\end{exercice}

\begin{corrige}[Exercice]
1) La méiose divise le stock chromosomique par deux : mégaspore, oosphere et noyau polaire possèdent chacun $n = 10$ chromosomes (toutes les cellules du sac embryonnaire sont haploïdes). 2) Chaque cellule-mère du sac embryonnaire ne donne qu'\textbf{un seul} sac embryonnaire fonctionnel (trois mégaspores dégénèrent) : 50 sacs embryonnaires. Chaque cellule-mère du pollen donne \textbf{quatre} microspores, donc quatre grains de pollen : $50 \times 4 = 200$ grains de pollen. Comme chez les Mammifères, la gamétogenèse mâle a un rendement quatre fois supérieur.
\end{corrige}

\begin{attention}[Pièges de rédaction au Baccalauréat]
\begin{itemize}
\item Le grain de pollen n'est \textbf{pas} le gamète mâle : c'est une structure qui \emph{contient} (ou produira) les gamètes mâles ; le gamète femelle est l'oosphère, pas le sac embryonnaire entier.
\item Les mitoses qui forment le sac embryonnaire ou le grain de pollen \textbf{ne changent pas la ploïdie} : elles produisent des noyaux $n$ à partir de cellules $n$. Seule la méiose réduit la ploïdie.
\item Ne confonds pas « microspore » (plante, $n$) et « spermatide » (animal, $n$) : les termes sont spécifiques à chaque groupe.
\end{itemize}
\end{attention}

En définitive, que ce soit chez les Mammifères ou chez les Spermaphytes, la gamétogenèse aboutit au même résultat fondamental : la production de gamètes \textbf{haploïdes} et \textbf{génétiquement diversifiés} (grâce au brassage de la méiose). Reste à comprendre comment ces gamètes se rencontrent et restaurent la diploïdie : ce sera l'objet de la section suivante, consacrée à la fécondation, avec la remarquable double fécondation des Spermaphytes.

\coursec{La fécondation chez les Mammifères et la double fécondation chez les Spermaphytes}

Dans la section précédente, tu as vu comment la gamétogenèse produit, à partir des cellules de la lignée germinale, des gamètes haploïdes ($n$) : les spermatozoïdes et l'ovocyte chez les Mammifères, les noyaux spermatiques et l'oosphère chez les Spermaphytes. Ces gamètes, chargés chacun d'une quantité d'ADN égale à $Q/2$, sont des cellules à durée de vie limitée et incapables de se diviser seules. Leur destin est de se rencontrer : c'est la \cle{fécondation}, l'événement qui clôt le cycle de la reproduction sexuée et qui, en restaurant la diploïdie, assure la continuité de l'espèce. Nous allons décrire ce phénomène chez les Mammifères, puis chez les Spermaphytes où il présente une originalité remarquable : la \cle{double fécondation}.

\courssub{Définition générale de la fécondation}

\begin{definition}[Fécondation]
La \cle{fécondation} est l'union d'un gamète mâle haploïde ($n$) et d'un gamète femelle haploïde ($n$), qui donne naissance à une cellule unique, la \cle{cellule-œuf} ou \cle{zygote}, diploïde ($2n$). Elle comporte deux temps successifs : la \emph{plasmogamie} (fusion des cytoplasmes et des membranes plasmatiques) et la \emph{caryogamie} ou \cle{amphimixie} (fusion des deux noyaux et mélange de leurs patrimoines génétiques sur un fuseau commun).
\end{definition}

La fécondation réalise ainsi deux choses essentielles : elle \textbf{restaure le nombre diploïde de chromosomes} (divisé par deux lors de la méiose) et elle \textbf{brasse les patrimoines génétiques} de deux parents, créant un individu génétiquement nouveau. C'est le pendant obligatoire de la méiose : sans fécondation, la méiose conduirait à l'extinction de la lignée ; sans méiose, la fécondation doublerait le nombre de chromosomes à chaque génération.

\courssub{La fécondation chez les Mammifères}

\textbf{Situation concrète.} Dans une maternité de l'hôpital central de Yaoundé, on estime qu'une grossesse débute lorsqu'un seul spermatozoïde, parmi les quelques centaines de millions éjaculés, rencontre l'ovocyte. Comment cette rencontre a-t-elle lieu, et pourquoi un seul spermatozoïde parvient-il à féconder l'ovocyte ?

\textbf{Le lieu de la rencontre.} Lors du rapport sexuel, les spermatozoïdes sont déposés dans le vagin. Ils traversent le col de l'utérus, la cavité utérine, puis pénètrent dans les \cle{trompes de Fallope} (oviductes). C'est dans le \textbf{tiers externe de la trompe} (l'ampoule) que se fait la rencontre avec l'ovocyte II, expulsé de l'ovaire lors de l'ovulation et capté par le pavillon de la trompe. Sur des centaines de millions de spermatozoïdes, seuls quelques centaines atteignent l'ovocyte.

\textbf{Les étapes de la fécondation.} L'ovocyte II, bloqué en métaphase II de la méiose, est entouré de deux enveloppes : la \cle{zone pellucide} (matrice extracellulaire glycoprotéique, riche en ZP3, récepteur spécifique du spermatozoïde) et, plus à l'extérieur, la \cle{couronne radiée} (cellules folliculaires adhérant à la zone pellucide). Les étapes sont les suivantes :

\begin{enumerate}
\item \textbf{Approche et réaction acrosomique} : au contact de la couronne radiée, l'\cle{acrosome} (vésicule dérivée de l'appareil de Golgi, située à l'extrémité antérieure de la tête du spermatozoïde) libère ses enzymes (hyaluronidase pour dissocier les cellules de la couronne radiée, acrosine et protéases pour digérer la zone pellucide). C'est la \cle{réaction acrosomique}, déclenchée par la fixation de protéines de la membrane du spermatozoïde sur les récepteurs ZP3 de la zone pellucide.
\item \textbf{Franchissement des enveloppes et plasmogamie} : le spermatozoïde traverse la zone pellucide et sa membrane plasmique (région post-acrosomique) fusionne avec celle de l'ovocyte : c'est la plasmogamie.
\item \textbf{Blocage de la polyspermie} : dès la fusion membranaire, l'ovocyte déclenche une \cle{réaction corticale} : les \cle{granules corticaux} (vésicules sous-corticales) fusionnent avec la membrane plasmique et libèrent leur contenu (protéases, glycosaminoglycanes) dans l'espace perivitellin. Cela modifie la structure de la zone pellucide (coupure de ZP2, masquage de ZP3), la rendant impénétrable et non réceptrice pour les autres spermatozoïdes. Ainsi, \textbf{un seul spermatozoïde féconde l'ovocyte} (savoir admis).
\item \textbf{Entrée du noyau mâle et achèvement de la méiose II} : la tête du spermatozoïde (noyau + centrioles proximaux) pénètre dans le cytoplasme de l'ovocyte (le flagelle et les mitochondries restent généralement à l'extérieur ou sont dégradés). Cette entrée, via un facteur soluble (PLC-zêta), déclenche des oscillations calciques qui lèvent le blocage en métaphase II : l'ovocyte achève sa méiose II, expulse le \cle{second globule polaire} et devient l'\emph{ovule} définitif, à noyau haploïde (pronucleus femelle).
\item \textbf{Amphimixie} : le noyau du spermatozoïde se décondense et forme le \cle{pronucleus mâle} ; le noyau de l'ovule forme le \cle{pronucleus femelle}. Les deux pronuclei migrent l'un vers l'autre (guidés par les microtubules organisés par les centrioles paternels) ; leurs enveloppes nucléaires se désorganisent et leurs chromosomes s'alignent sur un fuseau mitotique commun : c'est la caryogamie (amphimixie). Le \cle{zygote} diploïde ($2n$) est formé et entre aussitôt en première division de segmentation (mitose).
\end{enumerate}

\begin{popfigure}
\begin{tikzpicture}[scale=0.95, every node/.style={font=\small}]
% Styles
\tikzset{
    zona/.style={thick, draw=poppink, fill=poppinkL},
    corona/.style={thick, draw=poporange, fill=poporangeL!30},
    cyto/.style={fill=popcream},
    sperm/.style={fill=popblue, draw=popdark},
    nucleus/.style={fill=poppurpleL, draw=poppurple, thick},
    pnuc/.style={fill=popblueL, draw=popblue, thick},
    arrow/.style={->, thick, popdark},
    label/.style={font=\scriptsize, align=center, popdark}
}

% PANEL 1 : Approche et réaction acrosomique
\begin{scope}[shift={(0,0)}]
% Corona radiata
\foreach \angle/\r in {30/1.3, 90/1.3, 150/1.3, 210/1.3, 270/1.3, 330/1.3}
    \node[corona, circle, minimum size=0.3cm] at (\angle:\r) {};
% Zona pellucida
\draw[zona] (0,0) circle (1.1);
% Perivitelline space
\draw[thin, dashed, popmuted] (0,0) circle (0.95);
% Oocyte membrane + cortex
\draw[thick, popdark, cyto] (0,0) circle (0.85);
% Cortical granules
\foreach \a in {10,50,90,130,170,210,250,290,330}
    \fill[popred] (\a:0.82) circle (0.04);
% Oocyte nucleus (Metaphase II plate)
\draw[nucleus] (0,0) ellipse (0.25 and 0.1);
\node[label] at (0,0) {Métaphase II};
% Spermatozoid approaching
\draw[sperm] (2.2,0.3) ellipse (0.3 and 0.12);
\draw[popblue, thick] (2.5,0.3) -- (3.0,0.4);
% Acrosome reaction visualization
\draw[poporange, decorate, decoration={snake, amplitude=1pt, segment length=3pt}] (1.9,0.3) -- (1.2,0.2);
\node[label, poporange, align=center] at (2.5,0.8){Libération enzymes\\ (hyaluronidase, acrosine)};
\node[label] at (0,-1.5) {1. Approche \& réaction acrosomique};
\node[label, left] at (-1.3,1.1) {Couronne radiée};
\node[label, left] at (-1.3,0.5) {Zone pellucide};
\node[label, left] at (-1.3,-0.1) {Espace perivitellin};
\node[label, left] at (-1.3,-0.7) {Granules corticaux};
\end{scope}

% PANEL 2 : Entrée, réaction corticale, fin méiose II
\begin{scope}[shift={(6.5,0)}]
% Zona
\draw[zona] (0,0) circle (1.1);
% Perivitelline space modified
\draw[thin, dashed, popmuted] (0,0) circle (0.95);
% Oocyte membrane
\draw[thick, popdark, cyto] (0,0) circle (0.85);
% Cortical granules released -> modified ZP
\draw[popred, thick] (0,1.1) -- (0,1.3) node[above, label, popred] {Réaction corticale : ZP modifiée};
% Sperm head fused
\draw[sperm] (0.1,0.1) ellipse (0.25 and 0.12);
% Male pronucleus forming
\draw[pnuc] (0.1,0.1) circle (0.18);
\node[label] at (0.1,0.1) {PNm};
% Centrioles
\draw[popblue, thick] (0.15,0.0) -- (0.25,-0.05) node[right, label, popblue] {Centrioles};
% Female meiosis II completing
\draw[nucleus] (-0.3,0.3) ellipse (0.2 and 0.1);
\node[label] at (-0.3,0.3) {Anaphase II};
% 2nd Polar body
\draw[nucleus] (-0.6,0.6) circle (0.12);
\node[label] at (-0.8,0.6) {2\textsuperscript{e} GP};
\node[label, align=center] at (0,-1.5){2. Plasmogamie, réaction corticale,\\fin méiose II $\to$ Ovule};
\end{scope}

% PANEL 3 : Amphimixie
\begin{scope}[shift={(13,0)}]
% Zona
\draw[zona] (0,0) circle (1.1);
% Membrane
\draw[thick, popdark, cyto] (0,0) circle (0.85);
% Two pronuclei
\draw[pnuc] (-0.25,0) circle (0.22);
\node[label] at (-0.25,0) {PNm ($n$)};
\draw[pnuc] (0.25,0) circle (0.22);
\node[label] at (0.25,0) {PNf ($n$)};
% Microtubules / Spindle forming
\draw[popblue, thick, ->] (-0.05,0.25) -- (-0.05,0.45);
\draw[popblue, thick, ->] (0.05,0.25) -- (0.05,0.45);
\node[label, popblue] at (0,0.6) {Fuseau mitotique commun};
% Centrosomes
\draw[popblue, fill=popblue] (-0.25,-0.3) circle (0.05);
\draw[popblue, fill=popblue] (0.25,-0.3) circle (0.05);
\node[label, align=center] at (0,-1.5){3. Amphimixie :\\Zygote $2n$};
\end{scope}

% Arrows between panels
\draw[arrow, very thick] (3.3,-1.2) -- (4.8,-1.2);
\draw[arrow, very thick] (9.8,-1.2) -- (11.3,-1.2);
\end{tikzpicture}
\legende{Étapes de la fécondation chez les Mammifères. \textbf{1} : Le spermatozoïde traverse la couronne radiée et la zone pellucide grâce à la réaction acrosomique. \textbf{2} : Fusion membranaire (plasmogamie), réaction corticale (blocage polyspermie), entrée du noyau mâle et des centrioles, achèvement de la méiose II (expulsion du 2\textsuperscript{e} globule polaire). \textbf{3} : Formation des deux pronuclei (mâle PNm, femelle PNf), migration, rupture des enveloppes et alignement sur un fuseau commun (amphimixie) $\to$ Zygote diploïde $2n$.}
\end{popfigure}

\begin{attention}[Pièges classiques au Bac : Mammifères]
\begin{itemize}
\item \textbf{Polyspermie} : N'écris jamais « plusieurs spermatozoïdes entrent ». Un seul pénètre grâce au blocage rapide (dépolarisation membranaire, secondes) et au blocage lent (réaction corticale, minutes) modifiant la ZP.
\item \textbf{Flagelle} : Seule la \textbf{tête} (noyau + centrioles) pénètre dans l'ovocyte. Le flagelle et la gaine mitochondriale restent à l'extérieur et sont dégradés.
\item \textbf{Vocabulaire précis} : Distingue \emph{ovocyte II} (cellule bloquée en métaphase II, fécondable), \emph{ovule} (cellule après expulsion du 2\textsuperscript{e} globule polaire, produit fini de l'ovogenèse) et \emph{zygote} (cellule diploïde issue de l'amphimixie).
\item \textbf{Centrioles} : Ils sont d'origine \textbf{paternelle} (apportés par le spermatozoïde) et organisent le premier fuseau mitotique du zygote.
\end{itemize}
\end{attention}

\courssub{La fécondation chez les Spermaphytes : la double fécondation}

\textbf{Situation concrète.} Dans une plantation de maïs de la région de l'Ouest, chaque grain de l'épi résulte de la fécondation d'un ovule. Or chez le maïs, les ovules sont au bout de longues « soies » (styles) pouvant dépasser \SI{30}{cm} : comment le gamète mâle, immobile et non flagellé, parvient-il jusqu'à l'oosphère ?

\textbf{Étape 1 : la pollinisation.} Contrairement aux Mammifères, les gamètes mâles des Spermaphytes (Angiospermes) ne nagent pas. Le gamète mâle est transporté à l'intérieur du \cle{grain de pollen} (gamétophyte mâle). La \cle{pollinisation} est le dépôt du grain de pollen sur le \cle{stigmate} (partie réceptrice du pistil), assurée par le vent (\cle{anémophilie} : maïs, graminées) ou par les insectes (\cle{entomophilie} : cotonnier, cacaoyer, manguier).

\textbf{Étape 2 : germination du pollen et formation du tube pollinique.} Sur le stigmate humide et réceptif, le grain de pollen germe : la \cle{intine} (paroi interne) s'extériorise par un pore et forme le \cle{tube pollinique}. Ce tube s'allonge par croissance apicale (extrémité) dans le style, guidé par des signaux chimiques (attractants) sécrétés par les synergides du sac embryonnaire, jusqu'au \cle{micropyle} de l'ovule. Pendant ce trajet, le \cle{noyau générateur} du grain de pollen (haploïde $n$) subit une \textbf{mitose} qui produit \textbf{deux noyaux spermatiques haploïdes ($n$)} : ce sont les véritables gamètes mâles. Le \cle{noyau végétatif} (ou tube) contrôle la croissance du tube mais ne participe pas à la fécondation.

\begin{exemplebox}[Performance du tube pollinique de maïs]
Chez le maïs (\emph{Zea mays}), les soies (styles) mesurent jusqu'à \SI{35}{cm}. Le tube pollinique croît à une vitesse moyenne de \SI{1,2}{cm.h^{-1}} (jusqu'à \SI{2}{cm.h^{-1}}), atteignant l'ovule en \SI{12}{h} à \SI{24}{h}. C'est l'une des croissances cellulaires les plus rapides du monde végétal, nécessitant un métabolisme intense et un apport continu en sucres par le style.
\end{exemplebox}

\textbf{Étape 3 : la double fécondation.} Le tube pollinique pénètre dans le \cle{sac embryonnaire} (gamétophyte femelle, type \emph{Polygonum} à 7 cellules / 8 noyaux le plus fréquent) par le micropyle, généralement entre les deux synergides, et y éclate, libérant les deux noyaux spermatiques. Se produisent alors \textbf{deux fusions distinctes et quasi simultanées} :

\begin{enumerate}
\item Le \textbf{premier noyau spermatique} ($n$) fusionne avec le noyau de l'\cle{oosphère} ($n$) : c'est la fécondation proprement dite, qui donne le \cle{zygote} diploïde ($2n$), futur embryon.
\item Le \textbf{second noyau spermatique} ($n$) fusionne avec les \textbf{deux noyaux polaires} ($n + n$) de la \cle{cellule centrale} : il en résulte un noyau \textbf{triploïde ($3n$)}, noyau initial de l'\cle{endosperme} (ou albumen en développement).
\end{enumerate}

\begin{definition}[Double fécondation]
La \cle{double fécondation} est le phénomène propre aux Angiospermes (Spermaphytes) au cours duquel les deux noyaux spermatiques d'un même grain de pollen réalisent deux fusions distinctes dans le sac embryonnaire : l'un avec l'oosphère forme le zygote diploïde ($2n$), l'autre avec les deux noyaux polaires forme un noyau triploïde ($3n$) à l'origine de l'endosperme (albumen).
\end{definition}

\begin{propriete}[Rôle de l'endosperme / albumen — admis]
L'\cle{endosperme} (tissu triploïde $3n$ issu de la seconde fusion) constitue le \textbf{tissu nutritif de réserve} de l'embryon dans la graine (amidon, protéines, lipides). Il se développe \emph{après} la fécondation, ce qui garantit qu'aucune réserve n'est gaspillée si la fécondation échoue. Selon les espèces, il persiste à maturité (graines \emph{albuminées} : maïs, blé, ricin) ou est entièrement consommé par l'embryon qui stocke les réserves dans ses cotylédons (graines \emph{exalbuminées} : haricot, arachide, cacaoyer).
\end{propriete}

\begin{savaistu}[Une découverte russe... au jardin des lis]
La double fécondation a été découverte en \textbf{1898} par le botaniste russe \textbf{Sergueï Nawaschin} (et indépendamment par Léon Guignard en France), en observant au microscope des ovules de lis (\emph{Lilium martagon}). Il y vit deux noyaux spermatiques réaliser deux fusions simultanées dans le même sac embryonnaire — une innovation évolutive majeure qui distingue les Angiospermes de tous les autres végétaux (Gymnospermes, Fougères, Bryophytes) où il n'y a qu'une seule fusion.
\end{savaistu}

\begin{popfigure}
\begin{tikzpicture}[scale=1.05, every node/.style={font=\small}]
% Styles
\tikzset{
    ovule/.style={thick, draw=popdark, fill=popgreenL!20},
    sac/.style={thick, draw=popgreen, fill=popcream},
    tube/.style={very thick, draw=popblue},
    nuc/.style={fill=popblue, draw=popdark, thick, circle, minimum size=0.28cm},
    nucf/.style={fill=poppinkL, draw=poppink, thick, circle, minimum size=0.35cm},
    nucp/.style={fill=poppurpleL, draw=poppurple, thick, circle, minimum size=0.25cm},
    arrowf/.style={->, very thick, poporange},
    label/.style={font=\scriptsize, align=center, popdark},
    labell/.style={font=\scriptsize, align=left, popdark}
}

% Ovule outer shape
\draw[ovule] (0,0) ellipse (2.8 and 3.6);
% Embryo sac
\draw[sac] (0,0.1) ellipse (1.9 and 2.8);

% Micropyle
\draw[popdark, fill=white, thick] (0,-3.5) circle (0.2);
\node[label] at (0,-3.9) {Micropyle};

% CHALAZAL END (top) : Antipodes (3 cells) - often degenerate, optional to draw for clarity
% \node[label] at (0,2.8) {Pôle chalazal (antipodes)};

% CENTRAL CELL : 2 Polar Nuclei (often fused before fertilization)
\draw[nucp] (-0.4,1.1) circle (0.22);
\draw[nucp] (0.4,1.1) circle (0.22);
\node[label] at (0,1.1) {$n$ \quad $n$};
\node[label, align=center] at (0,1.6){2 noyaux polaires\\(cellule centrale)};

% MICROPYLAR END (bottom) : Egg Apparatus
% Synergids (2) - flanking the egg
\draw[fill=poppinkL!50, draw=poppink, thick] (-0.55,-1.8) circle (0.22);
\draw[fill=poppinkL!50, draw=poppink, thick] (0.55,-1.8) circle (0.22);
\node[label] at (-0.55,-1.8) {Sy};
\node[label] at (0.55,-1.8) {Sy};
\node[label] at (-1.2,-1.8) {Synergides};

% Oosphere (Egg cell)
\draw[nucf] (0,-2.3) circle (0.35);
\node[label] at (0,-2.3) {Oosphère $n$};

% POLLEN TUBE
\draw[tube] (0.15,-6.0) -- (0.15,-3.5) -- (0.15,-2.0);
\draw[tube] (0.45,-6.0) -- (0.45,-3.5) -- (0.45,-2.0);
\draw[tube] (0.15,-2.0) -- (0.45,-2.0); % burst tip
\node[labell, popblue, align=center] at (0.7,-4.0){Tube pollinique\\(croissance apicale,\\noyau végétatif)};

% SPERM NUCLEI released
\draw[nuc] (0.3,-1.7) circle (0.14);
\draw[nuc] (0.3,-1.4) circle (0.14);
\node[labell, popblue, align=center] at (0.7,-1.55){2 noyaux\\spermatiques $n$};

% FUSION ARROWS
% Fusion 1: Sperm + Egg
\draw[arrowf] (0.3,-1.65) -- (0.1,-2.1);
% Fusion 2: Sperm + Polar nuclei
\draw[arrowf] (0.5,-1.5) .. controls (1.2,-0.5) .. (0.5,1.1);

% PRODUCTS LABELS
% Zygote
\node[label, align=center] at (-3.8,-2.3) {Zygote $2n$\\$\to$ Embryon};
\draw[->, thick, popdark] (-2.8,-2.3) -- (-0.7,-2.3);

% Endosperm nucleus (3n)
\node[label, align=center] at (4.2,1.1) {Noyau $3n$\\$\to$ Endosperme\\(Albumen)};
\draw[->, thick, popdark] (3.2,1.1) -- (1.1,1.1);

% FUTURE STRUCTURES
\node[label, font=\small\bfseries, popdark] at (0,4.2) {Ovule $\to$ Graine \quad | \quad Ovaire $\to$ Fruit};
\end{tikzpicture}
\legende{Schéma de la double fécondation chez les Spermaphytes (type \emph{Polygonum}). Le tube pollinique (bleu) chemine dans le style, pénètre dans le sac embryonnaire par le micropyle et libère deux noyaux spermatiques ($n$). \textbf{Fusion 1 (orange)} : noyau spermatique + oosphère ($n$) $\to$ zygote ($2n$), futur embryon. \textbf{Fusion 2 (orange)} : noyau spermatique + 2 noyaux polaires ($n+n$) $\to$ noyau triploïde ($3n$), à l'origine de l'endosperme (albumen). Sy = Synergides.}
\end{popfigure}

\begin{attention}[Pièges classiques au Bac : Spermaphytes]
\begin{itemize}
\item \textbf{Tube pollinique $\neq$ Gamète} : Erreur majeure : « le tube pollinique féconde l'oosphère ». Le tube est une \textbf{structure de transport} (haploïde $n$, issu du noyau végétatif). Les \textbf{gamètes mâles sont les deux noyaux spermatiques} ($n$) issus de la mitose du noyau générateur.
\item \textbf{Grain de pollen $\neq$ Gamète} : Le grain de pollen est le \textbf{gamétophyte mâle} (organisme haploïde multicellulaire : 2 ou 3 cellules). Il \emph{produit} et \emph{transporte} les gamètes.
\item \textbf{Noyaux polaires} : Ils sont au nombre de \textbf{deux} (souvent fusionnés en un noyau diploïde $2n$ avant l'arrivée du tube). La fusion avec le noyau spermatique ($n$) donne bien $3n$.
\item \textbf{Double fécondation = Double fusion} : Ne pas confondre avec « deux fécondations successives » ou « fécondation de deux ovules ». C'est un seul grain de pollen, deux fusions simultanées dans le même sac embryonnaire.
\item \textbf{Devenir} : Zygote $\to$ Embryon ; Noyau $3n$ $\to$ Endosperme ; Intégruments (téguments de l'ovule) $\to$ Téguments de la graine ; Paroi de l'ovaire $\to$ Péricarpe (fruit).
\end{itemize}
\end{attention}

\textbf{Devenir post-fécondation.} Après la double fécondation, les transformations s'enchaînent sans discontinuité : le \textbf{zygote} se divise par mitoses et se différencie en \cle{embryon} (radicule, plumule, cotylédon(s)) ; le noyau $3n$ se divise et forme l'\cle{endosperme} (albumen) ; l'\textbf{ovule} entier se transforme en \cle{graine} (ses téguments deviennent les téguments de la graine, le funicule devient le pédoncule) ; l'\textbf{ovaire} se transforme en \cle{fruit} (sa paroi devient le péricarpe, charnu ou sec). Manger un grain de maïs, une fève de cacao ou une tomate, c'est consommer le produit direct d'une double fécondation.

\courssub{Comparaison des deux modèles de fécondation}

\begin{center}
\begin{tabularx}{\linewidth}{|l|X|X|}
\hline
\rowcolor{popblueL} \textbf{Critère} & \textbf{Mammifères} & \textbf{Spermaphytes (Angiospermes)} \\
\hline
Lieu de la rencontre & Tiers externe de la trompe de Fallope (ampoule) & Sac embryonnaire de l'ovule, dans l'ovaire \\
\hline
Transport du gamète mâle & Flagelle du spermatozoïde (gamète mobile, chimiotactisme) & Pollinisation (vent/insectes) puis tube pollinique (gamète immobile transporté) \\
\hline
Nature du gamète mâle & Spermatozoïde (cellule différenciée, flagellée, $n$) & Deux noyaux spermatiques ($n$) issus de la mitose du noyau générateur (pas de cytoplasme propre) \\
\hline
Nombre de fusions & Une seule (fécondation simple) & Deux (double fécondation) \\
\hline
Produits nucléaires formés & Zygote ($2n$) uniquement & Zygote ($2n$) + Noyau initial de l'endosperme ($3n$) \\
\hline
Déclenchement méiose femelle & L'entrée du spermatozoïde achève la méiose II & La méiose est achevée \emph{avant} la fécondation (oosphère déjà $n$) \\
\hline
Blocage polyspermie & Réaction corticale (modification Zone Pellucide) & Éclatement du tube pollinique + modifications du sac embryonnaire (synergides) \\
\hline
Devenir immédiat & Zygote $\to$ segmentation dans la trompe $\to$ blastocyste $\to$ implantation & Zygote $\to$ embryon ; Ovule $\to$ graine ; Ovaire $\to$ fruit \\
\hline
\end{tabularx}
\end{center}

Dans les deux cas, l'essentiel est conservé : \textbf{deux noyaux haploïdes fusionnent pour restaurer la diploïdie ($2n$)}. La double fécondation est une innovation évolutive des Angiospermes qui garantit la production d'un tissu nutritif (endosperme) \emph{uniquement si la fécondation a lieu} — une économie de ressources remarquable comparée aux Gymnospermes où le tissu nutritif (haploïde, gamétophyte femelle) se développe avant la fécondation.

\courssub{Évolution de la quantité d'ADN au cours de la fécondation}

Suivons la quantité d'ADN \emph{par noyau}, en prenant $Q$ comme quantité d'ADN d'une cellule somatique diploïde en phase G1 (avant réplication). Chaque gamète haploïde ($n$) contient une quantité d'ADN égale à $Q/2$. Au moment de l'amphimixie, les deux patrimoines se réunissent dans un cytoplasme commun : le zygote contient $Q/2 + Q/2 = Q$. Le zygote entre alors immédiatement dans le cycle cellulaire : la phase S (réplication de l'ADN) fait passer la quantité de $Q$ à $2Q$ (chaque chromosome est constitué de deux chromatides) avant la première division de segmentation (mitose), qui redistribue l'ADN pour ramener la quantité à $Q$ dans chaque cellule-fille (blastomères).

\begin{popfigure}
\begin{center}
\begin{tikzpicture}
\begin{axis}[
width=0.9\linewidth, height=6.5cm,
xlabel={Temps (étapes successives)},
ylabel={Quantité d'ADN par noyau},
xmin=0, xmax=10, ymin=0, ymax=2.3,
xtick={1,3,5,7,9},
xticklabels={Gamètes haploïdes, Amphimixie, Zygote (G1), Fin Phase S, 1\iere{} segmentation},
ytick={0.5,1,2},
yticklabels={$Q/2$,$Q$,$2Q$},
grid=major, grid style={popline!50},
legend style={at={(0.02,0.98)},anchor=north west,font=\scriptsize},
clip=false
]
% Step function: constant then jump
\addplot[popcurve, const plot, very thick] coordinates {
    (0,0.5) (2,0.5)   % Gametes
    (2,1) (4,1)       % Fusion -> Zygote G1
    (4,1) (6,1)       % G1
    (6,2) (8,2)       % S phase -> G2/M
    (8,1) (10,1)      % Mitosis -> Daughter cells G1
};
\addplot[only marks, mark=*, poporange, mark size=2.5pt] coordinates {
    (2,0.5) (2,1) (6,1) (6,2) (8,2) (8,1)
};
\legend{Quantité d'ADN par noyau}
% Annotations
\node[font=\scriptsize, popdark, align=center] at (axis cs:2,1.25) {Fusion\\$Q/2+Q/2=Q$};
\node[font=\scriptsize, popdark, align=center] at (axis cs:6,2.25) {Réplication ADN\\(Phase S)};
\node[font=\scriptsize, popdark, align=center] at (axis cs:8.8,0.72) {Mitose :\\$2Q \to Q$};
\node[font=\scriptsize, popgreen, align=center] at (axis cs:1,0.25) {Spermatozoïde\\$n, Q/2$};
\node[font=\scriptsize, poppink, align=center] at (axis cs:1,0.25) {Ovocyte II\\$n, Q/2$};
\end{axis}
\end{tikzpicture}
\end{center}
\legende{Évolution de la quantité d'ADN au cours de la fécondation et de la première segmentation. Chaque gamète haploïde ($n$) apporte $Q/2$. L'amphimixie restaure $Q$ dans le zygote ($2n$, G1). La réplication (phase S) porte la quantité à $2Q$ (chromosomes bichromatidiens). La première division de segmentation (mitose) redistribue $Q$ à chaque cellule-fille. Le cycle cellulaire somatique a repris : l'alternance $n \xrightarrow{\text{fécondation}} 2n \xrightarrow{\text{méiose}} n$ est bouclée.}
\end{popfigure}

\begin{exoresolu}[Exploiter une courbe d'évolution de la quantité d'ADN]
\textbf{Énoncé.} Chez une espèce animale, la quantité d'ADN d'une cellule somatique en phase G1 est $Q = \num{6,6}$ pg (picogrammes). 
\begin{enumerate}
\item Quelle est la quantité d'ADN d'un spermatozoïde et d'un ovocyte II de cette espèce ? Justifie.
\item Quelle est la quantité d'ADN du zygote juste après l'amphimixie, puis à la fin de la première phase S (début G2) ?
\item Explique pourquoi la méiose et la fécondation sont deux phénomènes complémentaires indispensables à la pérennité de l'espèce.
\end{enumerate}
\tcblower
\textbf{Solution détaillée.}
\begin{enumerate}
\item Les gamètes (spermatozoïde et ovocyte II) sont haploïdes ($n$), issus de la méiose. La méiose divise par deux le nombre de chromosomes et la quantité d'ADN par rapport à la cellule souche diploïde en G1 ($2n, Q$). Donc chaque gamète contient \textbf{$Q/2 = \num{3,3}$ pg d'ADN}.
\item Juste après l'amphimixie (fusion des deux noyaux haploïdes), le zygote est diploïde ($2n$) et se trouve en phase G1 (pas encore de réplication). Sa quantité d'ADN est la somme des deux apports : \textbf{$Q/2 + Q/2 = Q = \num{6,6}$ pg}. À la fin de la phase S (réplication de l'ADN), chaque chromosome est dupliqué (deux chromatides) : la quantité d'ADN double. Elle vaut \textbf{$2Q = \num{13,2}$ pg}. Elle retombera à $Q$ dans chaque blastomère après la première mitose de segmentation.
\item \textbf{Complémentarité :} 
\begin{itemize}
\item La \textbf{méiose} réduit le nombre de chromosomes ($2n \to n$) et la quantité d'ADN ($Q \to Q/2$) lors de la formation des gamètes. Sans elle, la fécondation doublerait le caryotype à chaque génération ($2n + 2n \to 4n$, puis $8n$, etc.), ce qui est incompatible avec la viabilité.
\item La \textbf{fécondation} restaure le nombre de chromosomes ($n + n \to 2n$) et la quantité d'ADN ($Q/2 + Q/2 \to Q$) lors de la formation du zygote. Sans elle, les gamètes haploïdes ne pourraient pas développer un nouvel individu diploïde (chez la plupart des Mammifères et Spermaphytes).
\end{itemize}
Ensemble, elles maintiennent \textbf{constant le caryotype de l'espèce} (nombre et forme des chromosomes, quantité d'ADN) à travers les générations successives, tout en brassant le patrimoine génétique (croisement, assortiment indépendant, fécondation aléatoire). C'est le fondement cellulaire et moléculaire de la \textbf{pérennité de l'espèce}.
\end{enumerate}
\end{exoresolu}

\begin{aretenir}[L'essentiel de la section : Fécondation et Pérennité]
\begin{itemize}
\item \textbf{Fécondation} = Union de deux gamètes haploïdes ($n$) $\to$ Zygote diploïde ($2n$). Deux temps : Plasmogamie (fusion cytoplasmes) + Amphimixie (fusion noyaux).
\item \textbf{Mammifères} : Rencontre dans la trompe (ampoule). Réaction acrosomique (enzymes) $\to$ Franchissement Zone Pellucide $\to$ Plasmogamie $\to$ Réaction corticale (blocage polyspermie) $\to$ Entrée noyau mâle + centrioles $\to$ Fin méiose II (2\textsuperscript{e} globule polaire) $\to$ Amphimixie (pronuclei) $\to$ Zygote $2n$.
\item \textbf{Spermaphytes} : Pollinisation $\to$ Germination pollen $\to$ Tube pollinique (transport, noyau végétatif) $\to$ Mitose noyau générateur $\to$ 2 noyaux spermatiques ($n$). \textbf{Double fécondation} (Nawaschin, 1898) : 1\textsuperscript{er} noyau $n$ + Oosphère $n$ $\to$ Zygote $2n$ (Embryon) ; 2\textsuperscript{e} noyau $n$ + 2 noyaux polaires $n+n$ $\to$ Noyau $3n$ (Endosperme/Albumen). Devenir : Ovule $\to$ Graine, Ovaire $\to$ Fruit.
\item \textbf{Quantité d'ADN} : Gamètes $Q/2$ $\xrightarrow{\text{amphimixie}}$ Zygote $Q$ (G1) $\xrightarrow{\text{Phase S}}$ $2Q$ (G2/M) $\xrightarrow{\text{Mitose}}$ $Q$ par cellule-fille.
\item \textbf{Pérennité} : Méiose (réduction) + Fécondation (restauration) = Constante du caryotype intergénérationnelle + Brassage génétique (diversité).
\end{itemize}
\end{aretenir}

\coursec{Méiose et fécondation : deux phénomènes complémentaires au service de la pérennité de l'espèce}

Dans les sections précédentes, nous avons décrit la fécondation chez les Mammifères — fusion d'un spermatozoïde et d'un ovocyte II pour former un zygote diploïde — et la double fécondation chez les Spermaphytes, où deux noyaux mâles fécondent respectivement l'oosphère (noyau + noyau mâle $\rightarrow$ zygote $2n$) et la cellule centrale (deux noyaux polaires + noyau mâle $\rightarrow$ noyau tertiaire $3n$, futur albumen). Une question de synthèse s'impose : pourquoi la nature a-t-elle sélectionné deux phénomènes aussi complexes que la méiose et la fécondation, alors que la mitose suffit à multiplier les cellules ? Cette section conclusives répond à cette question : la \cle{pérennité de l'espèce} repose sur la complémentarité obligée de ces deux phénomènes.

\courssub{Raisonnement par l'absurde : que se passerait-il sans méiose ou sans fécondation ?}

Pour saisir la nécessité de chaque phénomène, adoptons une démarche logique classique en science : le \cle{raisonnement par l'absurde}. On suppose qu'un des deux phénomènes disparaît, et l'on examine les conséquences sur le stock chromosomique.

\textbf{Première hypothèse : suppression de la méiose.} Si les gamètes étaient produits par mitose, ils conserveraient le stock chromosomique diploïde ($2n$). À la fécondation, le zygote recevrait $2n + 2n = 4n$ chromosomes ; la génération suivante aurait $8n$, puis $16n$... Chez l'Homme ($2n = 46$), on passerait de 46 à 92, puis 184, puis 368 chromosomes en trois générations. Le stock chromosomique doublerait à chaque génération : c'est matériellement impossible (encombrement nucléaire, durée de réplication, ségrégation mitotique défaillante), et l'espèce disparaîtrait. La méiose est donc \textbf{indispensable} pour réduire de moitié le nombre de chromosomes avant chaque fécondation.

\textbf{Seconde hypothèse : suppression de la fécondation.} Si la méiose existait seule, elle produirait indéfiniment des cellules haploïdes ($n$). Or chez les espèces diploïdes comme les Mammifères et les Spermaphytes, un individu haploïde ne possède qu'un seul exemplaire de chaque gène : tout allèle délétère récessif s'exprime immédiatement (pas d'allèle « de secours » dominant). De plus, l'individu haploïde ne peut pas réaliser de méiose (absence de chromosomes homologues à apparier en prophase I) : il est stérile. La lignée s'éteint. La fécondation est donc \textbf{indispensable} pour restaurer la diploïdie.

\begin{aretenir}[Conclusion du raisonnement par l'absurde]
Sans méiose, le nombre de chromosomes doublerait à chaque génération (impossible physiquement). Sans fécondation, la méiose conduirait à des individus haploïdes non viables et stériles chez les espèces diploïdes. Les deux phénomènes sont donc \cle{complémentaires} et \cle{indissociables} : aucun ne peut assurer seul la reproduction sexuée.
\end{aretenir}

\courssub{Le cycle chromosomique : une alternance rigoureuse $2n \rightleftarrows n$}

La vie d'une espèce à reproduction sexuée peut se résumer à un cycle chromosomique simple et parfaitement équilibré :
\[ \text{Zygote } (2n) \xrightarrow{\text{développement (mitoses)}} \text{Adulte } (2n) \xrightarrow{\text{méiose}} \text{Gamètes } (n) \xrightarrow{\text{fécondation}} \text{Zygote } (2n) \]

La méiose divise par deux le nombre de chromosomes ($2n \rightarrow n$) ; la fécondation le multiplie par deux ($n + n \rightarrow 2n$). Les deux phénomènes se \textbf{compensent exactement} : la méiose est le « moins » qui équilibre le « plus » de la fécondation. Le caryotype de l'espèce reste ainsi constant au fil des générations.

\begin{popfigure}
\begin{center}
\begin{tikzpicture}[scale=1, every node/.style={font=\small}]
% Cercle du cycle
\draw[popline, thick] (0,0) circle (2.8);
% Noeuds
\node[draw=popblue, fill=popblueL, rounded corners, thick, minimum width=2.8cm, minimum height=1.0cm] (zyg) at (90:2.8) {\textbf{Zygote} $2n$};
\node[draw=popgreen, fill=popgreenL, rounded corners, thick, minimum width=2.8cm, minimum height=1.0cm] (adu) at (-30:2.8) {\textbf{Adulte} $2n$};
\node[draw=poporange, fill=poporangeL, rounded corners, thick, minimum width=2.8cm, minimum height=1.0cm] (gam) at (210:2.8) {\textbf{Gamètes} $n$};
% Flèches courbes
\draw[-{Stealth[length=3mm]}, very thick, popgreen] (zyg.east) .. controls (3.6,1.8) and (3.6,-0.5) .. (adu.north) node[midway, right, align=left, text=popdark]{Développement\\ (mitoses successives)};
\draw[-{Stealth[length=3mm]}, very thick, poporange] (adu.west) .. controls (0.5,-2.5) .. (gam.east) node[midway, below, text=popdark]{\textbf{Méiose} : $2n \rightarrow n$};
\draw[-{Stealth[length=3mm]}, very thick, poppurple] (gam.north) .. controls (-3.6,-0.5) and (-3.6,1.8) .. (zyg.west) node[midway, left, align=right, text=popdark]{\textbf{Fécondation}\\ $n + n \rightarrow 2n$};
% Centre
\node[align=center, text=popdark, font=\normalsize\bfseries] at (0,0) {\textbf{Cycle}\\ \textbf{chromosomique}\\ $2n \rightleftarrows n$};
\end{tikzpicture}
\end{center}
\legende{Schéma circulaire du cycle chromosomique. La méiose (flèche orange) réduit le stock de $2n$ à $n$ ; la fécondation (flèche violette) le restaure à $2n$ ; le développement de l'embryon en adulte se fait par mitoses (flèche verte), qui conservent $2n$ à chaque division.}
\end{popfigure}

\courssub{Courbe synthétique : l'évolution de la quantité d'ADN sur un cycle vital complet}

Nous avons étudié séparément la courbe d'ADN de la méiose et celle de la fécondation. Rassemblons-les sur un cycle vital complet, de la gamétogenèse d'un parent jusqu'à l'adulte de la génération suivante. On note $q$ la quantité d'ADN d'une cellule haploïde ($n$ chromosomes à une chromatide) ; la cellule diploïde en $G_1$ contient $2q$, et $4q$ après réplication (fin de phase $S$).

\begin{popfigure}
\begin{center}
\begin{tikzpicture}
\begin{axis}[
  width=0.98\linewidth, height=7.5cm,
  xlabel={Étapes du cycle vital}, ylabel={Quantité d'ADN par noyau},
  xmin=0, xmax=10, ymin=0, ymax=5,
  xtick={0.5,1.5,2.5,3.5,4.5,5.5,6.5,7.5,8.5,9.5},
  xticklabels={$G_1$, $S$, $G_2$, Méiose I, Méiose II, Gamète, Fécondation, $S$ zygote, Mitoses, Adulte},
  x tick label style={rotate=35, anchor=east, font=\scriptsize},
  ytick={1,2,4}, yticklabels={$q$, $2q$, $4q$},
  grid=major, grid style={popline!40},
  clip=false,
]
% Courbe continue avec sauts verticaux aux divisions/fécondation
\addplot[popcurve, thick, mark=none] coordinates {
(0,2) (1,2) (2,4) (3,4) (3,2) (4,2) (4,1) (5,1) (6,1) (6,2) (7,2) (8,4) (9,4) (9,2) (10,2)
};
% Annotations
\node[font=\scriptsize, text=popdark, align=center] at (axis cs:0.5,2.6) {$G_1$ : $2q$};
\node[font=\scriptsize, text=popdark, align=center] at (axis cs:2.5,4.4) {Réplication\\ ($2q \rightarrow 4q$)};
\node[font=\scriptsize, text=popdark, align=center] at (axis cs:3.5,3.1) {Méiose I\\ ($4q \rightarrow 2q$)};
\node[font=\scriptsize, text=popdark, align=center] at (axis cs:4.5,1.6) {Méiose II\\ ($2q \rightarrow q$)};
\node[font=\scriptsize, text=popdark, align=center] at (axis cs:5.5,0.5) {Gamète $n$\\ ($q$)};
\node[font=\scriptsize, text=popdark, align=center] at (axis cs:6.5,2.6) {Fécondation\\ ($q \rightarrow 2q$)};
\node[font=\scriptsize, text=popdark, align=center] at (axis cs:8.5,3.1) {Mitoses : alternance\\ $2q \rightleftarrows 4q$};
\end{axis}
\end{tikzpicture}
\end{center}
\legende{Courbe synthétique de la quantité d'ADN par noyau au cours d'un cycle vital complet. Phase $G_1$ parentale : $2q$ ; phase $S$ : réplication continue ($2q \rightarrow 4q$) ; méiose I : chute brutale à $2q$ (ségrégation homologues) ; méiose II : chute à $q$ (ségrégation chromatides) ; fécondation : saut à $2q$ (fusion noyaux) ; phase $S$ zygotique : réplication à $4q$ ; mitoses du développement : oscillations $2q \leftrightarrow 4q$ avec retour à $2q$ en $G_1$ de l'adulte.}
\end{popfigure}

\begin{attention}[Lecture de la courbe : pièges fréquents]
Ne confonds pas \textbf{quantité d'ADN} et \textbf{nombre de chromosomes} :
\begin{itemize}
\item Après la réplication (fin $S$, $G_2$, prophase I), la cellule contient $4q$ d'ADN mais toujours $2n$ chromosomes (chacun constitué de deux chromatides).
\item À l'issue de la méiose I, la quantité d'ADN est $2q$ mais le nombre de chromosomes est $n$ (chacun à deux chromatides).
\item Le saut de $q$ à $2q$ à la fécondation correspond au passage de $n$ à $2n$ chromosomes (chacun à une chromatide).
\item La quantité d'ADN ne redescend \emph{jamais} en dessous de $q$ dans le cycle : $q$ est le plancher haploïde.
\end{itemize}
\end{attention}

\courssub{Le rôle conservateur : la stabilité du caryotype de l'espèce}

Grâce à la compensation exacte entre méiose et fécondation, chaque espèce conserve, de génération en génération, un \cle{caryotype} constant : nombre, forme et taille des chromosomes. L'Homme garde $2n = 46$ chromosomes, le maïs $2n = 20$, le chimpanzé $2n = 48$. Cette constance garantit la transmission intégrale de l'\cle{information génétique} caractéristique de l'espèce : un enfant humain naît humain, un plant de maïs donne du maïs. C'est le premier pilier de la pérennité : la \textbf{stabilité}.

\courssub{Le rôle créateur de diversité : des individus tous différents}

Si la reproduction ne faisait que conserver, tous les individus d'une espèce seraient génétiquement identiques (des clones). Or il suffit d'observer une fratrie pour constater le contraire. D'où vient cette diversité ? De trois sources combinées, dont deux opèrent \emph{pendant la méiose} :

\begin{enumerate}
\item Le \cle{brassage interchromosomique} : à l'anaphase I de la méiose, les paires de chromosomes homologues se répartissent au hasard vers les pôles. Une cellule à $2n$ chromosomes (donc $n$ paires) peut produire $2^{n}$ types de gamètes différents par ce seul mécanisme. Chez l'Homme ($n=23$), $2^{23} \approx \num{8,4}$ millions de gamètes distincts par individu.
\item Le \cle{brassage intrachromosomique} (ou \cle{crossing-over}) : lors de la prophase I, des échanges de fragments (chiasmas) ont lieu entre chromatides non-sœurs de chromosomes homologues, créant des chromosomes « mosaïques » inédits. Le nombre de combinaisons alléliques devient alors pratiquement infini.
\item La \cle{rencontre aléatoire des gamètes} à la fécondation : n'importe quel spermatozoïde (parmi des millions) peut féconder l'ovule (unique par cycle chez la femme).
\end{enumerate}

\begin{exemplebox}[Un calcul vertigineux chez l'Homme]
Pour un couple humain, sans même compter les crossing-over : le père peut produire $2^{23}$ types de spermatozoïdes, la mère $2^{23}$ types d'ovocytes. Le nombre de zygotes génétiquement distincts possibles est donc :
\[ 2^{23} \times 2^{23} = 2^{46} \approx \num{7,0e13} \]
soit plus de \textbf{70 000 milliards} de combinaisons possibles. Avec les crossing-over, ce nombre devient astronomique : à part les vrais jumeaux (issus d'un même zygote), deux êtres humains génétiquement identiques n'ont jamais existé et n'existeront jamais.
\end{exemplebox}

\begin{attention}[Erreur fréquente au Bac]
N'affirme jamais que « la fécondation crée la diversité génétique ». La fécondation ne \textbf{crée} rien : elle \textbf{combine} au hasard des gamètes déjà diversifiés par la méiose (brassages inter- et intrachromosomiques). La source première de la diversité est la méiose ; la fécondation en amplifie l'expression. Formulation correcte attendue : « la diversité génétique résulte des brassages méiotiques \emph{et} de la rencontre aléatoire des gamètes à la fécondation ».
\end{attention}

\courssub{Diversité génétique et adaptation : une assurance-vie pour l'espèce}

À quoi sert cette diversité ? Elle est la \textbf{matière première de l'adaptation et de l'évolution}. Lorsque le milieu change (climat, parasites, maladies, pression anthropique), certains individus, grâce à leurs combinaisons génétiques particulières, survivent et se reproduisent mieux que d'autres : l'espèce, elle, traverse la crise.

\begin{exemplebox}[Le maïs au Cameroun : la diversité qui protège]
Dans les champs de l'Ouest et du Nord-Ouest camerounais, les paysans cultivent des variétés de maïs génétiquement diverses (variétés locales, variétés améliorées comme CMS 8704 ou ATP). Face à un parasite comme la chenille légionnaire d'automne (\emph{Spodoptera frugiperda}), qui ravage les cultures depuis 2016, toutes les plantes ne réagissent pas de la même façon : certaines combinaisons génétiques confèrent une résistance partielle. Une parcelle génétiquement uniforme (monoculture clonale) pourrait être anéantie en une saison ; une population diverse conserve toujours des individus capables de produire des graines. La diversité génétique issue de la reproduction sexuée est donc une véritable \textbf{assurance-vie} pour l'espèce — et pour la sécurité alimentaire des populations.
\end{exemplebox}

\courssub{Quand la méiose déraille : la leçon de la trisomie 21 (complément hors programme strict)}

L'importance capitale du bon déroulement de la méiose apparaît de façon frappante lorsqu'elle connaît un accident. Ce complément, bien qu'au-delà du strict programme de Terminale D, est un excellent argument de synthèse au Bac.

\begin{savaistu}[La trisomie 21, conséquence d'une non-disjonction méiotique]
La \cle{trisomie 21} (syndrome de Down) résulte le plus souvent d'une \textbf{non-disjonction} : les deux chromosomes 21 (homologues en anaphase I, ou chromatides sœurs en anaphase II) ne se séparent pas lors de la méiose, le plus fréquemment lors de l'ovogenèse chez la mère. Il en résulte un gamète à $n + 1 = 24$ chromosomes, porteur de deux chromosomes 21. Après fécondation par un gamète normal ($n = 23$), le zygote possède $2n + 1 = 47$ chromosomes, dont trois chromosomes 21. Cet exemple montre que la méiose n'assure la pérennité de l'espèce que si son mécanisme — apparaiement, séparation des homologues puis des chromatides — s'exécute avec une fidélité extrême. La fréquence des non-disjonctions augmente avec l'âge maternel, ce qui justifie le dépistage prénatal proposé pendant la grossesse dans les maternités bien équipées.
\end{savaistu}

\courssub{Bilan de la leçon et méthode pour le Bac}

\begin{propriete}[Conclusion générale de la leçon]
La méiose et la fécondation sont deux phénomènes \cle{complémentaires} et \cle{indissociables} de la reproduction sexuée :
\begin{itemize}
\item la méiose \textbf{réduit} le stock chromosomique de moitié ($2n \rightarrow n$) et \textbf{brasse} l'information génétique (brassages inter- et intrachromosomiques) ;
\item la fécondation \textbf{restaure} la diploïdie ($n + n \rightarrow 2n$) et \textbf{mélange} les patrimoines génétiques de deux parents.
\end{itemize}
Ensemble, ils assurent la \cle{pérennité de l'espèce} en conciliant deux exigences apparemment contradictoires : la \textbf{stabilité} du caryotype et de l'information génétique (rôle conservateur) et la \textbf{diversité} des individus (rôle créateur), matière première de l'adaptation et de l'évolution.
\end{propriete}

\begin{methode}[Structurer une question de synthèse au Bac]
Face à une question du type « Montrer que la méiose et la fécondation sont indispensables à la pérennité de l'espèce », organise ta réponse en deux axes distincts :
\begin{enumerate}
\item \textbf{Axe stabilité (conservation du caryotype)} : la méiose réduit le stock chromosomique ($2n \rightarrow n$), la fécondation le restaure ($n + n \rightarrow 2n$) ; les deux se compensent exactement, donc le caryotype de l'espèce est maintenu constant de génération en génération. Appuie-toi sur le raisonnement par l'absurde (doublement sans méiose, stérilité sans fécondation) et sur la courbe d'ADN (retour au $2q$ initial).
\item \textbf{Axe diversité (moteur de l'évolution)} : les brassages méiotiques (interchromosomique à l'anaphase I, intrachromosomique par crossing-over en prophase I) produisent des gamètes tous différents ; la fécondation en combine deux au hasard ; il en résulte des individus génétiquement uniques, gage d'adaptation de l'espèce aux changements du milieu.
\end{enumerate}
Conclus toujours par la formule-bilan : \emph{stabilité + diversité = pérennité de l'espèce}.
\end{methode}

\begin{exoresolu}[Exercice type Bac : analyse de documents]
\textbf{Énoncé.} Chez une espèce animale, on mesure la quantité d'ADN par noyau dans trois types de cellules : cellules de la peau ($\SI{6,6}{pg}$), spermatozoïdes ($\SI{3,3}{pg}$), cellules germinales en fin de phase $S$ ($\SI{13,2}{pg}$). (1) Attribuer à chaque valeur le stade correspondant du cycle chromosomique. (2) Expliquer, à l'aide de deux phénomènes biologiques précis, comment cette espèce maintient constante la quantité d'ADN de ses cellules somatiques au fil des générations, tout en produisant des individus génétiquement différents.
\tcblower
\textbf{Solution.} (1) Les cellules de la peau sont des cellules somatiques diploïdes en $G_1$ : $2q = \SI{6,6}{pg}$. Les spermatozoïdes sont des gamètes haploïdes : $q = \SI{3,3}{pg}$, soit exactement la moitié. Les cellules germinales en fin de phase $S$ ont répliqué leur ADN : $4q = \SI{13,2}{pg}$, soit le double des cellules somatiques. Les rapports $1:2:4$ ($q:2q:4q$) confirment le cycle $q \rightarrow 2q \rightarrow 4q \rightarrow q$.
(2) \emph{Stabilité} : la méiose ramène la quantité d'ADN de $4q$ à $q$ (deux divisions successives sans réplication intermédiaire après une seule phase $S$), produisant des gamètes haploïdes ; la fécondation réunit deux gamètes ($q + q = 2q$), restaurant la diploïdie. Les deux phénomènes se compensent exactement : chaque génération repart de $2q = \SI{6,6}{pg}$ dans ses cellules somatiques. \emph{Diversité} : les brassages interchromosomique (ségrégation aléatoire des paires d'homologues en anaphase I) et intrachromosomique (crossing-over en prophase I) rendent chaque gamète génétiquement unique, et la rencontre des gamètes est aléatoire ; chaque zygote est donc une combinaison inédite. Conclusion : la complémentarité méiose-fécondation assure à la fois la constance du stock génétique et la diversité des individus, conditions de la pérennité de l'espèce.
\end{exoresolu}

\begin{exercice}[Pour t'entraîner]
Chez le maïs ($2n = 20$), calcule le nombre de types de gamètes qu'une plante peut produire par brassage interchromosomique seul, puis le nombre de zygotes distincts issus du croisement de deux plantes (sans crossing-over). Rédige ensuite un paragraphe de dix lignes maximum montrant que méiose et fécondation sont « deux phénomènes complémentaires au service de la pérennité de l'espèce ».
\end{exercice}

\begin{corrige}[Exercice : maïs et pérennité]
Nombre de paires de chromosomes homologues : $n = 10$. Nombre de gamètes distincts par plante par brassage interchromosomique : $2^{10} = \num{1024}$. Nombre de zygotes distincts possibles par croisement de deux plantes : $2^{10} \times 2^{10} = 2^{20} = \num{1048576}$, soit plus d'un million de combinaisons (sans crossing-over). 

Paragraphe attendu : La méiose réduit le stock chromosomique de $2n$ à $n$ et brasse les allèles (brassages inter- et intrachromosomiques) ; la fécondation restaure $2n$ en unissant deux gamètes au hasard. Sans méiose, le stock doublerait à chaque génération ; sans fécondation, les individus haploïdes seraient non viables et stériles. Les deux phénomènes se compensent exactement (stabilité du caryotype maintenue) tout en générant des individus génétiquement uniques (diversité), matière première de l'adaptation : c'est ainsi qu'ils assurent ensemble la pérennité de l'espèce.
\end{corrige}

\newpage

\coursec{Activité d'intégration}

\courssub{Objectifs de l'activité}

À l'issue de cette activité d'intégration, tu seras capable de :
\begin{itemize}
\item interpréter une courbe d'évolution de la quantité d'\cle{ADN} dans les cellules de la lignée germinale et identifier les populations cellulaires correspondantes ;
\item reconnaître, ordonner et nommer les étapes de la \cle{méiose} à partir de documents microscopiques ;
\item exploiter un caryotype pour vérifier la stabilité du nombre de \cle{chromosomes} d'une génération à l'autre ;
\item calculer le nombre de gamètes génétiquement différents qu'un individu peut produire ;
\item argumenter, par un texte structuré, la \cle{complémentarité} entre méiose et fécondation dans le maintien du caryotype et le renouvellement de la diversité génétique, conditions de la \cle{pérennité de l'espèce}.
\end{itemize}

\courssub{Situation}

\textbf{Le mystère du caryotype constant : enquête sur deux générations de zébus de l'Adamaoua.}

\medskip
À Ngaoundéré, dans la région de l'Adamaoua, un éleveur de zébus Goudali croise un mâle reproducteur avec une femelle de race améliorée (croisée Montbéliarde) dans le but d'augmenter la production laitière de son troupeau. Les deux parents possèdent le même nombre de chromosomes : $2n = 60$ (soit $n = 30$ paires de chromosomes homologues). Sur trois ans, le troupeau s'enrichit de 12 veaux. L'éleveur fait deux constats qui l'intriguent :

\begin{itemize}
\item \textbf{Constat 1 :} un laboratoire vétérinaire de Ngaoundéré a réalisé le caryotype de trois veaux : tous possèdent exactement 60 chromosomes, comme leurs parents. Le nombre chromosomique est \emph{stable} d'une génération à l'autre.
\item \textbf{Constat 2 :} pourtant, aucun veau ne ressemble exactement à un autre ni à ses parents : couleur de la robe, taille de la bosse, forme des cornes, résistance aux parasites varient d'un individu à l'autre. La descendance est \emph{diverse}.
\end{itemize}

Comment la reproduction sexuée peut-elle à la fois \textbf{conserver} le caryotype de l'espèce et \textbf{brasser} les caractères héréditaires ? Pour répondre, tu disposes de quatre documents.

\medskip
\textbf{Document 1 — Les caryotypes.} Les caryotypes du père, de la mère et d'un veau montrent chacun 30 paires de chromosomes homologues, soit $2n = 60$. Chez le veau, chaque paire est constituée d'un chromosome d'origine paternelle et d'un chromosome d'origine maternelle.

\medskip
\textbf{Document 2 — Courbe de la quantité d'ADN dans les cellules testiculaires du mâle.} On mesure la quantité d'ADN par noyau (en unités arbitraires, u.a.) dans un échantillon de cellules du testicule. On distingue quatre populations de cellules :

\begin{center}
\begin{tabularx}{\linewidth}{|c|>{\centering\arraybackslash}X|>{\centering\arraybackslash}X|>{\centering\arraybackslash}X|>{\centering\arraybackslash}X|}
\hline
\rowcolor{popblueL} Population & A & B & C & D \\
\hline
Quantité d'ADN par noyau (u.a.) & 2 & 4 & 2 & 1 \\
\hline
\end{tabularx}
\end{center}

On note $Q$ la quantité d'ADN d'un gamète (spermatozoïde) : $Q = 1$ u.a. Les populations A et C ont la même quantité d'ADN mais n'ont pas le même contenu chromosomique.

\medskip
\textbf{Document 3 — Microphotographies d'étapes de la méiose.} Quatre clichés de cellules en division, pris dans le désordre, montrent : (cliché 1) des chromosomes bivalents alignés sur la plaque équatoriale ; (cliché 2) des chromosomes à deux chromatides rassemblés en bivalents avec des figures d'enroulement (chiasmas) ; (cliché 3) des chromosomes à une chromatide migrant vers les pôles dans des cellules déjà haploïdes ; (cliché 4) des chromosomes à deux chromatides migrant vers les pôles, les centromères n'étant pas divisés.

\medskip
\textbf{Document 4 — Expérience de blocage de la méiose.} Chez un animal modèle, on administre une substance qui bloque la méiose : les gamètes produits restent diploïdes ($2n$). La fécondation entre deux gamètes $2n$ donne des œufs tétraploïdes ($4n = 120$) qui ne sont pas viables : le développement embryonnaire s'arrête précocement.

\courssub{Consignes}

\begin{enumerate}
\item \textbf{Exploitation de la courbe d'ADN.} En justifiant chaque réponse, identifie les quatre populations cellulaires A, B, C et D du document 2 (nature des cellules et stade de la gamétogenèse). Précise à quel moment du cycle cellulaire la quantité d'ADN passe de 2 à 4 u.a.
\item \textbf{Étapes de la méiose.} Remets les quatre clichés du document 3 dans l'ordre chronologique, nomme chaque étape (division et phase) et décris le comportement des chromosomes à chacune d'elles.
\item \textbf{Analyse des caryotypes.} À l'aide du document 1, explique l'origine des 60 chromosomes du veau et justifie la formule : $n + n = 2n$.
\item \textbf{Calcul de la diversité.} Sachant que le zébu possède $n = 30$ paires de chromosomes homologues et qu'à l'anaphase I la répartition des chromosomes paternels et maternels est indépendante et aléatoire, calcule le nombre de types de gamètes différents qu'un parent peut produire (sans tenir compte du crossing-over). Donne le résultat en notation scientifique. Que devient ce nombre si l'on considère les deux parents réunis ?
\item \textbf{Interprétation de l'expérience de blocage.} À partir du document 4, montre que la méiose est indispensable à la stabilité du caryotype au fil des générations.
\item \textbf{Synthèse argumentée.} Rédige un texte argumenté de 15 à 20 lignes répondant à la question de l'éleveur : \og Pourquoi tous mes veaux ont-ils 60 chromosomes tout en étant tous différents ? \fg{} Tu concluras sur la nécessité des deux phénomènes — méiose et fécondation — pour la pérennité de l'espèce.
\end{enumerate}

\courssub{Ressources à mobiliser}

\begin{aretenir}[Les outils de la leçon à réactiver]
\begin{itemize}
\item \textbf{Méiose :} deux divisions successives. La première (division réductionnelle) sépare les chromosomes homologues : passage de $2n$ chromosomes à deux chromatides à $n$ chromosomes à deux chromatides. La seconde (division équationnelle) sépare les chromatides : $n$ chromosomes à une chromatide.
\item \textbf{Quantité d'ADN :} $Q$ dans un gamète ; $2Q$ dans une cellule diploïde en $G_1$ ; $4Q$ après la réplication (phase $S$) et jusqu'à la fin de la méiose I ; $2Q$ dans les cellules haploïdes à chromosomes bichromatidiens (entre méiose I et méiose II) ; $Q$ après la méiose II.
\item \textbf{Fécondation :} fusion de deux gamètes haploïdes ($n$) qui restaure la diploïdie : $n + n = 2n$.
\item \textbf{Diversité génétique :} brassage interchromosomique (anaphase I) : $2^{n}$ combinaisons possibles par parent ; brassage intrachromosomique (crossing-over en prophase I) qui multiplie encore la diversité.
\item \textbf{Pérennité :} la méiose divise par deux le stock chromosomique, la fécondation le restaure : leur alternance maintient le caryotype de l'espèce tout en créant des individus génétiquement uniques.
\end{itemize}
\end{aretenir}

\courssub{Démarche et corrigé commenté}

\begin{exoresolu}[Le mystère du caryotype constant — corrigé intégral]
\textbf{Question 1.} Identifier les quatre populations cellulaires du testicule et situer le passage de 2 à 4 u.a. d'ADN.

\textbf{Question 2.} Ordonner et nommer les étapes de la méiose (clichés 1 à 4).

\textbf{Question 3.} Expliquer l'origine des 60 chromosomes du veau.

\textbf{Question 4.} Calculer le nombre de gamètes différents possibles chez le zébu ($n = 30$).

\textbf{Question 5.} Montrer, à l'aide du document 4, que la méiose est indispensable.

\textbf{Question 6.} Rédiger la synthèse argumentée.
\tcblower
\textbf{1. Identification des populations cellulaires.}

On raisonne sur deux critères : la \emph{quantité d'ADN} et le \emph{nombre de chromosomes} (ploidie).

\begin{itemize}
\item \textbf{Population D (1 u.a.) :} quantité $Q$. Ce sont les \textbf{spermatides} (ou spermatozoïdes après spermiogenèse) : cellules haploïdes à $n = 30$ chromosomes à une chromatide.
\item \textbf{Population B (4 u.a.) :} quantité $4Q$. Ce sont les \textbf{spermatocytes I} (et les spermatogonies en fin de phase $S$/$G_2$) : cellules diploïdes dont l'ADN a été répliqué, à $2n = 60$ chromosomes à deux chromatides.
\item \textbf{Population A (2 u.a.) :} quantité $2Q$ dans une cellule \emph{diploïde} : ce sont les \textbf{spermatogonies} en phase $G_1$ (avant réplication), à $2n = 60$ chromosomes à une chromatide.
\item \textbf{Population C (2 u.a.) :} quantité $2Q$ dans une cellule \emph{haploïde} : ce sont les \textbf{spermatocytes II}, issus de la méiose I, à $n = 30$ chromosomes à deux chromatides (d'où $2Q$ d'ADN malgré l'haploïdie).
\end{itemize}

Le passage de 2 à 4 u.a. correspond à la \textbf{réplication de l'ADN pendant la phase $S$} de l'interphase qui précède la méiose : chaque chromosome, d'abord à une chromatide, devient un chromosome à deux chromatides identiques.

\medskip
\textbf{2. Ordre chronologique des étapes.}

Ordre correct : \textbf{cliché 2 $\rightarrow$ cliché 1 $\rightarrow$ cliché 4 $\rightarrow$ cliché 3}.

\begin{itemize}
\item \textbf{Cliché 2 — Prophase I :} les chromosomes homologues s'apparient en \textbf{bivalents} ; les chiasmas témoignent du \textbf{crossing-over} (brassage intrachromosomique).
\item \textbf{Cliché 1 — Métaphase I :} les bivalents s'alignent sur la plaque équatoriale ; l'orientation de chaque paire est aléatoire (source du brassage interchromosomique).
\item \textbf{Cliché 4 — Anaphase I :} les chromosomes \emph{homologues} se séparent et migrent vers les pôles opposés, chacun restant formé de deux chromatides (les centromères ne se divisent pas) : c'est la \textbf{division réductionnelle}, $2n \rightarrow n$.
\item \textbf{Cliché 3 — Anaphase II :} dans des cellules déjà haploïdes, les centromères se divisent et les chromatides sœurs migrent vers les pôles : c'est la \textbf{division équationnelle}, qui aboutit à quatre cellules haploïdes à $n$ chromosomes à une chromatide.
\end{itemize}

\medskip
\textbf{3. Origine des 60 chromosomes du veau.}

Le père produit par méiose des spermatozoïdes haploïdes ($n = 30$) ; la mère produit des ovules haploïdes ($n = 30$). À la fécondation, les deux noyaux gamétiques fusionnent :
\[
n + n = 30 + 30 = 60 = 2n
\]
Dans chaque paire de chromosomes homologues du veau, un chromosome vient du père, l'autre de la mère : c'est exactement ce que montre le caryotype du document 1. La diploïdie est ainsi \emph{restaurée}, ni augmentée ni diminuée.

\medskip
\textbf{4. Calcul de la diversité gamétique.}

À l'anaphase I, pour chacune des 30 paires, le chromosome paternel peut migrer vers l'un ou l'autre pôle, indépendamment des autres paires. Le nombre de combinaisons est donc :
\[
2^{n} = 2^{30} = \num{1073741824} \approx \num{1,07e9} \text{ types de gamètes par parent.}
\]
Pour le couple, le nombre de combinaisons possibles pour un veau est :
\[
2^{30} \times 2^{30} = 2^{60} \approx \num{1,15e18}.
\]
Ce chiffre, déjà considérable, est encore \emph{sous-estimé} car il ne tient pas compte du crossing-over de la prophase I, qui crée des chromosomes recombinés uniques. Il est donc pratiquement impossible que deux veaux (hors vrais jumeaux) soient génétiquement identiques.

\medskip
\textbf{5. La méiose est indispensable à la stabilité du caryotype.}

Le document 4 fournit la preuve par l'absurde : si la méiose est bloquée, les gamètes restent diploïdes ($2n = 60$). Leur fusion donne un œuf tétraploïde :
\[
2n + 2n = 4n = 120 \text{ chromosomes, non viable.}
\]
À chaque génération, le nombre de chromosomes doublerait ($2n \rightarrow 4n \rightarrow 8n \ldots$), ce qui est incompatible avec la vie. La méiose, en divisant par deux le stock chromosomique, est donc le \emph{préalable obligé} à toute fécondation viable.

\medskip
\textbf{6. Synthèse argumentée (correction-type).}

\emph{Tous les veaux de l'éleveur possèdent 60 chromosomes parce que la reproduction sexuée alterne deux phénomènes complémentaires. La méiose, qui se déroule dans les gonades des parents, comporte deux divisions : la première sépare les chromosomes homologues et fait passer les cellules de $2n = 60$ à $n = 30$ chromosomes ; la seconde sépare les chromatides. Elle produit ainsi des gamètes haploïdes. Lors de la fécondation, un spermatozoïde ($n = 30$) fusionne avec un ovule ($n = 30$) : l'œuf redevient diploïde ($2n = 60$), comme le confirme le caryotype des veaux. Si la méiose n'avait pas lieu, les gamètes seraient diploïdes et donneraient des descendants tétraploïdes non viables (document 4) : la méiose garantit donc la stabilité du caryotype. Mais cette même méiose est aussi une source de diversité : le crossing-over de la prophase I recompose les chromosomes, et la répartition aléatoire des homologues à l'anaphase I permet à chaque parent de produire plus d'un milliard de types de gamètes ($2^{30} \approx \num{1,07e9}$). La fécondation, en réunissant au hasard deux de ces gamètes uniques, crée un individu génétiquement nouveau : voilà pourquoi les veaux sont tous différents. En conclusion, la méiose (réduction et brassage) et la fécondation (restauration de la diploïdie et rencontre de deux lots génétiques) sont deux phénomènes indissociables : leur complémentarité assure à la fois la stabilité du caryotype et le renouvellement de la diversité génétique, les deux conditions de la pérennité de l'espèce.}
\end{exoresolu}

\courssub{Bilan de l'activité}

Cette enquête sur le troupeau de zébus de l'Adamaoua t'a permis de mobiliser simultanément tous les savoir-faire de la leçon et de dégager trois idées-force :

\begin{itemize}
\item \textbf{La méiose réduit et brasse.} En deux divisions, elle fait passer les cellules germinales de $2n$ à $n$ chromosomes et la quantité d'ADN de $4Q$ à $Q$, tout en créant, par crossing-over et répartition aléatoire des homologues, une diversité gamétique immense ($2^{30} \approx \num{1,07e9}$ gamètes différents chez le zébu).
\item \textbf{La fécondation restaure et combine.} La fusion de deux gamètes haploïdes rétablit la diploïdie ($n + n = 2n = 60$) et associe deux lots génétiques uniques, ce qui explique à la fois la constance du caryotype et l'originalité de chaque veau.
\item \textbf{Les deux phénomènes sont indissociables.} Sans méiose, le nombre de chromosomes doublerait à chaque génération (descendants $4n$ non viables) ; sans fécondation, pas de restauration de la diploïdie ni de brassage entre individus. Leur alternance au cours du cycle de vie garantit la \cle{pérennité de l'espèce} : stabilité du patrimoine chromosomique et renouvellement permanent de la diversité génétique, condition de l'adaptation des populations à leur environnement.
\end{itemize}

\begin{savaistu}[Chez les Spermaphytes : la double fécondation]
Chez les plantes à fleurs (Angiospermes), la fécondation est \textbf{double} : un spermatozoïde féconde l'oosphère ($n$) pour former le zygote ($2n$) qui donnera l'embryon ; l'autre spermatozoïde fusionne avec les deux noyaux polaires ($2n$) pour former le noyau tertiaire ($3n$) qui donnera l'albumen (tissu nutritif). La méiose a lieu chez les anthères (microsporogenèse $\rightarrow$ grains de pollen) et dans l'ovule (macrosporogenèse $\rightarrow$ sac embryonnaire). La complémentarité méiose/fécondation assure ici aussi la pérennité de l'espèce.
\end{savaistu}

\begin{attention}[Pièges fréquents à éviter dans ce type d'activité]
\begin{itemize}
\item Ne confonds pas \emph{quantité d'ADN} et \emph{nombre de chromosomes} : un spermatocyte II contient $n$ chromosomes mais $2Q$ d'ADN, car chaque chromosome comporte encore deux chromatides.
\item La réduction du nombre de chromosomes a lieu en \textbf{anaphase I} (séparation des homologues), pas en méiose II : la seconde division est équationnelle.
\item Dans le calcul $2^{n}$, $n$ désigne le nombre de \emph{paires} de chromosomes homologues ($n = 30$), pas le nombre diploïde ($2n = 60$).
\item La fécondation ne \og crée \fg{} pas la diversité à elle seule : elle \emph{combine} des gamètes déjà diversifiés par la méiose.
\end{itemize}
\end{attention}

\newpage

\coursec{Méthodes et exercices résolus}

\courssub{Fiche méthode 1 : Réaliser la dissection d'un Mammifère et d'une fleur, et schématiser les organes reproducteurs}

\begin{methode}[Dissection et schéma d'observation des organes reproducteurs]
\begin{enumerate}
\item \textbf{Préparer le matériel} : animal fraîchement euthanasié (rat, cobaye) ou fleur fraîche (hibiscus, flamboyant), planche à disséquer, pinces fines, ciseaux, scalpel, aiguilles, loupe binoculaire, gants, blouse.
\item \textbf{Chez le Mammifère} : écarter la peau et la paroi abdominale, repérer les \cle{testicules} (souvent descendus dans les bourses scrotales chez l'adulte) ou les \cle{ovaires} (petits organes ovoïdes fixés près des reins par le mésovaire), puis suivre les voies génitales (épididyme, canal déférent, vésicules séminales, prostate ; trompes utérines, utérus, vagin).
\item \textbf{Chez la fleur} : écarter sépales (calice) et pétales (corolle), identifier l'\cle{androcée} (étamines = filet + anthère) et le \cle{gynécée} (pistil = stigmate, style, ovaire contenant les ovules). Ouvrir l'ovaire à l'aiguille pour observer les ovules (macrosporanges).
\item \textbf{Observer puis schématiser} : le schéma se fait au crayon à papier, traits nets et continus, \textbf{sans ombre ni couleur superflue}, avec un titre souligné précisant l'objet, l'espèce et la vue (ex. : « Coupe longitudinale de la fleur d'\emph{Hibiscus sabdariffa} »), et des légendes disposées de part et d'autre, reliées par des traits horizontaux non croisés.
\item \textbf{Respecter les proportions} et indiquer le grossissement (ex. : $\times 10$) ou l'échelle si l'observation est microscopique.
\end{enumerate}
\textbf{Réflexes} : une légende = un nom exact d'organe (vocabulaire anatomique précis) ; jamais de flèches croisées ; le titre est obligatoire et informatif.
\end{methode}

\begin{exoresolu}[Dissection d'une fleur de \emph{Hibiscus sabdariffa} (oseille de Guinée / Bissap)]
Au cours d'un TP au lycée de Ngaoundéré, un élève dissèque une fleur d'oseille de Guinée (\emph{Hibiscus sabdariffa}) et observe, de l'extérieur vers l'intérieur : 5 sépales verts soudés (calice gamosépalé), 5 pétales rouges (corolle dialypétale), un tube staminal portant de nombreuses anthères (androcée monadelphe), et au centre un pistil dont l'ovaire supère renferme plusieurs ovules.\\
1) Identifier les organes reproducteurs de cette fleur et préciser leur rôle.\\
2) Représenter schématiquement la coupe longitudinale de la fleur.
\tcblower
\textbf{1) Identification des organes reproducteurs.}\\
Les pièces florales stériles (calice, corolle) assurent la protection du bouton floral et l'attraction des pollinisateurs ; les organes reproducteurs (gonades) sont :
\begin{itemize}
\item l'\cle{androcée}, ensemble des étamines soudées par leurs filets en un \textbf{tube staminal} portant de nombreuses anthères : chaque anthère contient 4 sacs polliniques (microsporangites) où se déroule la \textbf{microsporogenèse} (méiose) puis la \textbf{microgamétogenèse} (mitoses) aboutissant à la formation des grains de pollen (gamétophytes mâles, $n$) ;
\item le \cle{gynécée} (pistil), formé du stigmate (surface réceptive, souvent gluante), du style et de l'ovaire supère renfermant les ovules (macrosporangites) : dans chaque ovule se déroule la \textbf{macrosporogenèse} (méiose) puis la \textbf{macrogamétogenèse} (mitoses) aboutissant à la formation du sac embryonnaire (gamétophyte femelle, $n$) contenant l'oosphère.
\end{itemize}
\textbf{2) Schéma de la coupe longitudinale.}
\begin{popfigure}
\begin{tikzpicture}[scale=0.95]
% Pedicel / Receptacle
\draw[thick,popgreen] (0,0) -- (0,0.8);
% Calice (sepal)
\draw[thick,popgreen] (-0.8,0.8) -- (-1.2,1.5);
\draw[thick,popgreen] (0.8,0.8) -- (1.2,1.5);
% Corolle (petals) - simplified
\draw[thick,poppink] (-1.2,1.5) .. controls (-2,2.5) and (-1.5,3.5) .. (0,3.8);
\draw[thick,poppink] (1.2,1.5) .. controls (2,2.5) and (1.5,3.5) .. (0,3.8);
% Staminal tube (Androecée monadelphe)
\draw[thick,poporange, fill=poporangeL!30] (-0.35,1.0) rectangle (0.35,3.0);
% Anthers on tube
\foreach \y in {1.3, 1.7, 2.1, 2.5, 2.9} {
    \fill[popgold] (-0.45, \y) circle (0.07);
    \fill[popgold] (0.45, \y) circle (0.07);
}
% Pistil (Gynécée)
\draw[thick,popblue, fill=popblueL!20] (-0.2,1.0) -- (-0.2,2.8) -- (0,3.2) -- (0.2,2.8) -- (0.2,1.0) -- cycle; % Ovary
\draw[thick,popblue] (0,2.8) -- (0,3.5); % Style
\fill[poppink] (0,3.6) circle (0.12); % Stigma
% Ovules inside ovary
\foreach \y in {1.3, 1.6, 1.9, 2.2, 2.5} {
    \fill[popgreen] (0, \y) circle (0.05);
}
% Legends
\draw[->,popmuted] (1.5,3.6) -- (0.5,3.5); \node[right,popmuted] at (1.5,3.6) {\small Stigmate};
\draw[->,popmuted] (1.5,3.2) -- (0.3,3.2); \node[right,popmuted] at (1.5,3.2) {\small Style};
\draw[->,popmuted] (1.5,2.0) -- (0.4,2.0); \node[right,popmuted] at (1.5,2.0) {\small Ovaire (supère)};
\draw[->,popmuted] (1.5,1.6) -- (0.3,1.6); \node[right,popmuted] at (1.5,1.6) {\small Ovules};
\draw[->,popmuted] (-1.5,3.0) -- (-0.5,3.0); \node[left,popmuted] at (-1.5,3.0) {\small Anthères sur tube staminal};
\draw[->,popmuted] (-1.5,1.2) -- (-0.5,1.2); \node[left,popmuted] at (-1.5,1.2) {\small Filets soudés (tube)};
\draw[->,popmuted] (-1.5,0.5) -- (-1.0,1.0); \node[left,popmuted] at (-1.5,0.5) {\small Sépale (calice)};
\draw[->,popmuted] (1.5,0.5) -- (1.0,1.0); \node[right,popmuted] at (1.5,0.5) {\small Pétale (corolle)};
\end{tikzpicture}
\legende{Coupe longitudinale schématique d'une fleur d'\emph{Hibiscus sabdariffa} (Malvacée) : androcée monadelphe (tube staminal + anthères) et gynécée (stigmate, style, ovaire supère à ovules nombreux).}
\end{popfigure}
\textbf{Conclusion} : les gonades végétales sont l'anthère (gonade mâle, site de la microsporogenèse) et l'ovaire (gonade femelle, site de la macrosporogenèse) ; elles abritent la méiose et la gamétogenèse, garantes de la reproduction sexuée.
\end{exoresolu}

\courssub{Fiche méthode 2 : Reconnaître les étapes de la méiose et décrire le comportement des chromosomes}

\begin{methode}[Identifier une étape de méiose sur une microphotographie ou un schéma]
\begin{enumerate}
\item \textbf{Compter les lots de chromosomes} visibles et repérer s'il y a une ou deux cellules : la méiose comporte \textbf{deux divisions successives} sans réplication d'ADN intercalaire (Méiose I réductionnelle, Méiose II équationnelle).
\item \textbf{Chercher les figures caractéristiques} :
\begin{itemize}
\item \textbf{Prophase I} : chromosomes bivalents (tétades) accolés par synapsis, chiasmas visibles (points d'enjambement) ; noyau volumineux.
\item \textbf{Métaphase I} : bivalents alignés sur la plaque équatoriale (plan médian), fuseau bipolaire.
\item \textbf{Anaphase I} : les chromosomes \emph{entiers} (chacun à 2 chromatides) migrent vers les pôles opposés — c'est l'étape \cle{réductionnelle} (séparation des homologues) ; $2n \to n$.
\item \textbf{Prophase II / Métaphase II} : deux cellules filles, chromosomes (à 2 chromatides) alignés sur la plaque équatoriale de chacune.
\item \textbf{Anaphase II} : les chromatides sœurs se séparent (division équationnelle, comme en mitose) ; $n$ chromosomes à 1 chromatide par pôle.
\item \textbf{Télophase II} : 4 cellules haploïdes à $n$ chromosomes simples (1 chromatide), décondensation, réapparition des enveloppes nucléaires.
\end{itemize}
\item \textbf{Compter les chromosomes} : passer de $2n$ à $n$ (nombre de centromères) indique la méiose I ; rester à $n$ (mais passer de 2 chromatides à 1 chromatide par chromosome) indique la méiose II.
\item \textbf{Rédiger la description} : nommer l'étape, décrire ce qu'on voit (alignement, migration, nombre, forme), puis conclure sur l'effet (réduction, brassage, conservation).
\end{enumerate}
\textbf{Réflexes} : « bivalent » = 2 chromosomes homologues = 4 chromatides (tétade) ; le brassage interchromosomique (indépendant) se joue en anaphase I ($2^{n}$ combinaisons possibles) ; le brassage intrachromosomique (enjambements) se matérialise par les chiasmas en prophase I.
\end{methode}

\begin{exoresolu}[Analyse de figures de méiose chez le criquet (\emph{Schistocerca gregaria})]
On observe au microscope des coupes de testicules de criquet ($2n = 24$ chromosomes). Sur quatre figures A, B, C, D, on note :\\
A : 12 bivalents accolés deux à deux, avec des points de contact en X ;\\
B : deux groupes de 12 chromosomes à deux chromatides migrant vers les pôles opposés d'une même cellule ;\\
C : 12 chromosomes à deux chromatides alignés sur la plaque équatoriale de chacune des deux cellules ;\\
D : 4 cellules contenant chacune 12 chromosomes à une chromatide.\\
1) Identifier chaque étape et la justifier.\\
2) Calculer le nombre de combinaisons chromosomiques possibles dans les gamètes par brassage interchromosomique.
\tcblower
\textbf{1) Identification des étapes.}
\begin{itemize}
\item \textbf{A = Prophase I (stade pachytène/diplotène)} : les chromosomes homologues s'apparient en \cle{bivalents} (12 bivalents = 24 chromosomes = $2n$) ; les points en X sont les \cle{chiasmas}, témoins visibles des enjambements (brassage intrachromosomique).
\item \textbf{B = Anaphase I} : les homologues de chaque bivalent se séparent et migrent vers les pôles opposés, chacun encore formé de 2 chromatides sœurs ; c'est la division \textbf{réductionnelle} : chaque pôle ne recevra que $n = 12$ chromosomes.
\item \textbf{C = Métaphase II} : les deux cellules issues de la méiose I alignent leurs $n = 12$ chromosomes (chacun à 2 chromatides) sur la plaque équatoriale.
\item \textbf{D = Télophase II (fin de méiose)} : 4 cellules haploïdes ($n = 12$) à chromosomes simples (1 chromatide), futures spermatides.
\end{itemize}
\textbf{2) Nombre de combinaisons gamétiques.}\\
Le nombre de combinaisons par brassage interchromosomique (ségrégation indépendante des paires homologues) est $2^{n}$, avec $n = 12$ (nombre de paires) :
\[ 2^{n} = 2^{12} = 4\,096 \text{ combinaisons chromosomiques possibles par individu.} \]
\textbf{Conclusion} : la méiose réduit la ploïdie de $2n$ à $n$ (passage en anaphase I) et génère une diversité considérable (au moins 4096 types de spermatozoïdes chez ce criquet, sans compter les enjambements), ce qui assure la variabilité génétique indispensable à la pérennité de l'espèce.
\end{exoresolu}

\courssub{Fiche méthode 3 : Interpréter la courbe d'évolution de la quantité d'ADN au cours de la méiose et de la gamétogenèse}

\begin{methode}[Lire une courbe de quantité d'ADN par cellule (pg/cellule)]
\begin{enumerate}
\item \textbf{Repérer les paliers} : noter $Q$ la quantité d'ADN d'une cellule haploïde ($n$, chromosomes à 1 chromatide). Une cellule $2n$ non répliquée (G1) contient $2Q$ ; après réplication (phase S, G2), elle contient $4Q$.
\item \textbf{Identifier les transitions majeures} :
\begin{itemize}
\item $2Q \rightarrow 4Q$ : \textbf{réplication de l'ADN} (phase S, interphase pré-méiotique) ;
\item $4Q \rightarrow 2Q$ : \textbf{Méiose I} (séparation des chromosomes homologues, anaphase I) ;
\item $2Q \rightarrow Q$ : \textbf{Méiose II} (séparation des chromatides sœurs, anaphase II) ;
\item $Q + Q \rightarrow 2Q$ : \textbf{Fécondation} (fusion des noyaux gamétiques, caryogamie, retour à la diploïdie).
\end{itemize}
\item \textbf{Associer chaque palier à un type cellulaire} (ex. spermatogenèse) : spermatogonie G1 ($2Q$), spermatocyte I ($4Q$), spermatocyte II ($2Q$), spermatide / spermatozoïde ($Q$).
\item \textbf{Vérifier la cohérence} : chaque division (méiose I, méiose II) divise la quantité d'ADN par deux ; toute rupture de pente correspond à une division ou à la fécondation.
\end{enumerate}
\textbf{Réflexe crucial} : quantité d'ADN $\neq$ nombre de chromosomes. Exemple : un spermatocyte II a $n$ chromosomes mais $2Q$ d'ADN (chromosomes à 2 chromatides) ; un spermatozoïde a $n$ chromosomes et $Q$ d'ADN (chromosomes à 1 chromatide).
\end{methode}

\begin{exoresolu}[Courbe de l'ADN au cours de la spermatogenèse chez un Mammifère]
On mesure la quantité d'ADN par noyau au cours de la spermatogenèse ; on note $Q = \SI{3,2}{pg}$ la quantité d'ADN d'un spermatozoïde.\\
1) Donner la quantité d'ADN d'une spermatogonie en G1, d'un spermatocyte I, d'un spermatocyte II.\\
2) Tracer l'allure de la courbe d'évolution de la quantité d'ADN par cellule et justifier chaque palier.
\tcblower
\textbf{1) Quantités d'ADN.}
\begin{itemize}
\item Spermatogonie en G1 : $2n$ chromosomes simples, soit $2Q = 2 \times 3{,}2 = \SI{6,4}{pg}$ ;
\item Spermatocyte I (après phase S, en G2/prophase I) : $2n$ chromosomes à 2 chromatides, soit $4Q = \SI{12,8}{pg}$ ;
\item Spermatocyte II (après méiose I) : $n$ chromosomes à 2 chromatides, soit $2Q = \SI{6,4}{pg}$ ;
\item Spermatide et spermatozoïde : $n$ chromosomes simples, soit $Q = \SI{3,2}{pg}$.
\end{itemize}
\textbf{2) Courbe.}
\begin{popfigure}
\begin{tikzpicture}
\begin{axis}[width=0.9\linewidth,height=6cm,
xlabel={Étapes successives de la spermatogenèse}, ylabel={Quantité d'ADN par noyau (pg)},
xmin=0,xmax=6,ymin=0,ymax=14,
xtick={0.5,1.5,2.5,3.5,4.5,5.5},
xticklabels={Spermatogonie (G1), Phase S, Spermatocyte I, Spermatocyte II, Spermatide, Spermatozoïde},
xticklabel style={font=\scriptsize,rotate=25,anchor=east},
ytick={3.2,6.4,12.8},
yticklabels={$Q$,$2Q$,$4Q$},
grid=both,grid style={popline!40},
ymajorgrids=true, xmajorgrids=false]
\addplot[popcurve,mark=*,mark size=1.5pt,very thick] coordinates {
(0,6.4) (1,6.4) 
(1.5,12.8) (2,12.8) (2.5,12.8) 
(3,6.4) (3.5,6.4) (4,6.4) 
(4.5,3.2) (5,3.2) (5.5,3.2)
};
\end{axis}
\end{tikzpicture}
\legende{Évolution de la quantité d'ADN par noyau au cours de la spermatogenèse : plateau $2Q$ (G1), montée $2Q \to 4Q$ (réplication phase S), plateau $4Q$ (Spermatocyte I), chute $4Q \to 2Q$ (Méiose I), plateau $2Q$ (Spermatocyte II), chute $2Q \to Q$ (Méiose II), plateau $Q$ (Spermatides/Spermatozoïdes).}
\end{popfigure}
\textbf{Justification} : la montée $2Q \to 4Q$ traduit la réplication semi-conservative de l'ADN en phase S ; le plateau $4Q$ correspond à la croissance et la prophase I ; la chute $4Q \to 2Q$ correspond à la méiose I (séparation des homologues) ; le plateau $2Q$ correspond à l'interphase brève (pas de réplication) et prophase II ; la chute $2Q \to Q$ correspond à la méiose II (séparation des chromatides sœurs).\\
\textbf{Conclusion} : la gamétogenèse produit des cellules à $Q$ d'ADN (haploïdes) ; la fécondation ($Q_{\text{mâle}} + Q_{\text{femelle}} = 2Q = \SI{6,4}{pg}$) restaurera la quantité caractéristique de l'espèce diploïde.
\end{exoresolu}

\courssub{Fiche méthode 4 : Identifier les cellules de la lignée germinale chez les Mammifères à leurs différents stades}

\begin{methode}[Reconnaître les cellules de la gamétogenèse mammalienne (coupes histologiques ou schémas)]
\begin{enumerate}
\item \textbf{Connaître les quatre phases successives} : 
\begin{enumerate}
\item \textbf{Multiplication} (mitoses) : gonies ($2n$, $2Q$) ;
\item \textbf{Croissance} (phase G2 prolongée) : cytes I ($2n$, $4Q$) ;
\item \textbf{Maturation} (méiose) : cytes I $\to$ cytes II ($n$, $2Q$) $\to$ tides ($n$, $Q$) ;
\item \textbf{Différenciation} (spermiogenèse ou maturation ovocytaire) : gamètes fonctionnels ($n$, $Q$).
\end{enumerate}
\item \textbf{Utiliser les critères d'identification microscopique} :
\begin{center}
\begin{tabularx}{\linewidth}{|l|X|X|X|}\hline
\rowcolor{popblueL}
\textbf{Cellule} & \textbf{Ploïdie / ADN} & \textbf{Morphologie / Position} & \textbf{Effectif issu d'une cellule mère} \\ \hline
Spermatogonie / Oogonie & $2n$ / $2Q$ & Petite, noyau dense, cytoplasme basophile ; \textbf{périphérie} du tube séminifère / cortex ovarien & -- \\ \hline
Spermatocyte I / Ovocyte I & $2n$ / $4Q$ & \textbf{Grosse} cellule, noyau clair (chromosomes décondensés ou figures de méiose I) ; lumière du tube / follicule ovarien & 1 $\to$ 1 (entrées en méiose) \\ \hline
Spermatocyte II & $n$ / $2Q$ & Cellule \textbf{plus petite}, noyau plus dense, figures de méiose II ; vers la lumière & 1 cyt I $\to$ 2 cyt II \\ \hline
Spermatide & $n$ / $Q$ & Cellule ronde, noyau dense, acrosome naissant ; \textbf{près de la lumière} (côté Sertoli) & 1 cyt II $\to$ 2 spermatides \\ \hline
Spermatozoïde & $n$ / $Q$ & Cellule \textbf{flagellée} polarisée : tête (noyau + acrosome), pièce intermédiaire (mitochondries), flagelle ; \textbf{dans la lumière} & 1 spermatide $\to$ 1 spermatozoïde \\ \hline
Ovocyte II (Ovule) & $n$ / $Q$ & \textbf{Très grosse} cellule, abondantes réserves (vitellus), zona pellucida ; arrêté en \textbf{Métaphase II} & 1 ovocyte I $\to$ 1 ovocyte II + 1 GP1 \\ \hline
Globules polaires (GP1, GP2) & $n$ / $Q$ & Petites cellules, peu de cytoplasme, dégénèrent rapidement & 1 ovocyte I $\to$ 1 GP1 ; 1 ovocyte II $\to$ 1 GP2 \\ \hline
\end{tabularx}
\end{center}
\item \textbf{Vérifier les effectifs finaux} : 1 spermatocyte I $\to$ 4 spermatozoïdes fonctionnels ; 1 ovocyte I $\to$ 1 ovule fonctionnel + 3 globules polaires (dégénérés).
\end{enumerate}
\end{methode}

\begin{exoresolu}[Effectifs et asymétrie de la gamétogenèse mammalienne]
Chez un Mammifère mâle, on estime qu'un tube séminifère contient \num{100000} spermatocytes I entrant en méiose. Chez la femelle, \num{100000} ovocytes I entrent également en maturation (ovulation successive).\\
1) Calculer le nombre de gamètes fonctionnels obtenus dans chaque sexe.\\
2) Expliquer la différence observée au niveau cellulaire et adaptatif.
\tcblower
\textbf{1) Effectifs des gamètes fonctionnels.}
\begin{itemize}
\item \textbf{Spermatogenèse} : chaque spermatocyte I donne, par les deux divisions successives de la méiose, 4 spermatides qui se différencient toutes en 4 spermatozoïdes fonctionnels (cytogenèse égale) :
\[ N_{\text{spermatozoïdes}} = 100\,000 \times 4 = 400\,000 \text{ spermatozoïdes.} \]
\item \textbf{Ovogenèse} : chaque ovocyte I donne, par des divisions cytoplasmiques \textbf{inégales} (cytokinèse asymétrique), \textbf{un seul} ovocyte II (futur ovule) fonctionnel, volumineux, et 3 globules polaires (GP1 issu de méiose I, GP2 issu de méiose II, et parfois division du GP1) qui dégénèrent par apoptose :
\[ N_{\text{ovules}} = 100\,000 \times 1 = 100\,000 \text{ ovules.} \]
\end{itemize}
\textbf{2) Explication cellulaire et adaptative.}\\
\begin{itemize}
\item \textbf{Cellulaire} : lors de l'ovogenèse, le fuseau méiotique est déporté vers la périphérie ; la quasi-totalité du cytoplasme (organites, ARNm maternels, réserves nutritives : vitellus, cortical granules) est conservée dans une seule cellule (ovocyte II), assurant l'autonomie métabolique de l'embryon avant implantation. Les globules polaires ne reçoivent que l'excédent chromosomique.
\item \textbf{Adaptatif} : le mâle produit une multitude de gamètes mobiles, peu coûteux, compensant la faible probabilité de rencontre (concurrence spermatique) ; la femelle investit massivement dans un seul gamète sessile, riche en réserves, garantissant le démarrage du développement.
\end{itemize}
\textbf{Conclusion} : la méiose produit 4 gamètes mâles fonctionnels par cellule-mère contre 1 seul gamète femelle, différence liée aux rôles respectifs des gamètes (mobilité et nombre contre réserves et protection).
\end{exoresolu}

\courssub{Fiche méthode 7 : La gamétogenèse chez les Spermaphytes (Angiospermes) : micro- et macro-gamétogenèse}

\begin{methode}[Décrire la formation des gamètes chez les Angiospermes]
\begin{enumerate}
\item \textbf{Microgamétogenèse (dans l'anthère = gonade mâle)} :
\begin{enumerate}
\item \textbf{Microsporogenèse} : une cellule mère des microspores ($2n$) $\xrightarrow{\text{Méiose}}$ 4 microspores ($n$) tétragones.
\item \textbf{Microgamétogenèse} : chaque microspore $\xrightarrow{\text{Mitose 1}}$ grain de pollen à 2 cellules (1 noyau végétatif / tubulaire + 1 noyau génératif / reproducteur) $\xrightarrow{\text{Mitose 2 (souvent dans le tube pollinique)}}$ \textbf{grain de pollen mature à 3 cellules} ($n$) : 1 noyau tubulaire + 2 noyaux spermatiques (gamètes mâles). \textit{Paroi : exine (sporopollénine) + intine}.
\end{enumerate}
\item \textbf{Macrogamétogenèse (dans l'ovule = gonade femelle, à l'intérieur de l'ovaire)} :
\begin{enumerate}
\item \textbf{Macrosporogenèse} : une cellule mère des macrospores ($2n$) $\xrightarrow{\text{Méiose}}$ 4 macrospores ($n$) alignées $\to$ \textbf{3 dégénèrent}, 1 seule \textbf{macrospore fonctionnelle} (chalazale).
\item \textbf{Macrogamétogenèse} : la macrospore fonctionnelle $\xrightarrow{\text{3 mitoses successives}}$ \textbf{sac embryonnaire (gamétophyte femelle) à 7 cellules / 8 noyaux} (type \emph{Polygonum}, le plus fréquent) :
\begin{itemize}
\item Pôle micropylaire : 1 \textbf{oosphère} ($n$) + 2 \textbf{synergides} ($n$) (guidage tube pollinique) ;
\item Centre : 1 \textbf{cellule centrale} binucléée ($2n$ = 2 noyaux polaires $n+n$) ;
\item Pôle chalazal : 3 \textbf{antipodes} ($n$) (souvent éphémères).
\end{itemize}
\end{enumerate}
\item \textbf{Schéma récapitulatif} : savoir schématiser le grain de pollen mature (3 noyaux) et le sac embryonnaire mature (7 cellules / 8 noyaux) avec légendes précises.
\end{enumerate}
\textbf{Réflexe} : chez les Angiospermes, le gamétophyte mâle (pollen) est très réduit (3 cellules) et le gamétophyte femelle (sac embryonnaire) est aussi réduit (7 cellules) et dépendant du sporophyte (nutrition) — adaptation à la vie terrestre.
\end{methode}

\begin{exoresolu}[Analyse de la micro- et macro-gamétogenèse chez le Maïs (\emph{Zea mays})]
Chez le maïs ($2n = 20$), on observe une anthère et un ovule au microscope.\\
1) Préciser le nombre de chromosomes et de chromatides au stade : microspore mère, microspore, grain de pollen mature (3 noyaux), macrospore mère, macrospore fonctionnelle, oosphère, noyau polaire fusionné.\\
2) Expliquer pourquoi on observe 4 microspores fonctionnelles mais 1 seule macrospore fonctionnelle.
\tcblower
\textbf{1) Bilan chromosomique ($2n=20 \Rightarrow n=10$).}
\begin{center}
\begin{tabular}{|c|c|c|c|}\hline
\rowcolor{popblueL}
\textbf{Stade} & \textbf{Ploïdie} & \textbf{Nb Chromosomes} & \textbf{Nb Chromatides / Chr} \\ \hline
Microspore mère (MMC) & $2n$ & 20 & 2 (après phase S) \\ \hline
Microspore (tétade) & $n$ & 10 & 1 \\ \hline
Noyau végétatif (pollen) & $n$ & 10 & 1 \\ \hline
Noyaux spermatiques (pollen) & $n$ & 10 & 1 \\ \hline
Macrospore mère (MMC) & $2n$ & 20 & 2 (après phase S) \\ \hline
Macrospore fonctionnelle & $n$ & 10 & 1 \\ \hline
Oosphère (dans sac embryonnaire) & $n$ & 10 & 1 \\ \hline
Noyaux polaires fusionnés (cellule centrale) & $2n$ & 20 & 1 (chacun) \\ \hline
\end{tabular}
\end{center}
\textbf{2) Asymétrie de la cytokinèse.}\\
\begin{itemize}
\item \textbf{Microsporogenèse} : les 4 produits de la méiose (tétade) reçoivent un cytoplasme équivalent $\to$ 4 microspores viables $\to$ 4 grains de pollen. Cela maximise la dispersion (anémophilie pour le maïs).
\item \textbf{Macrosporogenèse} : la cytokinèse est \textbf{très inégale} (comme en ovogenèse animale) : 3 macrospores reçoivent très peu de cytoplasme et dégénèrent ; seule la macrospore chalazale (la plus éloignée du micropyle) conserve les réserves et devient fonctionnelle. Cela concentre les ressources dans un seul gamétophyte femelle protégé à l'intérieur de l'ovule.
\end{itemize}
\textbf{Conclusion} : la gamétogenèse végétale suit le même schéma méiotique que l'animal (réduction $2n \to n$) mais la différenciation des gamétophytes (mitoses post-méiotiques) est spécifique : le pollen est un organisme mâle mobile (tube pollinique) et le sac embryonnaire est une structure femelle pluricellulaire nourrie par le sporophyte maternel.
\end{exoresolu}

\courssub{Fiche méthode 5 : Décrire et comparer la fécondation chez les Mammifères et la double fécondation chez les Spermaphytes}

\begin{methode}[Décrire la fécondation et interpréter la courbe d'ADN associée]
\begin{enumerate}
\item \textbf{Chez les Mammifères (fécondation interne, dans la trompe utérine)} : ordonner les étapes chronologiques :
\begin{enumerate}
\item Migration des spermatozoïdes (capacitation dans le tractus femelle) $\to$ rencontre dans l'ampoule de la trompe (tiers supérieur).
\item \textbf{Réaction acrosomique} : libération d'enzymes (hyaluronidase, acrosine) dissolvant la corona radiata et la \textbf{zone pellucide} (glycoprotéines ZP3, récepteurs spécifiques).
\item Pénétration d'\textbf{un seul spermatozoïde} (blocage de la polyspermie : réaction corticale $\to$ durcissement zone pellucide).
\item \textbf{Achèvement de la méiose II} de l'ovocyte II (arrêté en Métaphase II) $\to$ formation de l'ovule ($n$, $Q$) + 2ème globule polaire.
\item \textbf{Fusion des noyaux} (caryogamie) : noyau spermatique ($n$) + noyau femelle ($n$) $\to$ noyau zygotique \textbf{diploïde $2n$ ($2Q$)}.
\item Formation du \textbf{zygote (œuf)} $\to$ premières mitoses (clivages) $\to$ blastocyste $\to$ implantation (nidation) dans l'endomètre utérin vers J7.
\end{enumerate}
\item \textbf{Chez les Spermaphytes Angiospermes (double fécondation)} : suivre les deux noyaux spermatiques :
\begin{enumerate}
\item Germination du grain de pollen sur le stigmate (réhydratation, reconnaissance pistil/pollen) $\to$ croissance du \textbf{tube pollinique} dans le style (guidage chimio-tropique par les synergides).
\item Arrivée au micropyle $\to$ libération des \textbf{deux noyaux spermatiques} dans le sac embryonnaire (après dégénérescence d'une synergide).
\item \textbf{Première fécondation} : 1er noyau spermatique ($n$) + \textbf{oosphère} ($n$) $\to$ \textbf{Zygote} \textbf{$2n$} (futur \textbf{embryon}).
\item \textbf{Seconde fécondation} : 2ème noyau spermatique ($n$) + \textbf{cellule centrale} (2 noyaux polaires fusionnés = $2n$) $\to$ cellule \textbf{triploïde $3n$} (futur \textbf{albumen} / endosperme, tissu de réserve).
\end{enumerate}
C'est la \cle{double fécondation}, synapomorphie des Angiospermes.
\item \textbf{Sur une courbe d'ADN (quantité par noyau)} :
\begin{itemize}
\item Mammifères : saut $Q \to 2Q$ au moment de la caryogamie.
\item Angiospermes : apparition simultanée de deux noyaux : zygote ($2Q$) et albumen ($3Q$).
\end{itemize}
\item \textbf{Comparer} : dans les deux cas, la fécondation \textbf{restaure la diploïdie de l'espèce} ($2n$) et \textbf{associe deux patrimoines génétiques différents} (brassage à la fécondation = hasard de la rencontre). Différence : double fécondation + albumen triploïde propre aux Angiospermes ; fécondation interne vs externe (pollinisation).
\end{enumerate}
\textbf{Réflexe « Piège Bac »} : chez les Angiospermes, l'albumen (endosperme) est \textbf{triploïde ($3n$)}, jamais diploïde. Seul le zygote est diploïde ($2n$).
\end{methode}

\begin{exoresolu}[Double fécondation chez le Maïs (\emph{Zea mays}, culture majeure au Cameroun)]
Chez le maïs, $2n = 20$ chromosomes. Un grain de pollen germe sur la soie (stigmate) et son tube pollinique libère deux noyaux spermatiques dans le sac embryonnaire.\\
1) Donner le nombre de chromosomes de l'oosphère, des deux noyaux polaires fusionnés, du zygote et de l'albumen.\\
2) Expliquer pourquoi on parle de double fécondation et préciser le devenir de chaque produit.
\tcblower
\textbf{1) Nombres de chromosomes ($n = 10$).}
\begin{itemize}
\item Oosphère (gamète femelle) : $n = 10$ chromosomes (1 chromatide) ;
\item Noyaux polaires fusionnés (cellule centrale) : $n + n = 2n = 20$ chromosomes ;
\item Zygote (issu 1ère fécondation) : $n_{\text{mâle}} + n_{\text{femelle}} = 2n = 20$ chromosomes ;
\item Albumen / Endosperme (issu 2ème fécondation) : $n_{\text{mâle}} + 2n_{\text{central}} = 3n = 30$ chromosomes.
\end{itemize}
\textbf{2) Double fécondation et devenir des produits.}\\
On parle de \textbf{double fécondation} car les \emph{deux} noyaux spermatiques du grain de pollen réalisent chacun une fusion nucléaire distincte :
\begin{itemize}
\item le premier noyau spermatique ($n = 10$) fusionne avec l'oosphère ($n = 10$) : c'est la \textbf{fécondation proprement dite} (syngamie), qui donne le \cle{zygote} ($2n = 20$), à l'origine de l'\textbf{embryon} de la future plantule (radicule, plumule, scutellum, coléoptile) ;
\item le second noyau spermatique ($n = 10$) fusionne avec la cellule centrale diploïde ($2n = 20$) : il en résulte une cellule \textbf{triploïde} ($3n = 30$) dont les divisions mitotiques rapides donnent l'\cle{albumen} (endosperme), tissu de réserve nutritive riche en amidon (source de l'amidon du maïs consommé au Cameroun sous forme de pâte, bouillie, farine).
\end{itemize}
\textbf{Conclusion} : la double fécondation, propre aux Angiospermes, produit simultanément l'embryon ($2n$) et son tissu nourricier ($3n$), tout en restaurant la diploïdie de l'espèce à chaque génération. C'est une optimisation énergétique : l'albumen ne se développe que si la fécondation a eu lieu.
\end{exoresolu}

\courssub{Fiche méthode 6 : Expliquer la complémentarité méiose--fécondation dans la pérennité de l'espèce}

\begin{methode}[Construire un raisonnement scientifique sur la stabilité du caryotype et la diversité génétique]
\begin{enumerate}
\item \textbf{Poser le problème de la constance du nombre de chromosomes} : si les gamètes étaient diploïdes ($2n$), la fécondation doublerait le nombre de chromosomes à chaque génération ($2n \to 4n \to 8n \dots$) : instabilité létale.
\item \textbf{Montrer le rôle de la méiose (réduction + brassage)} :
\begin{itemize}
\item \textbf{Réduction} : passage de $2n$ à $n$ (anaphase I : séparation des homologues) $\to$ gamètes haploïdes.
\item \textbf{Brassage interchromosomique} : ségrégation indépendante des paires homologues en anaphase I $\to$ $2^{n}$ combinaisons chromosomiques différentes par gamète.
\item \textbf{Brassage intrachromosomique} : enjambements (crossing-over) en prophase I (chiasmas) $\to$ recombinaison d'allèles sur un même chromosome $\to$ diversité quasi infinie.
\end{itemize}
\item \textbf{Montrer le rôle de la fécondation (restauration + nouveau brassage)} :
\begin{itemize}
\item \textbf{Restauration} : fusion de deux noyaux haploïdes ($n + n = 2n$) $\to$ zygote diploïde, caryotype constant de l'espèce.
\item \textbf{Nouveau brassage} : rencontre aléatoire de deux gamètes issus d'individus différents (sauf autofécondation) $\to$ chaque zygote est une combinaison génétique unique.
\end{itemize}
\item \textbf{Conclure sur la pérennité} : l'alternance cyclique \textbf{Méiose ($2n \to n$) / Fécondation ($n + n \to 2n$)} maintient \textbf{constant} le nombre de chromosomes de l'espèce au fil des générations (\textbf{stabilité du caryotype}, condition de la cohérence spécifique) tout en générant une \textbf{diversité génétique maximale} (\textbf{nouveauté}, condition de l'adaptation aux changements environnementaux et de la survie évolutive de l'espèce).
\item \textbf{Chiffrer la diversité potentielle} (brassage interchromosomique seul) : nombre de zygotes différents possibles pour un couple $= 2^{n}_{\text{mâle}} \times 2^{n}_{\text{femelle}} = 2^{2n}$.
\end{enumerate}
\end{methode}

\begin{exoresolu}[Pérennité de l'espèce humaine (\emph{Homo sapiens})]
Chez l'Homme, $2n = 46$ chromosomes ($n = 23$ paires).\\
1) Montrer, à l'aide d'un schéma fonctionnel, comment le nombre de chromosomes reste constant d'une génération à l'autre.\\
2) Calculer le nombre de combinaisons génétiques possibles pour un enfant issu d'un couple, en ne considérant que le brassage interchromosomique. Commenter l'ampleur de ce nombre.
\tcblower
\textbf{1) Stabilité du stock chromosomique (Schéma fonctionnel).}
\begin{popfigure}
\begin{tikzpicture}[node distance=1.8cm, every node/.style={font=\small, align=center},
    box/.style={draw, thick, rounded corners, minimum width=2.2cm, minimum height=1cm},
    arrow/.style={->, thick, popdark, >=Stealth}]
\node[box, popblue, fill=popblueL, align=center] (p){Père\\ Somatique\\ $2n = 46$};
\node[box, poppink, fill=poppinkL, right=4cm of p, align=center] (m){Mère\\ Somatique\\ $2n = 46$};
\node[box, poporange, fill=poporangeL, below=1.5cm of p, align=center] (sp){Spermatogenèse\\ Méiose\\ $\downarrow$\\ Spermatozoïde\\ $n = 23$\\ $Q$ ADN};
\node[box, poporange, fill=poporangeL, below=1.5cm of m, align=center] (ov){Ovogenèse\\ Méiose\\ $\downarrow$\\ Ovocyte II (Ovule)\\ $n = 23$\\ $Q$ ADN};
\node[box, popgreen, fill=popgreenL, below=1.5cm of $(sp)!0.5!(ov)$, align=center] (z){Fécondation\\ (Trompe utérine)\\ $\downarrow$\\ Zygote\\ $2n = 46$\\ $2Q$ ADN};
\draw[arrow] (p) -- (sp) node[midway, left, font=\scriptsize, align=center]{Méiose\\$2n \to n$};
\draw[arrow] (m) -- (ov) node[midway, right, font=\scriptsize, align=center]{Méiose\\$2n \to n$};
\draw[arrow] (sp) -- (z);
\draw[arrow] (ov) -- (z);
\node[below=0.2cm of z, font=\small\itshape, popdark] {Caryogamie : $n + n \to 2n$};
\end{tikzpicture}
\legende{Alternance méiose (réduction $2n \to n$) et fécondation (restauration $n + n \to 2n$) : le nombre de chromosomes de l'espèce humaine reste constant à 46 (23 paires) de génération en génération.}
\end{popfigure}
La méiose produit des gamètes haploïdes ($n = 23$) ; la fécondation réunit deux gamètes et restaure la diploïdie ($2n = 46$) : le caryotype de l'espèce est stable au fil des générations.\\
\textbf{2) Diversité des zygotes (Brassage interchromosomique seul).}\\
Chaque parent produit $2^{n} = 2^{23} = 8\,388\,608$ types de gamètes génétiquement différents par ségrégation indépendante des chromosomes homologues. Le nombre de zygotes génétiquement distincts qu'un couple peut théoriquement produire est :
\[ 2^{23} \times 2^{23} = 2^{46} = 70\,368\,744\,177\,664 \approx 7{,}0 \times 10^{13} \text{ combinaisons,} \]
soit environ \textbf{70 000 milliards} de combinaisons possibles (70 milliards de milliards en échelle courte, 70 mille milliards en échelle longue).
\textbf{Commentaire} : ce nombre astronomique (bien supérieur à la population humaine totale historique) ne tient pas compte des enjambements (brassage intrachromosomique) qui multiplient encore la diversité par des facteurs considérables. Cela explique pourquoi deux frères et sœurs (sauf jumeaux monozygotes) sont toujours génétiquement uniques.
\textbf{Conclusion} : la méiose garantit la réduction chromosomique et le brassage des allèles (sources de variabilité), la fécondation restaure la diploïdie en associant deux patrimoines différents (source de nouveauté) ; leur complémentarité stricte assure à la fois la \textbf{stabilité du caryotype} (identité de l'espèce) et l'\textbf{immense diversité génétique} (potentiel d'adaptation), fondements indissociables de la pérennité de l'espèce.
\end{exoresolu}

\begin{attention}[Les 5 erreurs classiques qui coûtent des points au Baccalauréat]
\begin{enumerate}
\item \textbf{Confondre quantité d'ADN et nombre de chromosomes} : un spermatocyte II possède $n$ chromosomes mais $2Q$ d'ADN (chromosomes à deux chromatides) ; un spermatozoïde possède $n$ chromosomes et $Q$ d'ADN. \textbf{Toujours préciser les deux valeurs séparément} dans les réponses.
\item \textbf{Attribuer la réduction chromosomique à la méiose II} : c'est l'\textbf{anaphase I} (séparation des chromosomes homologues) qui est réductionnelle ($2n \to n$) ; la méiose II ne sépare que les chromatides sœurs (division équationnelle, $n \to n$).
\item \textbf{Écrire que l'albumen (endosperme) est diploïde} : chez les Angiospermes, l'albumen est \textbf{triploïde ($3n$)}, issu de la fusion d'un noyau spermatique ($n$) avec la cellule centrale ($2n$). Seul le zygote est diploïde ($2n$).
\item \textbf{Dire que l'ovogenèse produit 4 ovules} : elle ne produit qu'\textbf{un seul} ovocyte fonctionnel (ovule) par ovocyte I entrant en méiose ; les trois autres produits sont des globules polaires qui dégénèrent (apoptose).
\item \textbf{Schémas non conformes aux attendus du Bac} : oublier le titre (objet + espèce + vue), croiser les traits de légende, légender avec des termes vagues (« tube », « boule », « machin ») au lieu du vocabulaire anatomique exact (canal déférent, anthère, ovule, zone pellucide, tube pollinique, sac embryonnaire...). Un schéma sans titre ni légendes précises perd au minimum la moitié de ses points.
\end{enumerate}
\end{attention}

\newpage

\coursec{Exercices}

\courssub{Je vérifie mes connaissances}

\begin{exercice}[QCM : méiose et fécondation]
Pour chaque question, choisis la (ou les) bonne(s) réponse(s) et justifie brièvement ton choix.
\begin{enumerate}
\item La méiose est une division cellulaire qui :
\begin{enumerate}
\item[a)] produit deux cellules filles identiques à la cellule mère ;
\item[b)] produit quatre cellules filles haploïdes à partir d'une cellule mère diploïde ;
\item[c)] comporte deux divisions successives précédées chacune d'une réplication de l'ADN ;
\item[d)] comporte une seule réplication de l'ADN suivie de deux divisions.
\end{enumerate}
\item Chez un Mammifère dont le caryotype est $2n = 46$ chromosomes, un spermatozoïde contient :
\begin{enumerate}
\item[a)] 46 chromosomes ; \item[b)] 23 chromosomes ; \item[c)] 92 chromosomes ; \item[d)] 23 paires de chromosomes.
\end{enumerate}
\item La double fécondation chez les Spermaphytes produit :
\begin{enumerate}
\item[a)] deux zygotes ; \item[b)] un zygote diploïde et une cellule de l'albumen triploïde ; \item[c)] un zygote triploïde et un albumen diploïde ; \item[d)] deux embryons jumeaux.
\end{enumerate}
\item Au cours de la spermatogenèse, la spermiogenèse est :
\begin{enumerate}
\item[a)] une division supplémentaire des spermatides ;
\item[b)] la transformation des spermatides en spermatozoïdes sans division ;
\item[c)] la réplication de l'ADN des spermatogonies ;
\item[d)] la fusion de deux spermatides.
\end{enumerate}
\end{enumerate}
\end{exercice}

\begin{exercice}[Vrai ou faux justifié]
Indique si chacune des affirmations suivantes est vraie ou fausse, puis justifie ta réponse en une ou deux phrases.
\begin{enumerate}
\item La réplication de l'ADN a lieu avant chacune des deux divisions de la méiose.
\item Le crossing-over se produit au cours de la prophase I de la méiose, entre chromosomes homologues.
\item Chez les Spermaphytes, le grain de pollen est le gamète mâle.
\item La fécondation rétablit la diploïdie dans la cellule-\oe uf.
\item Chez les Mammifères, l'ovogenèse produit quatre ovules fonctionnels à partir d'une ovogonie.
\end{enumerate}
\end{exercice}

\begin{exercice}[Définitions]
Définis de manière précise et concise les termes suivants, en illustrant chacun par un exemple :
\begin{enumerate}
\item gonade ; \item méiose ; \item gamétogenèse ; \item fécondation ; \item double fécondation ; \item caryotype.
\end{enumerate}
\end{exercice}

\begin{exercice}[Calculs de ploïdie et de quantité d'ADN]
Chez le Rat, espèce utilisée dans les laboratoires camerounais, $2n = 42$ chromosomes et la quantité d'ADN d'un gamète est $Q = \SI{3,2}{\pico\gram}$ (picogrammes).
\begin{enumerate}
\item Donne le nombre de chromosomes et la quantité d'ADN d'une spermatogonie en phase G1.
\item Donne la quantité d'ADN d'une spermatogonie à la fin de la phase S, juste avant la première division de méiose.
\item Donne le nombre de chromosomes et la quantité d'ADN d'un spermatocyte II et d'une spermatide.
\item Quelle quantité d'ADN retrouve-t-on dans le noyau du zygote juste après la fécondation ? Justifie.
\end{enumerate}
\end{exercice}

\courssub{Je m'entraîne}

\begin{exercice}[Courbe d'évolution de la quantité d'ADN au cours de la spermatogenèse]
On a mesuré la quantité d'ADN par noyau dans les cellules de la lignée germinale mâle d'un Mammifère. Les résultats sont consignés dans le tableau ci-dessous ($Q$ = quantité d'ADN d'un gamète).
\begin{center}
\begin{tabularx}{\linewidth}{|l|X|X|X|X|X|}\hline
\rowcolor{popblueL} Stade & Spermatogonie (G1) & Spermatocyte I (début) & Spermatocyte II & Spermatide & Spermatozoïde \\ \hline
Quantité d'ADN & $2Q$ & $4Q$ & $2Q$ & $Q$ & $Q$ \\ \hline
\end{tabularx}
\end{center}
\begin{enumerate}
\item Trace la courbe d'évolution de la quantité d'ADN au cours de la spermatogenèse (quantité d'ADN en ordonnée, stades en abscisse).
\item Indique sur la courbe les phases de réplication de l'ADN, de première division, de deuxième division et de spermiogenèse.
\item Précise la ploïdie ($n$ ou $2n$) des cellules à chaque palier de la courbe.
\item Explique pourquoi la quantité d'ADN passe de $4Q$ à $2Q$ puis de $2Q$ à $Q$.
\end{enumerate}
\end{exercice}

\begin{exercice}[Remise en ordre d'électronographies de la méiose]
Six électronographies, notées A, B, C, D, E et F, montrent des étapes de la méiose chez une cellule animale à $2n = 4$ chromosomes :
\begin{itemize}
\item A : quatre noyaux haploïdes en voie de reconstitution, les chromosomes se décondensent ;
\item B : les chromosomes homologues appariés forment des bivalents ; des échanges de fragments sont visibles ;
\item C : les chromosomes à deux chromatides sont alignés sur la plaque équatoriale, chacun indépendamment ;
\item D : les bivalents sont alignés sur la plaque équatoriale ;
\item E : les chromosomes à une chromatide migrent vers les pôles ;
\item F : les chromosomes homologues, encore à deux chromatides, se séparent et migrent vers les pôles opposés.
\end{itemize}
\begin{enumerate}
\item Remets ces six étapes dans l'ordre chronologique de la méiose.
\item Nomme chaque étape (division et phase).
\item Justifie ton classement des étapes B, D et F par le comportement des chromosomes.
\item Précise à quelle(s) étape(s) on observe encore des chromosomes à deux chromatides.
\end{enumerate}
\end{exercice}

\begin{exercice}[Coupe de tube séminifère]
Le schéma d'une coupe transversale de tube séminifère de Rat montre, de la périphérie vers la lumière : des cellules de Sertoli, des spermatogonies, des spermatocytes I, des spermatocytes II, des spermatides et des spermatozoïdes dont les flagelles plongent dans la lumière.
\begin{enumerate}
\item Réalise un schéma légendé de cette coupe en respectant la disposition des cellules.
\item Identifie les cellules de la lignée germinale aux différents stades d'évolution.
\item Précise, pour chaque type cellulaire, sa ploïdie et sa quantité d'ADN (en fonction de $Q$).
\item Quel rôle jouent les cellules de Sertoli dans la spermatogenèse ?
\end{enumerate}
\end{exercice}

\begin{exercice}[Dissection et observation d'une fleur]
Au cours d'une séance de TP au lycée de Nkongsamba, des élèves disposent de fleurs de \emph{Hibiscus} (foléré) fraîchement cueillies.
\begin{enumerate}
\item Décris le protocole d'observation permettant de mettre en évidence les organes reproducteurs de la fleur (matériel, gestes, précautions).
\item Nomme la gonade mâle et la gonade femelle de la fleur et précise leur localisation.
\item Réalise un schéma légendé de la fleur en coupe longitudinale montrant les organes reproducteurs.
\item Où se forment les gamètes mâles et femelles chez cette Spermaphyte ?
\end{enumerate}
\end{exercice}

\courssub{Je me prépare au Bac}

\begin{exercice}[Type Bac : spermatogenèse et quantité d'ADN]
Dans un laboratoire de l'Université de Yaoundé I, on étudie la spermatogenèse chez un Mammifère dont le caryotype est $2n = 38$ chromosomes. La quantité d'ADN d'un spermatozoïde est $Q = \SI{3,0}{\pico\gram}$. On mesure la quantité d'ADN par noyau dans des cellules prélevées dans les tubes séminifères ; on obtient quatre populations de cellules contenant respectivement $3,0$ ; $6,0$ ; $12,0$ et $6,0$ picogrammes d'ADN, la population à \SI{6,0}{\pico\gram} regroupant en réalité deux types cellulaires morphologiquement distincts.
\begin{enumerate}
\item Attribue à chaque population de cellules son identité (spermatogonie, spermatocyte I, spermatocyte II, spermatide ou spermatozoïde) en justifiant par la quantité d'ADN.
\item Explique pourquoi deux types cellulaires différents peuvent contenir la même quantité d'ADN de \SI{6,0}{\pico\gram} ; comment les distinguer au microscope ?
\item Trace la courbe d'évolution de la quantité d'ADN au cours de la spermatogenèse complète, en indiquant les phases de réplication, les deux divisions de la méiose et la spermiogenèse.
\item Donne le nombre de chromosomes d'un spermatocyte I, d'un spermatocyte II et d'un spermatozoïde de cette espèce.
\item Un produit expérimental bloque la deuxième division de la méiose sans affecter la première. Quelle serait la ploïdie des spermatides obtenues ? Quelle serait la conséquence pour le zygote si un tel gamète participait à la fécondation ?
\end{enumerate}
\end{exercice}

\begin{exercice}[Type Bac : double fécondation chez les Spermaphytes]
Au champ expérimental de l'IRAD de Nkolbisson, on observe au microscope une coupe d'ovaire de Maïs (\emph{Zea mays}, $2n = 20$) au moment de la fécondation. Le document montre un tube pollinique pénétrant dans l'ovule par le micropyle et libérant deux noyaux mâles au voisinage du sac embryonnaire qui contient huit noyaux haploïdes, dont l'oosphère et les deux noyaux polaires.
\begin{enumerate}
\item Schématise et légende l'ovule au moment de la double fécondation : tube pollinique, micropyle, noyaux mâles, oosphère, noyaux polaires, sac embryonnaire.
\item Décris, étape par étape, le devenir des deux noyaux mâles libérés par le tube pollinique.
\item Indique la ploïdie et le nombre de chromosomes du zygote et de la cellule-mère de l'albumen issus de cette double fécondation.
\item Quel est le devenir du zygote et de l'albumen au cours du développement de la graine ? Quel rôle joue l'albumen pour la plantule ?
\item Compare, sous forme d'un tableau construit, la fécondation des Mammifères et celle des Spermaphytes selon les critères suivants : lieu de la fécondation, nombre de fusions nucléaires, nature des produits formés, devenir de ces produits.
\end{enumerate}
\end{exercice}

\begin{exercice}[Type Bac : méiose, fécondation et pérennité de l'espèce]
Des chercheurs du Centre Pasteur du Cameroun étudient un parasite dont le caryotype normal est $2n = 14$ chromosomes. Ils traitent des cellules germinales de ce parasite avec une substance qui bloque la méiose : les gamètes produits contiennent alors 14 chromosomes au lieu de 7. On fait ensuite s'unir deux de ces gamètes anormaux.
\begin{enumerate}
\item Quelle est la ploïdie des gamètes produits après blocage de la méiose ? Justifie.
\item Calcule le nombre de chromosomes du zygote issu de l'union de deux gamètes anormaux, puis celui du zygote qui serait issu de la génération suivante si le phénomène se répétait.
\item Démontre, par le raisonnement, que la méiose est indispensable à la constance du caryotype d'une génération à l'autre chez les espèces à reproduction sexuée.
\item Rappelle le rôle complémentaire de la fécondation dans cette constance.
\item Question de synthèse : en t'appuyant sur des faits précis (comportement des chromosomes, brassage génétique, évolution de la quantité d'ADN), montre que la méiose et la fécondation sont deux phénomènes complémentaires assurant à la fois la stabilité et la diversité génétique, conditions de la pérennité de l'espèce.
\end{enumerate}
\end{exercice}

\newpage

\coursec{Corrigés des exercices}

\coursec{Exercices corrigés : Méiose, fécondation et pérennité de l'espèce}

\begin{corrige}[Exercice 1 — QCM : méiose et fécondation]
\begin{enumerate}
\item \textbf{Bonnes réponses : b) et d).} La méiose est une succession de \cle{deux divisions} (méiose I et méiose II) précédées d'une \cle{seule réplication} de l'ADN (phase S de l'interphase) : une cellule mère diploïde ($2n$) donne quatre cellules filles haploïdes ($n$). La réponse a) décrit la mitose ; la réponse c) est fausse car il n'y a pas de réplication entre les deux divisions.
\item \textbf{Bonne réponse : b).} Le spermatozoïde est un gamète haploïde : $n = 46/2 = 23$ chromosomes. Attention au piège d) : les chromosomes d'un gamète ne sont \emph{pas} appariés, il n'y a plus de « paires ».
\item \textbf{Bonne réponse : b).} Chez les Spermaphytes, un noyau mâle ($n$) fusionne avec l'oosphère ($n$) pour donner le \cle{zygote} diploïde ($2n$) ; le second noyau mâle ($n$) fusionne avec les deux noyaux polaires ($n + n$) pour donner la cellule-mère de l'\cle{albumen}, triploïde ($3n$). Il n'y a qu'un seul zygote, donc a) et d) sont fausses.
\item \textbf{Bonne réponse : b).} La \cle{spermiogenèse} est la différenciation des spermatides (haploïdes) en spermatozoïdes : formation du flagelle, de l'acrosome, condensation du noyau, sans aucune division ni modification de la quantité d'ADN.
\end{enumerate}
\end{corrige}

\begin{corrige}[Exercice 2 — Vrai ou faux justifié]
\begin{enumerate}
\item \textbf{Faux.} La réplication de l'ADN a lieu une seule fois, pendant l'interphase qui précède la méiose I. Entre la première et la deuxième division, il n'y a pas de phase S : c'est ce qui permet le passage de $4Q$ à $Q$.
\item \textbf{Vrai.} Le \cle{crossing-over} (enjambement) se produit en prophase I (stade pachytène), entre chromatides de chromosomes \emph{homologues} appariés en bivalents. Il réalise un brassage intrachromosomique.
\item \textbf{Faux.} Le grain de pollen n'est pas le gamète : c'est une structure qui \emph{contient} (ou produira) les gamètes mâles. Les véritables gamètes mâles sont les deux noyaux spermatiques véhiculés par le tube pollinique.
\item \textbf{Vrai.} La fécondation est la fusion de deux gamètes haploïdes ($n + n$) : la cellule-\oe uf (zygote) redevient diploïde ($2n$), ce qui rétablit le caryotype caractéristique de l'espèce.
\item \textbf{Faux.} L'ovogenèse produit à partir d'une ovogonie \textbf{un seul ovule fonctionnel} (plus deux ou trois globules polaires qui dégénèrent), car les divisions sont inégales : le cytoplasme est conservé dans une seule cellule. C'est la spermatogenèse qui produit quatre gamètes fonctionnels.
\end{enumerate}
\end{corrige}

\begin{corrige}[Exercice 3 — Définitions]
\begin{enumerate}
\item \textbf{Gonade} : glande génitale qui produit les gamètes et sécrète les hormones sexuelles. \emph{Exemples :} le testicule et l'ovaire chez les Mammifères ; l'anthère (gonade mâle) et l'ovaire (gonade femelle) de la fleur chez les Spermaphytes.
\item \textbf{Méiose} : succession de deux divisions cellulaires précédées d'une seule réplication de l'ADN, qui transforme une cellule diploïde ($2n$) en quatre cellules haploïdes ($n$) génétiquement différentes. \emph{Exemple :} la méiose des spermatocytes I dans les tubes séminifères du Rat.
\item \textbf{Gamétogenèse} : ensemble des processus de formation des gamètes à partir des cellules germinales, incluant la multiplication, la croissance, la méiose et la différenciation. \emph{Exemples :} la spermatogenèse (formation des spermatozoïdes) et l'ovogenèse (formation de l'ovule).
\item \textbf{Fécondation} : union d'un gamète mâle et d'un gamète femelle, avec fusion de leurs noyaux haploïdes (caryogamie), donnant une cellule-\oe uf ou zygote diploïde. \emph{Exemple :} la rencontre spermatozoïde--ovule dans le tiers externe de la trompe de Fallope.
\item \textbf{Double fécondation} : particularité des Spermaphytes où les deux noyaux mâles du tube pollinique réalisent deux fusions : l'un avec l'oosphère (zygote, $2n$), l'autre avec les deux noyaux polaires (cellule-mère de l'albumen, $3n$). \emph{Exemple :} la fécondation chez le Maïs.
\item \textbf{Caryotype} : ensemble des chromosomes d'une cellule, classés par paires d'homologues selon leur taille et leur forme, caractéristique de l'espèce. \emph{Exemple :} le caryotype humain comporte $2n = 46$ chromosomes, soit 23 paires.
\end{enumerate}
\end{corrige}

\begin{corrige}[Exercice 4 — Calculs de ploïdie et de quantité d'ADN]
Données : $2n = 42$ ; le gamète contient $n = 21$ chromosomes et $Q = \SI{3,2}{\pico\gram}$ d'ADN.
\begin{enumerate}
\item \textbf{Spermatogonie en G1} : cellule diploïde non répliquée. Elle contient $2n = 42$ chromosomes (à une chromatide) et $2Q = 2 \times 3,2 = \SI{6,4}{\pico\gram}$ d'ADN.
\item \textbf{Fin de la phase S} : l'ADN a été répliqué, chaque chromosome comporte deux chromatides, mais le nombre de chromosomes reste $2n = 42$. Quantité d'ADN : $4Q = 4 \times 3,2 = \SI{12,8}{\pico\gram}$.
\item \textbf{Spermatocyte II} : issu de la première division (séparation des homologues), il est haploïde : $n = 21$ chromosomes, mais chacun à deux chromatides, donc $2Q = \SI{6,4}{\pico\gram}$ d'ADN. \textbf{Spermatide} : issue de la deuxième division (séparation des chromatides), elle contient $n = 21$ chromosomes à une chromatide, soit $Q = \SI{3,2}{\pico\gram}$ d'ADN.
\item \textbf{Zygote juste après fécondation} : fusion de deux gamètes ($n + n = 2n = 42$ chromosomes) ; quantité d'ADN : $Q + Q = 2Q = \SI{6,4}{\pico\gram}$. La fécondation rétablit la diploïdie et la quantité d'ADN caractéristique de l'espèce.
\end{enumerate}
\end{corrige}

\begin{attention}[Vérification des calculs]
$2 \times 3,2 = 6,4$ ; $4 \times 3,2 = 12,8$ ; $42/2 = 21$. Ne jamais confondre « nombre de chromosomes » (inchangé pendant la phase S) et « quantité d'ADN » (doublée pendant la phase S).
\end{attention}

\begin{corrige}[Exercice 5 — Courbe d'évolution de la quantité d'ADN au cours de la spermatogenèse]
\textbf{1. et 2. Courbe et annotation.}
\begin{popfigure}
\begin{center}
\begin{tikzpicture}
\begin{axis}[width=0.95\linewidth,height=6.5cm,
xlabel={Stades de la spermatogenèse},ylabel={Quantité d'ADN par noyau},
symbolic x coords={G1,S,SpI,SpII,Sde,Szo},xtick=data,
ymin=0,ymax=5,ytick={1,2,3,4},yticklabels={$Q$,$2Q$,$3Q$,$4Q$},
grid=major,grid style={popline!40}]
\addplot[popcurve,mark=*,mark size=1.5pt] coordinates {(G1,2) (S,4) (SpI,4) (SpII,2) (Sde,1) (Szo,1)};
\draw[poporange,thick,->] (axis cs:S,4.4) -- (axis cs:SpI,4.4) node[midway,above,font=\small]{méiose I};
\draw[popgreen,thick,->] (axis cs:SpII,1.5) -- (axis cs:Sde,1.5);
\node[popgreen,font=\small] at (axis cs:SpII,1.8) {};
\draw[popgold,thick,->] (axis cs:Sde,0.5) -- (axis cs:Szo,0.5) node[midway,below,font=\small]{spermiogenèse};
\node[poporange,font=\small] at (axis cs:G1,2.4) {réplication (phase S)};
\draw[poporange,thick,->] (axis cs:G1,2.2) -- (axis cs:S,3.9);
\node[popdark,font=\small] at (axis cs:SpII,2.35) {méiose II};
\end{axis}
\end{tikzpicture}
\end{center}
\legende{Évolution de la quantité d'ADN au cours de la spermatogenèse. G1 : spermatogonie ; S : phase de réplication ; SpI : spermatocyte I ; SpII : spermatocyte II ; Sde : spermatide ; Szo : spermatozoïde.}
\end{popfigure}

\textbf{3. Ploïdie à chaque palier :}
\begin{center}
\begin{tabularx}{\linewidth}{|l|X|X|X|X|X|}\hline
\rowcolor{popblueL} Stade & G1 & S / SpI & SpII & Spermatide & Spermatozoïde \\ \hline
Ploïdie & $2n$ & $2n$ & $n$ & $n$ & $n$ \\ \hline
ADN & $2Q$ & $4Q$ & $2Q$ & $Q$ & $Q$ \\ \hline
\end{tabularx}
\end{center}

\textbf{4. Explication des chutes :}
\begin{itemize}
\item Passage de $4Q$ à $2Q$ : la \cle{méiose I} sépare les chromosomes \emph{homologues} (disjonction en anaphase I) ; chaque cellule fille ne reçoit qu'un chromosome de chaque paire, encore à deux chromatides : elle devient haploïde ($n$) avec $2Q$ d'ADN.
\item Passage de $2Q$ à $Q$ : la \cle{méiose II} sépare les \emph{chromatides s\oe urs} ; chaque spermatide reçoit $n$ chromosomes à une chromatide, soit $Q$ d'ADN.
\end{itemize}
\end{corrige}

\begin{corrige}[Exercice 6 — Remise en ordre d'électronographies de la méiose]
\textbf{1. Ordre chronologique :} B $\rightarrow$ D $\rightarrow$ F $\rightarrow$ C $\rightarrow$ E $\rightarrow$ A.

\textbf{2. Nom des étapes :}
\begin{center}
\begin{tabularx}{\linewidth}{|l|X|X|}\hline
\rowcolor{popblueL} Cliché & Étape & Observation décisive \\ \hline
B & Prophase I & Bivalents + crossing-over \\ \hline
D & Métaphase I & Bivalents sur la plaque équatoriale \\ \hline
F & Anaphase I & Séparation des homologues (2 chromatides) \\ \hline
C & Métaphase II & Chromosomes alignés individuellement \\ \hline
E & Anaphase II & Migration des chromosomes à 1 chromatide \\ \hline
A & Télophase II & 4 noyaux haploïdes, décondensation \\ \hline
\end{tabularx}
\end{center}

\textbf{3. Justification du classement B--D--F :} en \textbf{B} (prophase I), les homologues s'apparient (synapsis) et échangent des fragments (crossing-over) : c'est la première étape reconnaissable. En \textbf{D} (métaphase I), les bivalents sont encore intacts mais alignés sur l'équateur : l'appariement persiste, donc D suit B. En \textbf{F} (anaphase I), les homologues se séparent en restant à deux chromatides : c'est la disjonction, qui suit nécessairement l'alignement des bivalents. L'ordre logique est donc appariement $\rightarrow$ alignement $\rightarrow$ séparation.

\textbf{4. Chromosomes à deux chromatides :} on les observe aux étapes \textbf{B, D, F et C} (toute la méiose I et la métaphase II). À partir de l'étape E (anaphase II), les chromatides s\oe urs se sont séparées : les chromosomes sont à une chromatide.
\end{corrige}

\begin{corrige}[Exercice 7 — Coupe de tube séminifère]
\textbf{1. et 2. Schéma légendé de la coupe.}
\begin{popfigure}
\begin{center}
\begin{tikzpicture}[scale=1]
\draw[fill=popcream,thick] (0,0) circle (3);
\draw[fill=popblueL] (0,0) circle (0.9);
\node[font=\small] at (0,0) {lumière};
% spermatozoïdes
\foreach \a in {30,90,150,210,270,330}{
\draw[popdark,fill=popdark] (\a:1.15) circle (0.09);
\draw[popdark] (\a:1.15) -- (\a:0.55);}
% spermatides
\foreach \a in {0,60,120,180,240,300}{
\draw[fill=popgold,draw=popdark] (\a:1.7) circle (0.18);}
% spermatocytes II
\foreach \a in {45,135,225,315}{
\draw[fill=popgreenL,draw=popdark] (\a:2.15) circle (0.22);}
% spermatocytes I
\foreach \a in {20,110,200,290}{
\draw[fill=poppinkL,draw=popdark] (\a:2.5) circle (0.27);}
% spermatogonies
\foreach \a in {70,160,250,340}{
\draw[fill=poporangeL,draw=popdark] (\a:2.82) circle (0.16);}
% Sertoli
\draw[fill=popgreen!50,draw=popdark] (55:2.45) ellipse (0.28 and 0.45);
\draw[fill=popgreen!50,draw=popdark] (235:2.45) ellipse (0.28 and 0.45);
% légendes
\node[font=\scriptsize,right] at (3.3,2.4) {1 : spermatogonie};
\node[font=\scriptsize,right] at (3.3,1.9) {2 : spermatocyte I};
\node[font=\scriptsize,right] at (3.3,1.4) {3 : spermatocyte II};
\node[font=\scriptsize,right] at (3.3,0.9) {4 : spermatide};
\node[font=\scriptsize,right] at (3.3,0.4) {5 : spermatozoïde};
\node[font=\scriptsize,right] at (3.3,-0.1) {6 : cellule de Sertoli};
\end{tikzpicture}
\end{center}
\legende{Coupe transversale de tube séminifère : les cellules de la lignée germinale sont disposées de la périphérie (spermatogonies) vers la lumière (spermatozoïdes), au gré de leur maturation.}
\end{popfigure}

\textbf{3. Ploïdie et quantité d'ADN :}
\begin{center}
\begin{tabularx}{\linewidth}{|l|X|X|}\hline
\rowcolor{popblueL} Cellule & Ploïdie & Quantité d'ADN \\ \hline
Spermatogonie & $2n$ & $2Q$ (G1) à $4Q$ (après S) \\ \hline
Spermatocyte I & $2n$ & $4Q$ \\ \hline
Spermatocyte II & $n$ & $2Q$ \\ \hline
Spermatide & $n$ & $Q$ \\ \hline
Spermatozoïde & $n$ & $Q$ \\ \hline
\end{tabularx}
\end{center}

\textbf{4. Rôle des cellules de Sertoli :} ce sont des cellules \emph{nourricières} et de soutien : elles nourrissent les cellules germinales en cours de maturation, les guident vers la lumière, phagocytent les déchets (corps résiduels de la spermiogenèse) et constituent la barrière hémato-testiculaire qui protège les gamètes en formation.
\end{corrige}

\begin{corrige}[Exercice 8 — Dissection et observation d'une fleur]
\textbf{1. Protocole d'observation :}
\begin{itemize}
\item \textbf{Matériel :} fleur fraîche de \emph{Hibiscus}, loupe binoculaire ou loupe à main, aiguilles fines, pince brucelles, lame et lamelle, scalpel ou lame de rasoir, boîte de Pétri, eau.
\item \textbf{Gestes :} observer d'abord la fleur entière à l'\oe il nu ; écarter délicatement le calice et la corolle avec les brucelles ; dégager la colonne staminale (étamines soudées en tube autour du pistil, caractéristique de \emph{Hibiscus}) ; prélever quelques anthères et les écraser entre lame et lamelle pour observer les grains de pollen ; pratiquer ensuite une coupe longitudinale de l'ovaire au scalpel pour observer les ovules à la loupe.
\item \textbf{Précautions :} manipuler le scalpel en coupant vers l'extérieur ; ne pas écraser les organes ; travailler sur fleur fraîche.
\end{itemize}

\textbf{2. Gonades :} la \cle{gonade mâle} est l'\textbf{anthère} (située au sommet des étamines, sur la colonne staminale) ; la \cle{gonade femelle} est l'\textbf{ovaire} (situé à la base du pistil, au centre de la fleur, surmonté du style et des stigmates ramifiés).

\textbf{3. Schéma de la fleur en coupe longitudinale.}
\begin{popfigure}
\begin{center}
\begin{tikzpicture}[scale=1.1]
% pétales
\draw[fill=poppinkL,draw=poppink] (-1.8,1.2) ellipse (0.5 and 1.4);
\draw[fill=poppinkL,draw=poppink] (1.8,1.2) ellipse (0.5 and 1.4);
% sépales
\draw[fill=popgreenL,draw=popgreen] (-1.2,-0.4) ellipse (0.35 and 0.9);
\draw[fill=popgreenL,draw=popgreen] (1.2,-0.4) ellipse (0.35 and 0.9);
% pistil
\draw[fill=popgoldL,draw=popdark] (0,-0.9) ellipse (0.55 and 0.6);
\draw[popdark,thick] (0,-0.3) -- (0,2.2);
\draw[fill=poporange,draw=popdark] (0,2.35) circle (0.18);
% ovules
\draw[fill=white,draw=popdark] (-0.25,-0.9) circle (0.12);
\draw[fill=white,draw=popdark] (0.25,-0.9) circle (0.12);
% étamines
\draw[popdark,thick] (-0.7,0.4) -- (-0.7,1.6);
\draw[popdark,thick] (0.7,0.4) -- (0.7,1.6);
\draw[fill=popgold,draw=popdark] (-0.7,1.75) ellipse (0.15 and 0.22);
\draw[fill=popgold,draw=popdark] (0.7,1.75) ellipse (0.15 and 0.22);
% réceptacle
\draw[popdark,thick] (0,-1.5) -- (0,-2.1);
% légendes
\node[font=\scriptsize,right] at (2.6,2.3) {1 : stigmate};
\draw[->,popmuted] (2.6,2.2) -- (0.2,2.35);
\node[font=\scriptsize,right] at (2.6,1.7) {2 : anthère (gonade mâle)};
\draw[->,popmuted] (2.6,1.6) -- (0.85,1.75);
\node[font=\scriptsize,right] at (2.6,1.1) {3 : filet de l'étamine};
\draw[->,popmuted] (2.6,1.0) -- (0.75,1.0);
\node[font=\scriptsize,right] at (2.6,0.5) {4 : style};
\draw[->,popmuted] (2.6,0.4) -- (0.05,0.8);
\node[font=\scriptsize,right] at (2.6,-0.1) {5 : pétale};
\draw[->,popmuted] (2.6,-0.2) -- (1.9,0.8);
\node[font=\scriptsize,left] at (-2.6,0.5) {6 : sépale};
\draw[->,popmuted] (-2.6,0.4) -- (-1.3,0.1);
\node[font=\scriptsize,left] at (-2.6,-0.3) {7 : ovaire (gonade femelle)};
\draw[->,popmuted] (-2.6,-0.4) -- (-0.6,-0.8);
\node[font=\scriptsize,left] at (-2.6,-1.1) {8 : ovule};
\draw[->,popmuted] (-2.6,-1.2) -- (-0.3,-0.95);
\node[font=\scriptsize,left] at (-2.6,-1.9) {9 : pédoncule floral};
\draw[->,popmuted] (-2.6,-2.0) -- (-0.05,-2.0);
\end{tikzpicture}
\end{center}
\legende{Coupe longitudinale schématique d'une fleur de Spermaphyte : les organes reproducteurs sont les étamines (anthères, gonade mâle) et le pistil (ovaire contenant les ovules, gonade femelle).}
\end{popfigure}

\textbf{4. Lieu de formation des gamètes :} les gamètes mâles (noyaux spermatiques) se forment dans les \textbf{anthères}, à partir des grains de pollen issus de la méiose des cellules-mères du pollen ; le gamète femelle (oosphère) se forme dans les \textbf{ovules} contenus dans l'ovaire, au sein du sac embryonnaire issu de la méiose d'une cellule-mère du sac embryonnaire.
\end{corrige}

\begin{corrige}[Exercice 9 — Type Bac : spermatogenèse et quantité d'ADN]
\textbf{1. Identification des populations} ($2n = 38$ ; $Q = \SI{3,0}{\pico\gram}$) :
\begin{center}
\begin{tabularx}{\linewidth}{|l|X|X|}\hline
\rowcolor{popblueL} ADN par noyau & Identité & Justification \\ \hline
\SI{3,0}{\pico\gram} ($Q$) & Spermatides et spermatozoïdes & $n$ chromosomes à 1 chromatide \\ \hline
\SI{6,0}{\pico\gram} ($2Q$) & Spermatogonies (G1) \emph{et} spermatocytes II & $2n$ non répliqué, ou $n$ à 2 chromatides \\ \hline
\SI{12,0}{\pico\gram} ($4Q$) & Spermatocytes I & $2n$ chromosomes à 2 chromatides \\ \hline
\end{tabularx}
\end{center}

\textbf{2. Deux types cellulaires à \SI{6,0}{\pico\gram} :} la spermatogonie en G1 contient $2n = 38$ chromosomes à \emph{une} chromatide ($2Q$), tandis que le spermatocyte II contient $n = 19$ chromosomes à \emph{deux} chromatides ($2Q$) : même masse d'ADN, organisation différente. Au microscope, on les distingue par le \textbf{nombre de chromosomes} (38 contre 19) et par l'\textbf{aspect} des chromosomes (simples contre doubles, en bâtonnets dédoublés).

\textbf{3. Courbe :} identique en forme à celle de l'exercice 5, avec les valeurs $2Q = 6,0$ ; $4Q = 12,0$ ; $2Q = 6,0$ ; $Q = 3,0$ picogrammes. On y reporte : montée de $6,0$ à $12,0$ (phase S, réplication), palier à $12,0$ (spermatocyte I), chute à $6,0$ (méiose I), chute à $3,0$ (méiose II), palier à $3,0$ (spermiogenèse, sans variation d'ADN).

\textbf{4. Nombre de chromosomes :} spermatocyte I : $2n = 38$ chromosomes (à deux chromatides) ; spermatocyte II : $n = 19$ chromosomes (à deux chromatides) ; spermatozoïde : $n = 19$ chromosomes (à une chromatide).

\textbf{5. Blocage de la méiose II :} la première division s'effectue normalement : les cellules obtenues sont haploïdes ($n = 19$) mais à chromosomes doubles (chaque chromosome comporte deux chromatides). Si la deuxième division est bloquée, les chromatides s\oe urs ne se séparent pas : les cellules issues de ce blocage conservent $2Q$ d'ADN et un stock chromosomique non réduit (19 chromosomes à deux chromatides, équivalent en ADN à une cellule $2n$ en G1). Un tel gamète non réduit, fusionnant avec un gamète normal ($n = 19$, $Q$), donnerait un zygote \textbf{triploïde ($3n = 57$ chromosomes, $3Q$ d'ADN)}, généralement non viable ou stérile : d'où la nécessité absolue de la réduction méiotique complète.

\textbf{Barème indicatif (sur 5 pts) :} identification justifiée des populations : 1,5 pt ; explication du double type à $2Q$ : 1 pt ; courbe correcte et annotée : 1 pt ; nombres de chromosomes : 0,75 pt ; conséquence du blocage : 0,75 pt.

\textbf{Points de vigilance :} citer explicitement que la phase S double l'ADN \emph{sans changer le nombre de chromosomes} ; distinguer toujours « quantité d'ADN » et « ploïdie » ; pour la question 5, raisonner sur la non-séparation des chromatides.
\end{corrige}

\begin{corrige}[Exercice 10 — Type Bac : double fécondation chez les Spermaphytes]
\textbf{1. Schéma de l'ovule au moment de la double fécondation.}
\begin{popfigure}
\begin{center}
\begin{tikzpicture}[scale=1.2]
% ovule
\draw[fill=popcream,draw=popdark,thick] (0,0) ellipse (1.6 and 2.4);
% sac embryonnaire
\draw[fill=white,draw=popblue,thick] (0,0.1) ellipse (1.0 and 1.8);
% micropyle
\draw[popdark,thick] (-0.35,2.35) -- (-0.25,1.95);
\draw[popdark,thick] (0.35,2.35) -- (0.25,1.95);
% tube pollinique
\draw[poporange,very thick] (0,3.4) -- (0,1.3);
\draw[poporange,very thick] (0.18,3.4) -- (0.18,1.35);
% noyaux mâles
\draw[fill=poporange,draw=popdark] (0.05,1.15) circle (0.13);
\draw[fill=poporange,draw=popdark] (0.2,0.85) circle (0.13);
% oosphère
\draw[fill=poppinkL,draw=popdark] (-0.35,0.9) circle (0.22);
% noyaux polaires
\draw[fill=popgreenL,draw=popdark] (-0.15,0.05) circle (0.16);
\draw[fill=popgreenL,draw=popdark] (0.2,-0.15) circle (0.16);
% antipodes
\draw[fill=popmuted!40,draw=popdark] (-0.3,-1.3) circle (0.12);
\draw[fill=popmuted!40,draw=popdark] (0.1,-1.45) circle (0.12);
\draw[fill=popmuted!40,draw=popdark] (0.35,-1.2) circle (0.12);
% légendes
\node[font=\scriptsize,right] at (2.0,2.8) {1 : tube pollinique};
\draw[->,popmuted] (2.0,2.7) -- (0.2,2.6);
\node[font=\scriptsize,right] at (2.0,2.2) {2 : micropyle};
\draw[->,popmuted] (2.0,2.1) -- (0.35,2.15);
\node[font=\scriptsize,right] at (2.0,1.2) {3 : noyaux mâles ($n$)};
\draw[->,popmuted] (2.0,1.1) -- (0.25,1.0);
\node[font=\scriptsize,right] at (2.0,0.3) {4 : noyaux polaires ($n$)};
\draw[->,popmuted] (2.0,0.2) -- (0.3,0.0);
\node[font=\scriptsize,left] at (-2.0,0.9) {5 : oosphère ($n$)};
\draw[->,popmuted] (-2.0,0.8) -- (-0.55,0.9);
\node[font=\scriptsize,left] at (-2.0,-0.3) {6 : sac embryonnaire};
\draw[->,popmuted] (-2.0,-0.4) -- (-1.0,-0.3);
\node[font=\scriptsize,left] at (-2.0,-1.3) {7 : antipodes};
\draw[->,popmuted] (-2.0,-1.4) -- (-0.4,-1.35);
\node[font=\scriptsize,left] at (-2.0,-2.1) {8 : paroi de l'ovule};
\draw[->,popmuted] (-2.0,-2.2) -- (-1.3,-1.9);
\end{tikzpicture}
\end{center}
\legende{Ovule de Spermaphyte au moment de la double fécondation : le tube pollinique franchit le micropyle et libère deux noyaux mâles dans le sac embryonnaire.}
\end{popfigure}

\textbf{2. Devenir des deux noyaux mâles :}
\begin{enumerate}
\item Le tube pollinique, guidé par le micropyle, pénètre dans le sac embryonnaire et libère les \textbf{deux noyaux mâles} ($n$).
\item \textbf{Première fusion :} un noyau mâle ($n$) fusionne avec l'\textbf{oosphere} ($n$) $\rightarrow$ \textbf{zygote} diploïde ($2n$).
\item \textbf{Seconde fusion :} l'autre noyau mâle ($n$) fusionne avec les \textbf{deux noyaux polaires} ($n + n$) $\rightarrow$ \textbf{cellule-mère de l'albumen} triploïde ($3n$).
\end{enumerate}

\textbf{3. Ploïdie et chromosomes} ($2n = 20$, donc $n = 10$) : zygote : $2n = 20$ chromosomes ; cellule-mère de l'albumen : $3n = 30$ chromosomes.

\textbf{4. Devenir et rôle :} le zygote se divise par mitoses et se différencie en \textbf{embryon} (future plantule) ; la cellule-mère de l'albumen se divise pour former l'\textbf{albumen}, tissu de réserve riche en amidon (chez le Maïs, graine albuminée). L'albumen \textbf{nourrit la plantule} pendant la germination, jusqu'à ce qu'elle devienne autotrophe grâce à la photosynthèse.

\textbf{5. Tableau comparatif :}
\begin{center}
\begin{tabularx}{\linewidth}{|l|X|X|}\hline
\rowcolor{popblueL} Critère & Mammifères & Spermaphytes \\ \hline
Lieu & Tiers externe de la trompe de Fallope & Sac embryonnaire de l'ovule (dans l'ovaire) \\ \hline
Nombre de fusions nucléaires & Une seule & Deux (double fécondation) \\ \hline
Produits formés & Zygote ($2n$) uniquement & Zygote ($2n$) + cellule-mère de l'albumen ($3n$) \\ \hline
Devenir des produits & Zygote $\rightarrow$ embryon puis f\oe tus (nourri via le placenta) & Zygote $\rightarrow$ embryon ; albumen $\rightarrow$ tissu de réserve de la graine \\ \hline
\end{tabularx}
\end{center}

\textbf{Barème indicatif (sur 5 pts) :} schéma légendé : 1 pt ; étapes de la double fécondation : 1 pt ; ploïdies et chromosomes : 0,75 pt ; devenir et rôle de l'albumen : 0,75 pt ; tableau comparatif : 1,5 pt.

\textbf{Points de vigilance :} ne pas écrire que le grain de pollen est le gamète ; préciser que l'albumen est \emph{triploïde} ($n + n + n$) ; dans le tableau, comparer ligne par ligne sans laisser de case vide.
\end{corrige}

\begin{corrige}[Exercice 11 — Type Bac : méiose, fécondation et pérennité de l'espèce]
\textbf{1. Ploïdie des gamètes après blocage :} les gamètes contiennent 14 chromosomes, soit le stock complet $2n$ : ils sont \textbf{diploïdes} (gamètes non réduits). La méiose étant bloquée, la réduction chromatique ($2n \rightarrow n$) n'a pas eu lieu.

\textbf{2. Calculs :} union de deux gamètes anormaux : $14 + 14 = 28$ chromosomes ($4n$). À la génération suivante, si le phénomène se répète : les gamètes contiendraient 28 chromosomes et le zygote $28 + 28 = 56$ chromosomes ($8n$). Le nombre de chromosomes \textbf{double à chaque génération}.

\textbf{3. Démonstration :} soit une espèce à $2n$ chromosomes se reproduisant par voie sexuée. La fécondation \emph{additionne} les stocks chromosomiques des deux gamètes. Sans méiose, les gamètes restent à $2n$ : le zygote serait à $4n$, puis $8n$, etc. Le caryotype doublerait indéfiniment, ce qui est incompatible avec la survie (déséquilibres génétiques létaux). Or on observe que le caryotype de chaque espèce est \emph{constant} d'une génération à l'autre (Homme : toujours $2n = 46$ ; Maïs : $2n = 20$). Il existe donc nécessairement un mécanisme de \textbf{réduction de moitié} du nombre de chromosomes avant la formation des gamètes : c'est la \cle{méiose} ($2n \rightarrow n$). La méiose est ainsi \emph{indispensable} à la constance du caryotype.

\textbf{4. Rôle complémentaire de la fécondation :} la fécondation ($n + n \rightarrow 2n$) \textbf{rétablit la diploïdie} dans le zygote : elle restitue les paires de chromosomes homologues (un chromosome d'origine paternelle, un d'origine maternelle) et relance le cycle. Sans elle, la méiose seule conduirait à des individus haploïdes. La constance du caryotype résulte de l'\emph{alternance} méiose (réduction) / fécondation (restauration) : $2n \xrightarrow{\text{méiose}} n \xrightarrow{\text{fécondation}} 2n$.

\textbf{5. Synthèse — stabilité et diversité :}
\begin{itemize}
\item \textbf{Stabilité (constance du caryotype) :} la méiose divise par deux le nombre de chromosomes et la quantité d'ADN ($2n$, $4Q \rightarrow n$, $Q$) grâce à une seule réplication suivie de deux divisions ; la fécondation restaure $2n$ et $2Q$ dans le zygote. Le cycle $2n \rightleftharpoons n$ garantit la permanence du caryotype spécifique : c'est la condition de la \emph{pérennité} de l'espèce.
\item \textbf{Diversité (brassage génétique) :} la méiose crée de la variabilité par le \emph{crossing-over} de la prophase I (brassage intrachromosomique) et par la \emph{disjonction aléatoire} des homologues en anaphase I (brassage interchromosomique : $2^{n}$ combinaisons possibles, soit $2^{23}$ chez l'Homme). La fécondation ajoute la \emph{rencontre au hasard} de deux gamètes parmi des millions de combinaisons : chaque zygote est génétiquement unique.
\item \textbf{Conclusion :} la méiose et la fécondation sont deux phénomènes \emph{complémentaires et indissociables} : la première réduit et brasse, la seconde restaure et associe. Ensemble, elles assurent à la fois la \textbf{stabilité} du patrimoine génétique de l'espèce et la \textbf{diversité} des individus, matière première de l'adaptation et de l'évolution : les deux conditions de la pérennité de l'espèce au sein du monde vivant.
\end{itemize}

\textbf{Barème indicatif (sur 5 pts) :} ploïdie des gamètes justifiée : 0,5 pt ; calculs (28 puis 56) : 0,75 pt ; démonstration par l'absurde / par le raisonnement : 1 pt ; rôle de la fécondation : 0,75 pt ; synthèse (stabilité + brassages + complémentarité) : 2 pts.

\textbf{Points de vigilance :} la démonstration attendue repose sur le \emph{raisonnement par l'absurde} (que se passerait-il sans méiose ?) ; citer les deux types de brassage en les nommant précisément ; conclure en revenant explicitement à la notion de \emph{pérennité de l'espèce}, cœur de la famille de situations du programme.
\end{corrige}

\newpage

\coursec{Fiche bilan}

\begin{popfigure}
\begin{center}
\begin{tikzpicture}[mindmap, grow cyclic,
  every node/.style={concept, text=white, font=\small\bfseries},
  root concept/.append style={concept color=popdark, font=\normalsize\bfseries},
  level 1/.append style={level distance=3.4cm, sibling angle=60},
  level 2/.append style={level distance=2.1cm, sibling angle=45, font=\scriptsize},
  concept connection/.append style={color=popmuted}]
\node[root concept]{Méiose \& Fécondation : pérennité de l'espèce}
  child[concept color=popblue]{ node {Gonades}
    child { node {Mammifères : testicules, ovaires} }
    child { node {Spermaphytes : étamines, pistil} }
  }
  child[concept color=poporange]{ node {Méiose}
    child { node {2 divisions, 4 cellules $n$} }
    child { node {Brassage inter et intra} }
  }
  child[concept color=poppurple]{ node {Quantité d'ADN}
    child { node {$2q \to 4q \to 2q \to q$} }
  }
  child[concept color=popgreen]{ node {Gamétogenèse}
    child { node {Spermato / ovogenèse} }
    child { node {Pollen, sac embryonnaire} }
  }
  child[concept color=poppink]{ node {Fécondation}
    child { node {Simple : Mammifères} }
    child { node {Double : Spermaphytes} }
  }
  child[concept color=popgold]{ node {Pérennité}
    child { node {Alternance $n \rightleftharpoons 2n$} }
    child { node {Diversité génétique} }
  };
\end{tikzpicture}
\end{center}
\legende{Carte mentale de la leçon : la méiose (réduction de ploïdie, brassage) et la fécondation (restauration de la diploïdie) sont les deux phénomènes complémentaires qui assurent la pérennité de l'espèce.}
\end{popfigure}

\begin{bilan}
\textbf{1. Définitions clés}
\begin{itemize}
  \item \cle{Gonade} : organe produisant les gamètes (testicule, ovaire chez les Mammifères ; étamine et ovaire de la fleur chez les Spermaphytes).
  \item \cle{Méiose} : succession de deux divisions cellulaires qui, à partir d'une cellule-mère diploïde ($2n$ chromosomes, $2q$ ADN), produit quatre cellules haploïdes ($n$, $q$).
  \item \cle{Gamétogenèse} : formation des gamètes ; spermatogenèse et ovogenèse chez les Mammifères ; pollinogenèse et formation du sac embryonnaire chez les Spermaphytes.
  \item \cle{Fécondation} : fusion de deux gamètes haploïdes ($n + n \to 2n$) formant un \oe uf ou zygote diploïde.
  \item \cle{Double fécondation} (Spermaphytes) : un noyau mâle + oosphère $\to$ zygote ($2n$) ; l'autre noyau mâle + deux noyaux polaires $\to$ albumen ($3n$).
\end{itemize}

\textbf{2. Comportement des chromosomes et de l'ADN (à connaître par c\oe ur)}
\begin{center}
\begin{tabularx}{\linewidth}{|l|X|X|}
\hline
\rowcolor{popblueL} \textbf{Étape} & \textbf{Événement chromosomique} & \textbf{Quantité d'ADN} \\
\hline
Interphase & Réplication de l'ADN (chromosomes à 2 chromatides) & $2q \to 4q$ \\
\hline
Méiose I (réductionnelle) & Séparation des chromosomes homologues (anaphase I) ; 2 cellules $n$ & $4q \to 2q$ \\
\hline
Méiose II (équationnelle) & Séparation des chromatides s\oe urs (anaphase II) ; 4 cellules $n$ & $2q \to q$ \\
\hline
Fécondation & Fusion des noyaux des gamètes & $q + q \to 2q$ \\
\hline
\end{tabularx}
\end{center}

\textbf{3. Sources de la diversité génétique}
\begin{itemize}
  \item \cle{Brassage interchromosomique} : répartition aléatoire des homologues en anaphase I ($2^{n}$ combinaisons).
  \item \cle{Brassage intrachromosomique} : crossing-over en prophase I entre chromatides de chromosomes homologues.
  \item Rencontre aléatoire des gamètes à la fécondation.
\end{itemize}

\textbf{4. Méthodes express}
\begin{itemize}
  \item \emph{Reconnaître une étape de méiose} : chromosomes à 2 chromatides associés en paires d'homologues $\to$ méiose I ; chromosomes à 2 chromatides isolés (cellule $n$) $\to$ méiose II ; chromosomes à 1 chromatide $\to$ fin de méiose.
  \item \emph{Lire la courbe d'ADN} : chute $4q \to 2q$ = fin méiose I ; chute $2q \to q$ = fin méiose II ; remontée $q \to 2q$ = fécondation.
  \item \emph{Identifier les cellules germinales} : spermatogonie ($2n$, $2q$) $\to$ spermatocyte I ($2n$, $4q$) $\to$ spermatocyte II ($n$, $2q$) $\to$ spermatide ($n$, $q$) $\to$ spermatozoïde ($n$, $q$).
\end{itemize}

\textbf{5. Conclusion fondamentale}
\begin{itemize}
  \item La méiose divise par deux le stock chromosomique ($2n \to n$) ; la fécondation le restaure ($n + n \to 2n$). Leur alternance maintient constant le nombre de chromosomes d'une génération à l'autre : c'est la condition de la \cle{pérennité de l'espèce}, tout en assurant la diversité des individus.
\end{itemize}
\end{bilan}

\begin{aretenir}[Checklist avant le Bac]
\begin{itemize}
  \item[$\square$] Je sais situer et nommer les gonades : testicules et ovaires (Mammifères), étamines et pistil (Spermaphytes).
  \item[$\square$] Je sais décrire les deux divisions de la méiose et le rôle de chacune (réductionnelle puis équationnelle).
  \item[$\square$] Je sais que la séparation des homologues a lieu en anaphase I et celle des chromatides s\oe urs en anaphase II.
  \item[$\square$] Je connais la séquence des quantités d'ADN : $2q \to 4q \to 2q \to q$, puis $q + q \to 2q$ à la fécondation.
  \item[$\square$] Je sais reconnaître les étapes de la méiose sur une micrographie ou un schéma.
  \item[$\square$] Je sais ordonner les cellules de la lignée germinale (spermatogonie $\to$ spermatozoïde ; ovogonie $\to$ ovocyte II).
  \item[$\square$] Je sais expliquer la double fécondation : zygote $2n$ et albumen $3n$.
  \item[$\square$] Je sais citer les deux brassages (inter et intrachromosomique) et leur conséquence : la diversité génétique.
  \item[$\square$] Je sais justifier que méiose et fécondation sont complémentaires pour maintenir $2n$ constant au fil des générations.
  \item[$\square$] Je rédige avec le vocabulaire exact : \emph{haploïde}, \emph{diploïde}, \emph{chromosomes homologues}, \emph{zygote}, sans confondre méiose et mitose.
\end{itemize}
\end{aretenir}

\findecours

\end{document}
